% chapters/ch09.tex
% chapters/ch09.tex
% chapters/ch09.tex
% 第~
% \ref{chap:9}~章 Sobolev 空间与 N 维椭圆边值问题的变分表述
% 源: 原书 PDF 第 269--329 页 (印刷页 263--323)

\chapter{Sobolev 空间与 N 维椭圆边值问题的变分表述}
\label{chap:9}

\section{\index{NweiSobolevkongjian@N 维 Sobolev 空间 (Sobolev space in dimension N)}
Sobolev 空间 $W^{1,p}(\Omega)$ 的定义与初等性质}
\label{sec:9.1}

设 $\Omega\subset\R^N$ 为开集, $p\in\R$ 满足 $1\le p\le\infty$.

\begin{definition}
\emph{Sobolev 空间} $W^{1,p}(\Omega)$ 定义为\footnote{若无混淆, 我们常写
$W^{1,p}$ 而不写 $W^{1,p}(\Omega)$.}
\begin{equation*}
W^{1,p}(\Omega)=\left\{u\in L^p(\Omega)\ \middle|\
\begin{array}{l}
\exists g_1,g_2,\dots,g_N\in L^p(\Omega)\ \text{使得}\\
\displaystyle\int_\Omega u\frac{\partial\varphi}{\partial x_i}
=-\int_\Omega g_i\varphi\quad\forall\varphi\in C_c^\infty(\Omega),\
\forall i=1,2,\dots,N
\end{array}
\right\}.
\end{equation*}
我们令
\begin{equation*}
H^1(\Omega)=W^{1,2}(\Omega).
\end{equation*}
对 $u\in W^{1,p}(\Omega)$, 我们定义\footnote{这个定义是有意义的: 由推论
 \ref{cor:4.24}, $g_i$ (几乎处处) 唯一.}
$\dfrac{\partial u}{\partial x_i}=g_i$, 并记
\begin{equation*}
\nabla u=\operatorname{grad}u=
\left(\frac{\partial u}{\partial x_1},\frac{\partial u}{\partial x_2},
\dots,\frac{\partial u}{\partial x_N}\right).
\end{equation*}
空间 $W^{1,p}(\Omega)$ 配以范数
\begin{equation*}
\norm{u}_{W^{1,p}}=\norm{u}_p+\sum_{i=1}^N
\norm{\frac{\partial u}{\partial x_i}}_p,
\end{equation*}
有时 (若 $1\le p<\infty$) 也配以等价范数
$(\norm{u}_p^p+\sum_{i=1}^N\norm{\frac{\partial u}{\partial x_i}}_p^p)^{1/p}$.
空间 $H^1(\Omega)$ 配以标量积
\begin{equation*}
(u,v)_{H^1}=(u,v)_{L^2}+\sum_{i=1}^N
\left(\frac{\partial u}{\partial x_i},\frac{\partial v}{\partial x_i}\right)_{L^2}
=\int_\Omega uv+\sum_{i=1}^N\frac{\partial u}{\partial x_i}
\frac{\partial v}{\partial x_i}.
\end{equation*}
相应的范数
\begin{equation*}
\norm{u}_{H^1}=\left(\norm{u}_2^2+\sum_{i=1}^N
\norm{\frac{\partial u}{\partial x_i}}_2^2\right)^{1/2}
\end{equation*}
与 $W^{1,2}$ 范数等价.
\end{definition}

\begin{keyproposition}
\label{prop:9.1}
$W^{1,p}(\Omega)$ 对每个 $1\le p\le\infty$ 都是 Banach 空间.当
$1<p<\infty$ 时它自反, 当 $1\le p<\infty$ 时它可分.$H^1(\Omega)$ 是可分
Hilbert 空间.
\end{keyproposition}

\begin{proof}
仿照命题  \ref{prop:8.1} 的证明, 取算子 $Tu=[u,\nabla u]$.
\end{proof}

\begin{remark}
\label{rem:9.1}
在 $W^{1,p}$ 的定义中, 我们本来同样可以用 $C_c^\infty(\Omega)$ 作为检验函数
$\varphi$ 的类 (而不用 $C_c^1(\Omega)$); 为证明这一点, 利用一列磨光核
$(\rho_n)$.
\end{remark}

\begin{remark}
\label{rem:9.2}
显然, 若 $u\in C^1(\Omega)\cap L^p(\Omega)$ 且
$\frac{\partial u}{\partial x_i}\in L^p(\Omega)$ 对一切 $i=1,2,\dots,N$
成立 (这里 $\frac{\partial u}{\partial x_i}$ 指 $u$ 的通常偏导数), 则
$u\in W^{1,p}(\Omega)$.此外, 通常的偏导数与 $W^{1,p}$ 意义下的偏导数一致,
所以记号是相容的.特别地, 若 $\Omega$ 有界, 则
$C^1(\overline\Omega)\subset W^{1,p}(\Omega)$ 对一切 $1\le p\le\infty$
成立.反之可以证明: 若 $u\in W^{1,p}(\Omega)$ (其中 $1\le p\le\infty$) 且
$\frac{\partial u}{\partial x_i}\in C(\Omega)$ 对一切 $i=1,2,\dots,N$ 成立
(这里 $\frac{\partial u}{\partial x_i}$ 指 $W^{1,p}$ 意义下的偏导数), 则
$u\in C^1(\Omega)$; 更精确地说, 存在函数 $\tilde u\in C^1(\Omega)$ 使得
$u=\tilde u$ 几乎处处成立.
\end{remark}

\begin{starremark}
\label{rem:9.3}
对每个 $u\in L^1_{\mathrm{loc}}(\Omega)$, 分布论赋予
$\frac{\partial u}{\partial x_i}$ 一个意义 ($\frac{\partial u}{\partial x_i}$
是巨大的分布空间 $\mathcal{D}'(\Omega)$ 中的元素, 这个空间特别地包含
$L^1_{\mathrm{loc}}(\Omega)$).用分布的语言可以说, $W^{1,p}(\Omega)$ 就是
那些使得所有偏导数 $\frac{\partial u}{\partial x_i}$ ($1\le i\le N$, 分布
意义下的) 都属于 $L^p(\Omega)$ 的函数 $u\in L^p(\Omega)$ 的集合.

当 $\Omega=\R^N$ 且 $p=2$ 时, 也可以用 Fourier 变换来定义 Sobolev 空间;
参见例如 J.-L.~Lions--E.~Magenes \bref{LionsMagenes}{1}, P.~Malliavin \bref{Malliavin}{1}, H.~Triebel \bref{Triebel}{1}, L.~Grafakos \bref{Grafakos}{1}.这里不采取这个观点.
\end{starremark}

\begin{remark}
\label{rem:9.4}
记住下述事实是有用的:

\noindent (i) 设 $(u_n)$ 为 $W^{1,p}$ 中的序列, 使得 $u_n\to u$ 在 $L^p$ 中,
且 $(\nabla u_n)$ 在 $(L^p)^N$ 中收敛到某个极限.则 $u\in W^{1,p}$ 且
$\norm{u_n-u}_{W^{1,p}}\to0$.当 $1<p\le\infty$ 时, 只须知道 $u_n\to u$
在 $L^p$ 中且 $(\nabla u_n)$ 在 $(L^p)^N$ 中有界, 就能推出 $u\in W^{1,p}$
(为什么?).

\noindent (ii) 给定定义在 $\Omega$ 上的函数 $f$, 我们记 $\bar f$ 为它在
$\Omega$ 外的延拓, 即
\begin{equation*}
\bar f(x)=
\begin{cases}
f(x) & \text{若}\ x\in\Omega,\\
0 & \text{若}\ x\in\R^N\setminus\Omega.
\end{cases}
\end{equation*}
设 $u\in W^{1,p}(\Omega)$, $\alpha\in C_c^1(\Omega)$.则\footnote{注意: 一般
说来 $\bar u\notin W^{1,p}(\R^N)$ (为什么?).}
\begin{equation*}
\alpha u\in W^{1,p}(\R^N)
\quad\text{且}\quad
\frac{\partial}{\partial x_i}(\alpha u)=\alpha\frac{\partial u}{\partial x_i}
+\frac{\partial\alpha}{\partial x_i}u.
\end{equation*}
事实上, 设 $\varphi\in C_c^1(\R^N)$; 我们有
\begin{align*}
\int_{\R^N}\alpha u\frac{\partial\varphi}{\partial x_i}
&=\int_\Omega\alpha u\frac{\partial\varphi}{\partial x_i}
=\int_\Omega u\left(\frac{\partial}{\partial x_i}(\alpha\varphi)
-\frac{\partial\alpha}{\partial x_i}\varphi\right)\\
&=-\int_\Omega\left(\frac{\partial u}{\partial x_i}\alpha\varphi
+u\frac{\partial\alpha}{\partial x_i}\varphi\right)
=-\int_{\R^N}\left(\alpha\frac{\partial u}{\partial x_i}
+\frac{\partial\alpha}{\partial x_i}u\right)\varphi.
\end{align*}
如果把 $\alpha\in C_c^1(\Omega)$ 换成
$\alpha\in C^1(\R^N)\cap L^\infty(\R^N)$, 其中
$\nabla\alpha\in L^\infty(\R^N)^N$ 且
$\supp\alpha\subset\R^N\setminus(\partial\Omega)$, 结论仍然成立.

下面是对一般开集 $\Omega$ 成立的第一个稠密性结果; 我们稍后在附加假设下
(推论  \ref{cor:9.8}) 建立一个更精确的结果.我们需要下述概念.
\end{remark}

\begin{definition}
设 $\Omega\subset\R^N$ 为开集.我们说 $\R^N$ 中的开集 $\omega$ \emph{强含于}
$\Omega$, 并记 $\omega\subset\subset\Omega$, 如果 $\overline\omega\subset\Omega$
且 $\overline\omega$ 紧.\footnote{$\overline\omega$ 表示 $\omega$ 在 $\R^N$
中的闭包.}
\end{definition}

\begin{keytheorem}[Friedrichs]
\label{thm:9.2}
设 $u\in W^{1,p}(\Omega)$, $1\le p<\infty$.则存在 $C_c^\infty(\R^N)$ 中的
序列 $(u_n)$ 使得
\begin{equation}
u_n|_\Omega\to u\quad\text{在}\ L^p(\Omega)\ \text{中}
\label{eq:9.1}
\end{equation}
且
\begin{equation}
\nabla u_n|_\omega\to\nabla u|_\omega
\quad\text{在}\ L^p(\omega)^N\ \text{中对一切}\ \omega\subset\subset\Omega.
\label{eq:9.2}
\end{equation}
若 $\Omega=\R^N$ 且 $u\in W^{1,p}(\R^N)$, $1\le p<\infty$, 则存在
$C_c^\infty(\R^N)$ 中的序列 $(u_n)$ 使得
\begin{equation*}
u_n\to u\quad\text{在}\ L^p(\R^N)\ \text{中}
\quad\text{且}\quad
\nabla u_n\to\nabla u\quad\text{在}\ L^p(\R^N)^N\ \text{中}.
\end{equation*}
\end{keytheorem}

在证明中我们要用到下述引理.

\begin{lemma}
\label{lem:9.1}
设 $\rho\in L^1(\R^N)$, $v\in W^{1,p}(\R^N)$, $1\le p\le\infty$.则
\begin{equation*}
\rho\star v\in W^{1,p}(\R^N)
\quad\text{且}\quad
\frac{\partial}{\partial x_i}(\rho\star v)=\rho\star\frac{\partial v}{\partial x_i}
\quad\forall i=1,2,\dots,N.
\end{equation*}
\end{lemma}

\begin{proof}[引理 \ref{lem:9.1}]
仿照引理  \ref{lem:8.4} 的证明.
\end{proof}

\begin{proof}[定理 \ref{thm:9.2}]
令
\begin{equation*}
\bar u(x)=
\begin{cases}
u(x) & \text{若}\ x\in\Omega,\\
0 & \text{若}\ x\in\R^N\setminus\Omega,
\end{cases}
\end{equation*}
并令 $v_n=\rho_n\star\bar u$ (其中 $(\rho_n)$ 是一列磨光核).我们知道
(见第~\ref{sec:4.4}~节) $v_n\in C^\infty(\R^N)$ 且 $v_n\to\bar u$ 在 $L^p(\R^N)$ 中.
我们断言: 对一切 $\omega\subset\subset\Omega$,
$\nabla v_n|_\omega\to\nabla u|_\omega$ 在 $L^p(\omega)^N$ 中.事实上,
给定 $\omega\subset\subset\Omega$, 固定函数 $\alpha\in C_c^1(\Omega)$,
$0\le\alpha\le1$, 使得 $\alpha=1$ 在 $\omega$ 的某个邻域上.当 $n$ 足够大时
我们有
\begin{equation}
\rho_n\star(\alpha u)=\rho_n\star\bar u\quad\text{在}\ \omega\ \text{上},
\label{eq:9.3}
\end{equation}
因为
\begin{align*}
\supp(\rho_n\star\alpha u-\rho_n\star\bar u)
&=\supp\bigl(\rho_n\star(1-\bar\alpha)\bar u\bigr)\\
&\subset\supp\rho_n+\supp\bigl((1-\bar\alpha)\bar u\bigr)
\subset B\left(0,\frac1n\right)+\supp(1-\bar\alpha)
\subset(\omega)^c
\end{align*}
对足够大的 $n$ 成立.由引理  \ref{lem:9.1} 与注记 \ref{rem:9.4}(ii) 我们有
\begin{equation*}
\frac{\partial}{\partial x_i}(\rho_n\star\alpha u)
=\rho_n\star\left(\alpha\frac{\partial u}{\partial x_i}
+\frac{\partial\alpha}{\partial x_i}u\right).
\end{equation*}
由此推出
\begin{equation*}
\frac{\partial}{\partial x_i}(\rho_n\star\alpha u)
\to\alpha\frac{\partial u}{\partial x_i}
+\frac{\partial\alpha}{\partial x_i}u\quad\text{在}\ L^p(\R^N)\ \text{中},
\end{equation*}
特别地
\begin{equation*}
\frac{\partial}{\partial x_i}(\rho_n\star\alpha u)
\to\frac{\partial u}{\partial x_i}\quad\text{在}\ L^p(\omega)\ \text{中}.
\end{equation*}
由 \eqref{eq:9.3} 我们有
\begin{equation*}
\frac{\partial}{\partial x_i}(\rho_n\star\bar u)
\to\frac{\partial u}{\partial x_i}\quad\text{在}\ L^p(\omega)\ \text{中}.
\end{equation*}
最后, 我们用一列截断函数 $(\zeta_n)$ 乘序列 $(v_n)$, 如定理  \ref{thm:8.7}
的证明中那样.\footnote{本章中我们总是用 $(\zeta_n)$ 表示一列截断函数, 即
固定函数 $\zeta\in C_c^\infty(\R^N)$, $0\le\zeta\le1$, 使得
\begin{equation*}
\zeta(x)=
\begin{cases}
1 & \text{若}\ \abs{x}\le1,\\
0 & \text{若}\ \abs{x}\ge2,
\end{cases}
\end{equation*}
并令 $\zeta_n(x)=\zeta(x/n)$, $n=1,2,\dots$.} 容易验证序列
$u_n=\zeta_nv_n$ 具有所要的性质, 即 $u_n\in C_c^\infty(\R^N)$,
$u_n\to u$ 在 $L^p(\Omega)$ 中, 且对每个 $\omega\subset\subset\Omega$,
$\nabla u_n\to\nabla u$ 在 $(L^p(\omega))^N$ 中.

当 $\Omega=\R^N$ 时, 序列 $u_n=\zeta_n(\rho_n\star u)$ 具有所要的性质.
\end{proof}

\begin{starremark}
\label{rem:9.5}
可以证明 (Meyers--Serrin 定理): 若 $u\in W^{1,p}(\Omega)$, $1\le p<\infty$,
则存在 $C^\infty(\Omega)\cap W^{1,p}(\Omega)$ 中的序列 $(u_n)$ 使得
$u_n\to u$ 在 $W^{1,p}(\Omega)$ 中; 这个结果的证明相当精细 (参见例如
R.~Adams \bref{Adams}{1} 或 A.~Friedman \bref{Friedman}{2}).一般说来,
若 $\Omega$ 是任意开集且 $u\in W^{1,p}(\Omega)$, 则未必存在
$C_c^1(\R^N)$ 中的序列 $(u_n)$ 使得 $u_n|_\Omega\to u$ 在 $W^{1,p}(\Omega)$
中.请把 Meyers--Serrin 定理 (它对任意开集成立) 与推论  \ref{cor:9.8}
(它假设 $\Omega$ 是正则的) 作比较.
\end{starremark}

下面是 $W^{1,p}$ 函数的一个简单刻画:

\begin{proposition}
\label{prop:9.3}
设 $u\in L^p(\Omega)$, $1<p\le\infty$.下述性质等价:

\noindent (i) $u\in W^{1,p}(\Omega)$;

\noindent (ii) 存在常数 $C$ 使得
\begin{equation*}
\left|\int_\Omega u\frac{\partial\varphi}{\partial x_i}\right|
\le C\norm{\varphi}_{L^{p'}(\Omega)}
\quad\forall\varphi\in C_c^\infty(\Omega),\
\forall i=1,2,\dots,N;
\end{equation*}

\noindent (iii) 存在常数 $C$ 使得对一切 $\omega\subset\subset\Omega$ 与一切
$h\in\R^N$, $\abs{h}<\dist(\omega,\partial\Omega)$, 有
\begin{equation*}
\norm{\tau_hu-u}_{L^p(\omega)}\le C\abs{h}.
\end{equation*}
(注意: 对 $x\in\omega$ 与 $\abs{h}<\dist(\omega,\partial\Omega)$,
$\tau_hu(x)=u(x+h)$ 有意义.)

此外, 在 (ii) 与 (iii) 中都可以取 $C=\norm{\nabla u}_{L^p(\Omega)}$.
若 $\Omega=\R^N$, 我们有
\begin{equation*}
\norm{\tau_hu-u}_{L^p(\R^N)}\le\abs{h}\norm{\nabla u}_{L^p(\R^N)}.
\end{equation*}
\end{proposition}

\begin{proof}
\noindent (i)$\Rightarrow$(ii). 显然.

\noindent (ii)$\Rightarrow$(i). 仿照命题  \ref{prop:8.3} 的证明.

\noindent (i)$\Rightarrow$(iii). 先设 $u\in C_c^\infty(\R^N)$.设 $h\in\R^N$,
令
\begin{equation*}
v(t)=u(x+th),\quad t\in\R.
\end{equation*}
则 $v'(t)=h\cdot\nabla u(x+th)$, 于是
\begin{equation*}
u(x+h)-u(x)=v(1)-v(0)=\int_0^1v'(t)\,dt
=\int_0^1h\cdot\nabla u(x+th)\,dt.
\end{equation*}
由此推出, 当 $1\le p<\infty$ 时,
\begin{equation*}
\abs{\tau_hu(x)-u(x)}^p\le\abs{h}^p\int_0^1
\abs{\nabla u(x+th)}^p\,dt,
\end{equation*}
且
\begin{align*}
\int_\omega\abs{\tau_hu(x)-u(x)}^p\,dx
&\le\abs{h}^p\int_\omega dx\int_0^1\abs{\nabla u(x+th)}^p\,dt\\
&=\abs{h}^p\int_0^1dt\int_\omega\abs{\nabla u(x+th)}^p\,dx
=\abs{h}^p\int_0^1dt\int_{\omega+th}\abs{\nabla u(y)}^p\,dy.
\end{align*}
若 $\abs{h}<\dist(\omega,\partial\Omega)$, 则存在开集
$\omega'\subset\subset\Omega$ 使得 $\omega+th\subset\omega'$ 对一切
$t\in[0,1]$ 成立, 于是
\begin{equation}
\norm{\tau_hu-u}_{L^p(\omega)}^p\le\abs{h}^p\int_{\omega'}
\abs{\nabla u}^p.
\label{eq:9.4}
\end{equation}
这就完成了 $u\in C_c^\infty(\R^N)$, $1\le p<\infty$ 情形的 (iii) 的证明.
现在设 $u\in W^{1,p}(\Omega)$, $1\le p<\infty$.由定理  \ref{thm:9.2}, 存在
$C_c^\infty(\R^N)$ 中的序列 $(u_n)$ 使得 $u_n\to u$ 在 $L^p(\Omega)$ 中,
且对一切 $\omega\subset\subset\Omega$, $\nabla u_n\to\nabla u$ 在
$L^p(\omega)^N$ 中.把 \eqref{eq:9.4} 用于 $(u_n)$ 并取极限, 即得对每个
$u\in W^{1,p}(\Omega)$ ($1\le p<\infty$) 成立的 (iii).当 $p=\infty$ 时,
对上文 (关于 $p<\infty$) 的结论令 $p\to\infty$.

\noindent (iii)$\Rightarrow$(ii). 设 $\varphi\in C_c^\infty(\Omega)$, 考虑开集
$\omega$ 使得 $\supp\varphi\subset\omega\subset\subset\Omega$.设 $h\in\R^N$,
$\abs{h}<\dist(\omega,\partial\Omega)$.由 (iii) 我们有
\begin{equation*}
\left|\int_\Omega(\tau_hu-u)\varphi\right|
\le C\abs{h}\norm{\varphi}_{L^{p'}(\Omega)}.
\end{equation*}
另一方面, 由于
\begin{equation*}
\int_\Omega(u(x+h)-u(x))\varphi(x)\,dx
=\int_\Omega u(y)\bigl(\varphi(y-h)-\varphi(y)\bigr)\,dy,
\end{equation*}
由此推出
\begin{equation*}
\left|\int_\Omega u(y)\frac{\varphi(y-h)-\varphi(y)}{\abs{h}}\,dy\right|
\le C\norm{\varphi}_{L^{p'}(\Omega)}.
\end{equation*}
取 $h=te_i$, $t\in\R$, 并令 $t\to0$, 即得 (ii).
\end{proof}

\begin{starremark}
\label{rem:9.6}
当 $p=1$ 时, 下述蕴涵仍然成立:
\begin{equation*}
\text{(i)}\Rightarrow\text{(ii)}\iff\text{(iii)}.
\end{equation*}
满足 $p=1$ 的 (ii) (或 (iii)) 的函数称为\emph{有界变差函数} (用分布的语言,
有界变差函数就是一个 $L^1$ 函数, 其所有一阶导数 (分布意义下的) 都是有界
测度).这个空间在许多应用中起着重要作用.人们在极小曲面理论中 (参见例如
E.~Giusti \bref{Giusti}{1} 以及 E.~DeGiorgi, M.~Miranda 等人的著作及其中
所引文献), 在弹性与塑性问题中 (有限变形函数, 参见例如
R.~Temam--G.~Strang \bref{TemamStrang}{2} 以及所引的 P.~Suquet 的工作), 在允许间断解的一阶拟线性方程即所谓守恒律中 (参见例如
A.~I.~Volpert \bref{Volpert}{1} 与 A.~Bressan \bref{Bressan}{1}),
都会遇到有界变差函数 (或具有类似性质的函数).关于这个广阔的课题, 还可参看
L.~Ambrosio--N.~Fusco--D.~Pallara \bref{Ambrosio}{1} 的书以及本章末的
评注 16.
\end{starremark}

\begin{remark}
\label{rem:9.7}
命题  \ref{prop:9.3} ((i)$\Rightarrow$(iii)) 表明, 任何
$u\in W^{1,\infty}(\Omega)$ 在 $\Omega$ 上都有连续的代表元.更精确地说,
若 $\Omega$ 连通, 则
\begin{equation}
\abs{u(x)-u(y)}\le\norm{\nabla u}_{L^\infty(\Omega)}
\dist_\Omega(x,y)\quad\forall x,y\in\Omega
\label{eq:9.5}
\end{equation}
(对上述连续代表元 $u$ 成立), 其中 $\dist_\Omega(x,y)$ 表示 $\Omega$ 中从
$x$ 到 $y$ 的测地距离; 特别地, 若 $\Omega$ 凸, 则
$\dist_\Omega(x,y)=\abs{x-y}$.由此还可推出: 若
$u\in W^{1,p}(\Omega)$ (其中 $1\le p\le\infty$, $\Omega$ 为某开集) 且
$\nabla u=0$ 在 $\Omega$ 上几乎处处成立, 则 $u$ 在 $\Omega$ 的每个连通分支上
为常数.
\end{remark}

\begin{proposition}[乘积的微分]
\label{prop:9.4}
设 $u,v\in W^{1,p}(\Omega)\cap L^\infty(\Omega)$, $1\le p\le\infty$.则
$uv\in W^{1,p}(\Omega)\cap L^\infty(\Omega)$ 且
\begin{equation*}
\frac{\partial}{\partial x_i}(uv)=\frac{\partial u}{\partial x_i}v
+u\frac{\partial v}{\partial x_i},\quad i=1,2,\dots,N.
\end{equation*}
\end{proposition}

\begin{proof}
如推论  \ref{cor:8.10} 的证明中那样, 只须考虑 $1\le p<\infty$ 的情形.由
定理  \ref{thm:9.2}, 存在 $C_c^\infty(\R^N)$ 中的序列 $(u_n),(v_n)$ 使得
\begin{equation*}
u_n\to u,\quad v_n\to v\quad\text{在}\ L^p(\Omega)\ \text{中且几乎处处在}\
\Omega\ \text{上},
\end{equation*}
\begin{equation*}
\nabla u_n\to\nabla u,\quad\nabla v_n\to\nabla v\quad\text{在}
\ L^p(\omega)^N\ \text{中对一切}\ \omega\subset\subset\Omega.
\end{equation*}
检查定理  \ref{thm:9.2} 的证明, 我们容易看出还有
\begin{equation*}
\norm{u_n}_{L^\infty(\R^N)}\le\norm{u}_{L^\infty(\Omega)}
\quad\text{且}\quad
\norm{v_n}_{L^\infty(\R^N)}\le\norm{v}_{L^\infty(\Omega)}.
\end{equation*}
另一方面,
\begin{equation*}
\int_\Omega u_nv_n\frac{\partial\varphi}{\partial x_i}
=-\int_\Omega\left(\frac{\partial u_n}{\partial x_i}v_n
+u_n\frac{\partial v_n}{\partial x_i}\right)\varphi
\quad\forall\varphi\in C_c^1(\Omega).
\end{equation*}
取极限 (利用控制收敛定理), 这成为
\begin{equation*}
\int_\Omega uv\frac{\partial\varphi}{\partial x_i}
=-\int_\Omega\left(\frac{\partial u}{\partial x_i}v
+u\frac{\partial v}{\partial x_i}\right)\varphi
\quad\forall\varphi\in C_c^1(\Omega).
\end{equation*}
\end{proof}

\begin{proposition}[复合的微分]
\label{prop:9.5}
设 $G\in C^1(\R)$ 满足 $G(0)=0$ 且 $\abs{G'(s)}\le M$ 对一切 $s\in\R$
成立, 其中 $M$ 为某常数.设 $u\in W^{1,p}(\Omega)$, $1\le p\le\infty$.则
\begin{equation*}
G\circ u\in W^{1,p}(\Omega)
\quad\text{且}\quad
\frac{\partial}{\partial x_i}(G\circ u)=(G'\circ u)
\frac{\partial u}{\partial x_i},\quad i=1,2,\dots,N.
\end{equation*}
\end{proposition}

\begin{proof}
我们有 $\abs{G(s)}\le M\abs{s}$ 对一切 $s\in\R$ 成立, 于是
$\abs{G\circ u}\le M\abs{u}$; 因此 $G\circ u\in L^p(\Omega)$, 并且
$(G'\circ u)\frac{\partial u}{\partial x_i}\in L^p(\Omega)$.剩下要验证
\begin{equation}
\int_\Omega(G\circ u)\frac{\partial\varphi}{\partial x_i}
=-\int_\Omega(G'\circ u)\frac{\partial u}{\partial x_i}\varphi
\quad\forall\varphi\in C_c^1(\Omega).
\label{eq:9.6}
\end{equation}
当 $1\le p<\infty$ 时, 取 $C_c^\infty(\R^N)$ 中的序列 $(u_n)$ 使得
$u_n\to u$ 在 $L^p(\Omega)$ 中且几乎处处在 $\Omega$ 上,
$\nabla u_n\to\nabla u$ 在 $L^p(\omega)^N$ 中对一切
$\omega\subset\subset\Omega$ 成立 (定理  \ref{thm:9.2}).我们有
\begin{equation*}
\int_\Omega(G\circ u_n)\frac{\partial\varphi}{\partial x_i}
=-\int_\Omega(G'\circ u_n)\frac{\partial u_n}{\partial x_i}\varphi
\quad\forall\varphi\in C_c^1(\Omega).
\end{equation*}
但 $G\circ u_n\to G\circ u$ 在 $L^p(\Omega)$ 中, 且
$(G'\circ u_n)\frac{\partial u_n}{\partial x_i}\to
(G'\circ u)\frac{\partial u}{\partial x_i}$ 在 $L^p(\omega)$ 中对一切
$\omega\subset\subset\Omega$ 成立 (由控制收敛), 于是 \eqref{eq:9.6}
得证.当 $p=\infty$ 时, 固定开集 $\Omega'$ 使得
$\supp\varphi\subset\Omega'\subset\subset\Omega$.则
$u\in W^{1,p}(\Omega')$ 对一切 $p<\infty$ 成立, 且 \eqref{eq:9.6} 由上述
结果推出.
\end{proof}

\begin{proposition}[变量替换公式]
\label{prop:9.6}
设 $\Omega$ 与 $\Omega'$ 为 $\R^N$ 中的两个开集, $H:\Omega'\to\Omega$ 为
双射, $x=H(y)$, 使得 $H\in C^1(\Omega')$, $H^{-1}\in C^1(\Omega)$,
$\operatorname{Jac}H\in L^\infty(\Omega')$,
$\operatorname{Jac}H^{-1}\in L^\infty(\Omega)$.\footnote{$\operatorname{Jac}H$
表示 Jacobi 矩阵 $\frac{\partial H_i}{\partial y_j}$; 于是它是
$L^\infty(\Omega')^{N\times N}$ 中的函数.} 设 $u\in W^{1,p}(\Omega)$,
$1\le p\le\infty$.则 $u\circ H\in W^{1,p}(\Omega')$ 且
\begin{equation*}
\frac{\partial}{\partial y_j}u(H(y))=\sum_i
\frac{\partial u}{\partial x_i}(H(y))\frac{\partial H_i}{\partial y_j}(y)
\quad\forall j=1,2,\dots,N.
\end{equation*}
\end{proposition}

\begin{proof}
当 $1\le p<\infty$ 时, 取 $C_c^\infty(\R^N)$ 中的序列 $(u_n)$ 使得
$u_n\to u$ 在 $L^p(\Omega)$ 中, 且 $\nabla u_n\to\nabla u$ 在
$L^p(\omega)^N$ 中对一切 $\omega\subset\subset\Omega$ 成立.于是
$u_n\circ H\to u\circ H$ 在 $L^p(\Omega')$ 中, 且
\begin{equation*}
\left(\frac{\partial u_n}{\partial x_i}\circ H\right)
\frac{\partial H_i}{\partial y_j}
\to\left(\frac{\partial u}{\partial x_i}\circ H\right)
\frac{\partial H_i}{\partial y_j}
\quad\text{在}\ L^p(\omega')\ \text{中对一切}\ \omega'\subset\subset\Omega'.
\end{equation*}
给定 $\psi\in C_c^1(\Omega')$, 我们有
\begin{equation*}
\int_{\Omega'}(u_n\circ H)\frac{\partial\psi}{\partial y_j}\,dy
=-\int_{\Omega'}\sum_i\left(\frac{\partial u_n}{\partial x_i}\circ H\right)
\frac{\partial H_i}{\partial y_j}\psi\,dy.
\end{equation*}
取极限即得所要结果.当 $p=\infty$ 时, 按命题  \ref{prop:9.5} 证明的末尾
同样的方式处理.
\end{proof}

\subsection*{空间 $W^{m,p}(\Omega)$}

设 $m\ge2$ 为整数, $p$ 为实数, $1\le p\le\infty$.我们归纳地定义
\begin{equation*}
W^{m,p}(\Omega)=\left\{u\in W^{m-1,p}(\Omega);\
\frac{\partial u}{\partial x_i}\in W^{m-1,p}(\Omega)\
\forall i=1,2,\dots,N\right\}.
\end{equation*}
等价地, 这些集合也可以这样引入:
\begin{equation*}
W^{m,p}(\Omega)=\left\{u\in L^p(\Omega)\ \middle|\
\begin{array}{l}
\forall\alpha,\ \abs{\alpha}\le m,\ \exists g_\alpha\in L^p(\Omega)\
\text{使得}\\
\displaystyle\int_\Omega uD^\alpha\varphi=(-1)^{\abs{\alpha}}
\int_\Omega g_\alpha\varphi\quad\forall\varphi\in C_c^\infty(\Omega)
\end{array}
\right\},
\end{equation*}
其中我们采用标准的多重指标记号
$\alpha=(\alpha_1,\alpha_2,\dots,\alpha_N)$, $\alpha_i\ge0$ 为整数,
\begin{equation*}
\abs{\alpha}=\sum_{i=1}^N\alpha_i
\quad\text{且}\quad
D^\alpha\varphi=\frac{\partial^{\abs{\alpha}}\varphi}
{\partial x_1^{\alpha_1}\partial x_2^{\alpha_2}\cdots\partial x_N^{\alpha_N}}.
\end{equation*}
我们令 $D^\alpha u=g_\alpha$.空间 $W^{m,p}(\Omega)$ 配以范数
\begin{equation*}
\norm{u}_{W^{m,p}}=\sum_{0\le\abs{\alpha}\le m}\norm{D^\alpha u}_p
\end{equation*}
是 Banach 空间.

空间 $H^m(\Omega)=W^{m,2}(\Omega)$ 配以标量积
\begin{equation*}
(u,v)_{H^m}=\sum_{0\le\abs{\alpha}\le m}(D^\alpha u,D^\alpha v)_{L^2}
\end{equation*}
是 Hilbert 空间.

\begin{remark}
\label{rem:9.8}
可以证明: 若 $\Omega$ ``足够光滑''且 $\Gamma=\partial\Omega$ 有界, 则
$W^{m,p}(\Omega)$ 上的范数等价于范数
\begin{equation*}
\norm{u}_p+\sum_{\abs{\alpha}=m}\norm{D^\alpha u}_p.
\end{equation*}
更精确地说, 可以证明: 对每个多重指标 $\alpha$, $0<\abs{\alpha}<m$, 以及
每个 $\varepsilon>0$, 存在常数 $C$ (依赖于 $\Omega,\varepsilon,\alpha$)
使得
\begin{equation*}
\norm{D^\alpha u}_p\le\varepsilon\sum_{\abs{\beta}=m}\norm{D^\beta u}_p
+C\norm{u}_p\quad\forall u\in W^{m,p}(\Omega)
\end{equation*}
(参见例如 R.~Adams \bref{Adams}{1}).
\end{remark}

\section{延拓算子}
\label{sec:9.2}

通过先考虑 $\Omega=\R^N$ 的情形来建立 $W^{1,p}(\Omega)$ 中函数的性质, 往往
是方便的 (例如见第~\ref{sec:9.3}~节的结果).因此, 能够把一个函数
$u\in W^{1,p}(\Omega)$ 延拓为函数 $\tilde u\in W^{1,p}(\R^N)$ 是有用的.
这并非总是可能的 (在一般区域 $\Omega$ 中).然而, 若 $\Omega$ ``光滑'',
这样的延拓可以构造出来.让我们先确切说明光滑开集的概念.

\begin{notation}
给定 $x\in\R^N$, 写
\begin{equation*}
x=(x',x_N),\quad x'\in\R^{N-1},\quad
x'=(x_1,x_2,\dots,x_{N-1}),
\end{equation*}
并令
\begin{equation*}
\abs{x'}=\left(\sum_{i=1}^{N-1}x_i^2\right)^{1/2}.
\end{equation*}
我们定义
\begin{align*}
&\R^N_+=\{x=(x',x_N);\ x_N>0\},\\
&Q=\{x=(x',x_N);\ \abs{x'}<1\ \text{且}\ \abs{x_N}<1\},\\
&Q_+=Q\cap\R^N_+,\\
&Q_0=\{x=(x',0);\ \abs{x'}<1\}.
\end{align*}
\end{notation}

\begin{definition}
我们说开集 $\Omega$ 属于 $C^1$ 类, 如果对每个 $x\in\partial\Omega=\Gamma$, 存在
$x$ 在 $\R^N$ 中的邻域 $U$ 以及双射 $H:Q\to U$, 使得
\begin{equation*}
H\in C^1(Q),\quad H^{-1}\in C^1(U),\quad
H(Q_+)=U\cap\Omega,\quad H(Q_0)=U\cap\Gamma.
\end{equation*}
映射 $H$ 称为\emph{局部坐标卡}.
\end{definition}

\begin{theorem}
\index{yantuosuanzi@延拓算子 (extension operator)}
\label{thm:9.7}
设 $\Omega$ 属于 $C^1$ 类且 $\Gamma$ 有界 (或者 $\Omega=\R^N_+$).则存在线性
延拓算子
\begin{equation*}
P:W^{1,p}(\Omega)\to W^{1,p}(\R^N)\quad(1\le p\le\infty),
\end{equation*}
使得对一切 $u\in W^{1,p}(\Omega)$,
\begin{equation*}
\begin{array}{ll}
Pu|_\Omega=u, & \text{(i)}\\[2pt]
\norm{Pu}_{L^p(\R^N)}\le C\norm{u}_{L^p(\Omega)}, & \text{(ii)}\\[2pt]
\norm{Pu}_{W^{1,p}(\R^N)}\le C\norm{u}_{W^{1,p}(\Omega)}, & \text{(iii)}
\end{array}
\end{equation*}
其中 $C$ 只依赖于 $\Omega$.
\end{theorem}

我们先证明一个关于反射延拓的简单而基本的引理.

\begin{lemma}
\label{lem:9.2}
给定 $u\in W^{1,p}(Q_+)$, $1\le p\le\infty$, 定义 $Q$ 上的函数
$u^\star$ 为反射延拓, 即
\begin{equation*}
u^\star(x',x_N)=
\begin{cases}
u(x',x_N) & \text{若}\ x_N>0,\\
u(x',-x_N) & \text{若}\ x_N<0.
\end{cases}
\end{equation*}
则 $u^\star\in W^{1,p}(Q)$ 且
\begin{equation*}
\norm{u^\star}_{L^p(Q)}\le2\norm{u}_{L^p(Q_+)},\quad
\norm{u^\star}_{W^{1,p}(Q)}\le2\norm{u}_{W^{1,p}(Q_+)}.
\end{equation*}
\end{lemma}

\begin{proof}
事实上, 我们将证明
\begin{equation}
\frac{\partial u^\star}{\partial x_i}
=\left(\frac{\partial u}{\partial x_i}\right)^\star
\quad\text{对}\ 1\le i\le N-1
\label{eq:9.7}
\end{equation}
且
\begin{equation}
\frac{\partial u^\star}{\partial x_N}
=\left(\frac{\partial u}{\partial x_N}\right)^{\square},
\label{eq:9.8}
\end{equation}
其中 $(\frac{\partial u}{\partial x_i})^\star$ 表示 $\frac{\partial u}
{\partial x_i}$ 的反射延拓, 且只要 $f$ 定义在 $Q_+$ 上, 我们就令
\begin{equation*}
f^{\square}(x',x_N)=
\begin{cases}
f(x',x_N) & \text{若}\ x_N>0,\\
-f(x',-x_N) & \text{若}\ x_N<0.
\end{cases}
\end{equation*}
我们将使用 $C^\infty(\R)$ 中的一列函数 $(\eta_k)$, 它定义为
\begin{equation*}
\eta_k(t)=\eta(kt),\quad t\in\R,\ k=1,2,\dots,
\end{equation*}
其中 $\eta$ 是任一固定函数, $\eta\in C^\infty(\R)$, 满足
\begin{equation*}
\eta(t)=
\begin{cases}
0 & \text{若}\ t<1/2,\\
1 & \text{若}\ t>1.
\end{cases}
\end{equation*}

\noindent\emph{\eqref{eq:9.7} 的证明.} 设 $\varphi\in C_c^1(Q)$.对
$1\le i\le N-1$, 我们有
\begin{equation}
\int_Q u^\star\frac{\partial\varphi}{\partial x_i}
=\int_{Q_+}u\frac{\partial\psi}{\partial x_i},
\label{eq:9.9}
\end{equation}
其中 $\psi(x',x_N)=\varphi(x',x_N)+\varphi(x',-x_N)$.函数 $\psi$ 一般
不属于 $C_c^1(Q_+)$, 因此不能作为检验函数 (在 $W^{1,p}$ 的定义中).另一方
面, $\eta_k(x_N)\psi(x',x_N)\in C_c^1(Q_+)$, 于是
\begin{equation*}
\int_{Q_+}u\frac{\partial}{\partial x_i}(\eta_k\psi)
=-\int_{Q_+}\frac{\partial u}{\partial x_i}\eta_k\psi.
\end{equation*}
由于 $\frac{\partial}{\partial x_i}(\eta_k\psi)=\eta_k
\frac{\partial\psi}{\partial x_i}$, 我们有
\begin{equation}
\int_{Q_+}u\eta_k\frac{\partial\psi}{\partial x_i}
=-\int_{Q_+}\frac{\partial u}{\partial x_i}\eta_k\psi.
\label{eq:9.10}
\end{equation}
在 \eqref{eq:9.10} 中令 $k\to\infty$ 取极限 (利用控制收敛), 得到
\begin{equation}
\int_{Q_+}u\frac{\partial\psi}{\partial x_i}
=-\int_{Q_+}\frac{\partial u}{\partial x_i}\psi.
\label{eq:9.11}
\end{equation}
结合 \eqref{eq:9.9} 与 \eqref{eq:9.11}, 我们得到
\begin{equation*}
\int_Q u^\star\frac{\partial\varphi}{\partial x_i}
=-\int_{Q_+}\frac{\partial u}{\partial x_i}\psi
=-\int_Q\left(\frac{\partial u}{\partial x_i}\right)^\star\varphi,
\end{equation*}
由此推出 \eqref{eq:9.7}.

\noindent\emph{\eqref{eq:9.8} 的证明.} 对每个 $\varphi\in C_c^1(Q)$, 我们有
\begin{equation}
\int_Q u^\star\frac{\partial\varphi}{\partial x_N}
=\int_{Q_+}u\frac{\partial\chi}{\partial x_N},
\label{eq:9.12}
\end{equation}
其中 $\chi(x',x_N)=\varphi(x',x_N)-\varphi(x',-x_N)$.注意
$\chi(x',0)=0$, 于是存在常数 $M$ 使得
$\abs{\chi(x',x_N)}\le M\abs{x_N}$ 在 $Q$ 上成立.由于
$\eta_k\chi\in C_c^1(Q_+)$, 我们有
\begin{equation}
\int_{Q_+}u\frac{\partial}{\partial x_N}(\eta_k\chi)
=-\int_{Q_+}\frac{\partial u}{\partial x_N}\eta_k\chi.
\label{eq:9.13}
\end{equation}
但
\begin{equation}
\frac{\partial}{\partial x_N}(\eta_k\chi)=\eta_k\frac{\partial\chi}
{\partial x_N}+k\eta'(kx_N)\chi.
\label{eq:9.14}
\end{equation}
我们断言
\begin{equation}
\int_{Q_+}uk\eta'(kx_N)\chi\to0\quad\text{当}\ k\to\infty.
\label{eq:9.15}
\end{equation}
事实上, 我们有
\begin{equation*}
\left|\int_{Q_+}uk\eta'(kx_N)\chi\right|
\le kMC\int_{0<x_N<1/k}\abs{u}x_N\,dx
\le MC\int_{0<x_N<1/k}\abs{u}\,dx,
\end{equation*}
其中 $C=\sup_{t\in[0,1]}\abs{\eta'(t)}$, 由此推出 \eqref{eq:9.15}.

由 \eqref{eq:9.13}, \eqref{eq:9.14} 与 \eqref{eq:9.15} 推出
\begin{equation*}
\int_{Q_+}u\frac{\partial\chi}{\partial x_N}
=-\int_{Q_+}\frac{\partial u}{\partial x_N}\chi.
\end{equation*}
最后, 我们有
\begin{equation}
\int_{Q_+}\frac{\partial u}{\partial x_N}\chi
=\int_Q\left(\frac{\partial u}{\partial x_N}\right)^{\square}\varphi.
\label{eq:9.16}
\end{equation}
结合 \eqref{eq:9.12} 与 \eqref{eq:9.16}, 得到 \eqref{eq:9.8}.这就完成了
引理  \ref{lem:9.2} 的证明.
\end{proof}

引理  \ref{lem:9.2} 的结论在把 $Q_+$ 换成 $\R^N_+$ 后仍然成立 (证明不变).
这就建立了 $\Omega=\R^N_+$ 情形的定理  \ref{thm:9.7}.

\begin{starremark}
\label{rem:9.9}
引理  \ref{lem:9.2} 对某些不属于 $C^1$ 类的开集 $\Omega$ 给出了延拓算子的
非常简单的构造.例如考虑正方形
\begin{equation*}
\Omega=\{x\in\R^2;\ 0<x_1<1,\ 0<x_2<1\}.
\end{equation*}
设 $u\in W^{1,p}(\Omega)$.通过四次相继的反射 (见图  \ref{fig:6}), 我们得到 $u$ 在
\begin{equation*}
\widetilde\Omega=\{x\in\R^2;\ -1<x_1<3,\ -1<x_2<3\}
\end{equation*}
中的延拓 $\tilde u\in W^{1,p}(\widetilde\Omega)$.

\begin{figure}[tbp]
	\centering
	\begin{tikzpicture}[scale=0.95]
		\draw[dashed] (0,0) rectangle (4,4);
		\draw[dashed] (2,0) -- (2,4);
		\draw[dashed] (1,0) -- (1,2);
		\draw[dashed] (0,2) -- (2,2);
		\draw (1,1) rectangle (2,2);
		\node at (1.28,1.48) {$\Omega$};
		\node at (0.45,1.15) {$2$};
		\node at (1.35,0.55) {$1$};
		\node at (1.12,3.15) {$3$};
		\node at (2.85,2.15) {$4$};
	\end{tikzpicture}
	\caption{}
	\label{fig:6}
\end{figure}

然后固定任一函数 $\psi\in C_c^1(\widetilde\Omega)$ 使得 $\psi=1$ 在
$\Omega$ 上.用 $Pu$ 表示把 $\psi\tilde u$ 在 $\widetilde\Omega$ 外延拓为
$0$ 所得的 $\R^2$ 上的函数.容易证明算子
$P:W^{1,p}(\Omega)\to W^{1,p}(\R^2)$ 满足 (i), (ii) 与 (iii).
\end{starremark}

下述引理非常有用.

\begin{lemma}[单位分解]
\label{lem:9.3}
设 $K$ 为 $\R^N$ 的紧子集, $U_1,U_2,\dots,U_k$ 为 $K$ 的开覆盖, 即
$K\subset\bigcup_{i=1}^kU_i$.则存在函数
$\theta_0,\theta_1,\theta_2,\dots,\theta_k\in C^\infty(\R^N)$ 使得
\begin{equation*}
0\le\theta_i\le1\quad\forall i=0,1,2,\dots,k
\quad\text{且}\quad\sum_{i=0}^k\theta_i=1\ \text{在}\ \R^N\ \text{上},
\end{equation*}
\begin{equation*}
\begin{array}{ll}
\supp\theta_i\ \text{紧且}\ \supp\theta_i\subset U_i
\quad\forall i=1,2,\dots,k, & \text{(i)}\\[2pt]
\supp\theta_0\subset\R^N\setminus K. & \text{(ii)}
\end{array}
\end{equation*}
若 $\Omega$ 为有界开集且 $K=\partial\Omega$, 则
$\theta_0|_\Omega\in C_c^\infty(\Omega)$.
\end{lemma}

\begin{proof}
这个引理是经典的; 类似的陈述可在例如 S.~Agmon \bref{Agmon}{1}, R.~Adams \bref{Adams}{1}, G.~Folland \bref{Folland}{1}, P.~Malliavin \bref{Malliavin}{1} 中找到.
\end{proof}

\begin{proof}[定理 \ref{thm:9.7}]
我们用局部坐标卡把 $\Gamma=\partial\Omega$ ``拉直'', 并使用单位分解.\footnote{下文中我们将经常使用这个技巧, 把在 $\R^N_+$ (或 $Q_+$) 上证明
的结果转移到光滑开集 $\Omega$ 上的同样结论.}更精确地说, 由于 $\Gamma$ 紧
且属于 $C^1$ 类, 存在 $\R^N$ 中的开集 $(U_i)_{1\le i\le k}$ 使得
$\Gamma\subset\bigcup_{i=1}^kU_i$, 以及双射 $H_i:Q\to U_i$ 使得
\begin{equation*}
H_i\in C^1(Q),\quad H_i^{-1}\in C^1(U_i),\quad
H_i(Q_+)=U_i\cap\Omega,\quad H_i(Q_0)=U_i\cap\Gamma.
\end{equation*}
考虑引理  \ref{lem:9.3} 中引入的函数
$\theta_0,\theta_1,\theta_2,\dots,\theta_k$.给定
$u\in W^{1,p}(\Omega)$, 写
\begin{equation*}
u=\sum_{i=0}^k\theta_iu=\sum_{i=0}^ku_i,
\end{equation*}
其中 $u_i=\theta_iu$.

现在我们把每个函数 $u_i$ 延拓到 $\R^N$, 对 $u_0$ 与 $(u_i)_{1\le i\le k}$
分别处理.

\noindent (a) \emph{$u_0$ 的延拓.} 我们用
\begin{equation*}
\bar u_0(x)=
\begin{cases}
u_0(x) & \text{若}\ x\in\Omega,\\
0 & \text{若}\ x\in\R^N\setminus\Omega
\end{cases}
\end{equation*}
定义 $u_0$ 到 $\R^N$ 的延拓.回忆
$\theta_0\in C^1(\R^N)\cap L^\infty(\R^N)$,
$\nabla\theta_0\in L^\infty(\R^N)$, 因为
$\nabla\theta_0=-\sum_{i=1}^k\nabla\theta_i$ 有紧支集, 且
$\supp\theta_0\subset\R^N\setminus\Gamma$.由此推出 (由注记
\ref{rem:9.4}(ii)) $\bar u_0\in W^{1,p}(\R^N)$ 且
\begin{equation*}
\frac{\partial}{\partial x_i}\bar u_0=\theta_0
\frac{\partial u}{\partial x_i}+\frac{\partial\theta_0}{\partial x_i}\bar u.
\end{equation*}
于是
\begin{equation*}
\norm{\bar u_0}_{W^{1,p}(\R^N)}\le C\norm{u}_{W^{1,p}(\Omega)}.
\end{equation*}

\noindent (b) \emph{$u_i$ 的延拓, $1\le i\le k$.} 考虑 $u$ 在
$U_i\cap\Omega$ 上的限制, 并借助 $H_i$ 把这个函数``转移''到 $Q_+$ 上.
更精确地说, 令 $v_i(y)=u(H_i(y))$, $y\in Q_+$.我们知道 (命题
 \ref{prop:9.6}) $v_i\in W^{1,p}(Q_+)$.然后用反射 (引理  \ref{lem:9.2})
定义 $v_i$ 在 $Q$ 上的延拓; 记之为 $v_i^\star$.我们知道
$v_i^\star\in W^{1,p}(Q)$.用 $H_i^{-1}$ 把 $v_i^\star$ ``转回''到 $U_i$
上, 记之为 $w_i$:
\begin{equation*}
w_i(x)=v_i^\star[H_i^{-1}(x)]\quad\text{对}\ x\in U_i.
\end{equation*}
则 $w_i\in W^{1,p}(U_i)$, $w_i=u$ 在 $U_i\cap\Omega$ 上, 且
\begin{equation*}
\norm{w_i}_{W^{1,p}(U_i)}\le C\norm{u}_{W^{1,p}(U_i\cap\Omega)}.
\end{equation*}
最后, 对 $x\in\R^N$ 令
\begin{equation*}
\hat u_i(x)=
\begin{cases}
\theta_i(x)w_i(x) & \text{若}\ x\in U_i,\\
0 & \text{若}\ x\in\R^N\setminus U_i,
\end{cases}
\end{equation*}
于是 $\hat u_i\in W^{1,p}(\R^N)$ (见注记 \ref{rem:9.4}(ii)),
$\hat u_i=u_i$ 在 $\Omega$ 上, 且
\begin{equation*}
\norm{\hat u_i}_{W^{1,p}(\R^N)}\le C\norm{u}_{W^{1,p}(U_i\cap\Omega)}.
\end{equation*}

\noindent (c) \emph{结论.} 算子
$Pu=\bar u_0+\sum_{i=1}^N\hat u_i$ 具有所有所要的性质.
\end{proof}

\begin{keycorollary}[稠密性]
\label{cor:9.8}
设 $\Omega$ 属于 $C^1$ 类, $u\in W^{1,p}(\Omega)$, $1\le p<\infty$.则存在
$C_c^\infty(\R^N)$ 中的序列 $(u_n)$ 使得 $u_n|_\Omega\to u$ 在
$W^{1,p}(\Omega)$ 中.换言之, $C_c^\infty(\R^N)$ 中函数在 $\Omega$ 上的限制
构成 $W^{1,p}(\Omega)$ 的稠密子空间.
\end{keycorollary}

\begin{proof}
先设 $\Gamma$ 有界.则由定理  \ref{thm:9.7} 存在延拓算子 $P$.序列\footnote{按
惯例, $(\rho_n)$ 是一列磨光核, $(\zeta_n)$ 是一列截断函数, 如定理
 \ref{thm:9.2} 的证明中那样.}
$\zeta_n(\rho_n\star Pu)$ 在 $W^{1,p}(\R^N)$ 中收敛到 $Pu$, 于是它解决了
问题.当 $\Gamma$ 无界时, 我们先考虑序列 $\zeta_nu$.给定 $\varepsilon>0$,
固定 $n_0$ 使得 $\norm{\zeta_{n_0}u-u}_{W^{1,p}}<\varepsilon$.然后可以构造
$\zeta_{n_0}u$ 的延拓 $v\in W^{1,p}(\R^N)$ (因为这只涉及 $\Gamma$ 与一个
大球的交集).最后任取 $w\in C_c^\infty(\R^N)$ 使得
$\norm{w-v}_{W^{1,p}(\R^N)}<\varepsilon$.
\end{proof}

\section{Sobolev 不等式}
\label{sec:9.3}

在第~
\ref{chap:8}~章中我们看到, 若 $\Omega$ 是一维的, 则对一切 $1\le p\le\infty$,
$W^{1,p}(\Omega)\subset L^\infty(\Omega)$ 且嵌入连续.在维数 $N\ge2$ 时,
这个包含只对 $p>N$ 成立; 当 $p\le N$ 时可以构造出不属于 $L^\infty$ 的
$W^{1,p}$ 函数 (见注记 \ref{rem:9.16}).然而, 一个本质上属于 Sobolev 的
重要结果表明: 若 $1\le p<N$, 则 $W^{1,p}(\Omega)\subset L^{p^\star}(\Omega)$
且嵌入连续, 其中某个 $p^\star\in(p,+\infty)$.这个结果通常称为
\emph{Sobolev 嵌入定理}.我们先考虑下述情形:

\runin{A. $\Omega=\R^N$ 的情形.}

\begin{keytheorem}[Sobolev, Gagliardo, Nirenberg]
\index{Sobolevbudengshi@Sobolev 不等式 (Sobolev inequality)}
\label{thm:9.9}
设 $1\le p<N$.则
\begin{equation*}
W^{1,p}(\R^N)\subset L^{p^\star}(\R^N),
\end{equation*}
其中 $p^\star$ 由
\begin{equation*}
\frac1{p^\star}=\frac1p-\frac1N
\end{equation*}
给出, 且存在常数\footnote{可以取 $C(p,N)=(N-1)p/(N-p)$, 但这个常数不是最优
的.最优常数是已知的 (但它并不简单!), 参见 Th.~Aubin \bref{Aubin}{1}, G.~Talenti \bref{Talenti}{1} 与 E.~Lieb \bref{Lieb}{1}.}
$C=C(p,N)$ 使得
\begin{equation}
\norm{u}_{p^\star}\le C\norm{\nabla u}_p
\quad\forall u\in W^{1,p}(\R^N).
\label{eq:9.17}
\end{equation}
\end{keytheorem}

\begin{remark}
\label{rem:9.10}
$p^\star$ 的值可以通过一个非常简单的伸缩论证得到 (物理学家所钟爱的伸缩论证
有时以最少的努力给出有用的信息).事实上, 假设存在常数 $C$ 与 $q$
($1\le q\le\infty$) 使得
\begin{equation}
\norm{u}_q\le C\norm{\nabla u}_p\quad\forall u\in C_c^\infty(\R^N).
\label{eq:9.18}
\end{equation}
则必有 $q=p^\star$.为看出这一点, 固定任一函数
$u\in C_c^\infty(\R^N)$, 把 $u_\lambda(x)=u(\lambda x)$ 代入
\eqref{eq:9.18}.我们得到
\begin{equation*}
\norm{u}_q\le C\lambda^{\left(1+\frac Nq-\frac Np\right)}
\norm{\nabla u}_p\quad\forall\lambda>0,
\end{equation*}
由此推出 $1+\frac Nq-\frac Np=0$, 即 $q=p^\star$ (只要 $u$ 不恒为零).
\end{remark}

定理  \ref{thm:9.9} 的证明依赖于下述引理:

\begin{lemma}
\label{lem:9.4}
设 $N\ge2$, $f_1,f_2,\dots,f_N\in L^{N-1}(\R^{N-1})$.对 $x\in\R^N$ 与
$1\le i\le N$, 令
\begin{equation*}
\tilde x_i=(x_1,x_2,\dots,x_{i-1},x_{i+1},\dots,x_N)\in\R^{N-1},
\end{equation*}
即从列表中删去 $x_i$.则函数
\begin{equation*}
f(x)=f_1(\tilde x_1)f_2(\tilde x_2)\cdots f_N(\tilde x_N),\quad x\in\R^N,
\end{equation*}
属于 $L^1(\R^N)$, 且
\begin{equation*}
\norm{f}_{L^1(\R^N)}\le\prod_{i=1}^N\norm{f_i}_{L^{N-1}(\R^{N-1})}.
\end{equation*}
\end{lemma}

\begin{proof}
$N=2$ 的情形是平凡的 (为什么?).让我们考虑 $N=3$ 的情形.我们有
\begin{align*}
\int_\R\abs{f(x)}\,dx_3
&=\abs{f_3(x_1,x_2)}\int_\R\abs{f_1(x_2,x_3)}\abs{f_2(x_1,x_3)}\,dx_3\\
&\le\abs{f_3(x_1,x_2)}\left(\int_\R\abs{f_1(x_2,x_3)}^2dx_3\right)^{1/2}
\left(\int_\R\abs{f_2(x_1,x_3)}^2dx_3\right)^{1/2}
\end{align*}
(由 Cauchy--Schwarz).再用一次 Cauchy--Schwarz 得到
\begin{equation*}
\int_{\R^3}\abs{f(x)}\,dx\le\norm{f_3}_{L^2(\R^2)}
\norm{f_1}_{L^2(\R^2)}\norm{f_2}_{L^2(\R^2)}.
\end{equation*}
一般情形用归纳法得到——先假设结论对 $N$ 成立, 再推出它对 $N+1$ 成立.
固定 $x_{N+1}\in\R$; 由 Hölder 不等式,
\begin{align*}
&\int_{\R^N}\abs{f(x)}\,dx_1dx_2\cdots dx_N\\
&\quad\le\norm{f_{N+1}}_{L^N(\R^N)}
\left(\int\abs{f_1f_2\cdots f_N}^{N'}dx_1dx_2\cdots dx_N\right)^{1/N'}
\end{align*}
(其中 $N'=N/(N-1)$).对函数 $\abs{f_1}^{N'},\abs{f_2}^{N'},\dots,
\abs{f_N}^{N'}$ 应用归纳假设, 我们得到
\begin{equation*}
\int_{\R^N}\abs{f_1}^{N'}\cdots\abs{f_N}^{N'}\,dx_1\cdots dx_N
\le\prod_{i=1}^N\norm{f_i}_{L^N(\R^{N-1})}^{N'},
\end{equation*}
由此推出
\begin{equation*}
\int_{\R^N}\abs{f(x)}\,dx_1\cdots dx_N
\le\norm{f_{N+1}}_{L^N(\R^N)}\prod_{i=1}^N
\norm{f_i}_{L^N(\R^{N-1})}.
\end{equation*}
现在让 $x_{N+1}$ 变化.每个函数
$x_{N+1}\mapsto\norm{f_i}_{L^N(\R^{N-1})}$ 都属于 $L^N(\R)$,
$1\le i\le N$.因此它们的乘积
$\prod_{i=1}^N\norm{f_i}_{L^N(\R^{N-1})}$ 属于 $L^1(\R)$ (见第~
\ref{chap:4}~章中
Hölder 不等式后的注记 2), 并且
\begin{equation*}
\int_{\R^{N+1}}\abs{f(x)}\,dx_1dx_2\cdots dx_Ndx_{N+1}
\le\prod_{i=1}^{N+1}\norm{f_i}_{L^N(\R^N)}.
\end{equation*}
\end{proof}

\begin{proof}[定理 \ref{thm:9.9}]
我们从 $p=1$ 且 $u\in C_c^1(\R^N)$ 的情形开始.我们有
\begin{align*}
\abs{u(x_1,x_2,\dots,x_N)}
&=\left|\int_{-\infty}^{x_1}\frac{\partial u}{\partial x_1}
(t,x_2,\dots,x_N)\,dt\right|\\
&\le\int_{-\infty}^{+\infty}\left|\frac{\partial u}{\partial x_1}
(t,x_2,\dots,x_N)\right|dt,
\end{align*}
类似地, 对每个 $1\le i\le N$,
\begin{equation*}
\abs{u(x_1,x_2,\dots,x_N)}\le\int_{-\infty}^{+\infty}
\left|\frac{\partial u}{\partial x_i}(x_1,x_2,\dots,x_{i-1},t,
x_{i+1},\dots,x_N)\right|dt
\overset{\text{def}}{=}f_i(\tilde x_i).
\end{equation*}
于是
\begin{equation*}
\abs{u(x)}^N\le\prod_{i=1}^Nf_i(\tilde x_i).
\end{equation*}
由引理  \ref{lem:9.4} 推出
\begin{equation*}
\int_{\R^N}\abs{u(x)}^{N/(N-1)}dx\le\prod_{i=1}^N
\norm{f_i}_{L^1(\R^{N-1})}^{1/(N-1)}
=\prod_{i=1}^N\left\|\frac{\partial u}{\partial x_i}\right\|_{L^1(\R^N)}^{1/(N-1)}.
\end{equation*}
因此我们有
\begin{equation}
\norm{u}_{L^{N/(N-1)}(\R^N)}\le\prod_{i=1}^N
\left\|\frac{\partial u}{\partial x_i}\right\|_{L^1(\R^N)}^{1/N}.
\label{eq:9.19}
\end{equation}
这就完成了 $p=1$ 且 $u\in C_c^1(\R^N)$ 情形下 \eqref{eq:9.17} 的证明.
现在转向 $1<p<N$ 的情形, 仍设 $u\in C_c^1(\R^N)$.设 $m\ge1$; 把
\eqref{eq:9.19} 用于 $\abs{u}^{m-1}u$ 而不是 $u$.我们得到
\begin{align}
\norm{u}_{mN/(N-1)}^m
&\le m\prod_{i=1}^N\left\|\abs{u}^{m-1}
\frac{\partial u}{\partial x_i}\right\|_1^{1/N}\notag\\
&\le m\norm{u}_{p'(m-1)}^{m-1}\prod_{i=1}^N
\left\|\frac{\partial u}{\partial x_i}\right\|_p^{1/N}.
\label{eq:9.20}
\end{align}
然后取 $m$ 使得 $mN/(N-1)=p'(m-1)$, 这给出
$m=(N-1)p^\star/N$ (由于 $1<p<N$, 有 $m\ge1$).我们得到
\begin{equation*}
\norm{u}_{p^\star}\le m\prod_{i=1}^N
\left\|\frac{\partial u}{\partial x_i}\right\|_p^{1/N},
\end{equation*}
从而
\begin{equation*}
\norm{u}_{p^\star}\le C\norm{\nabla u}_p\quad\forall u\in C^1(\R^N).
\end{equation*}
为完成证明, 设 $u\in W^{1,p}(\R^N)$, 并取 $C_c^1(\R^N)$ 中的序列
$(u_n)$ 使得 $u_n\to u$ 在 $W^{1,p}(\R^N)$ 中.必要时抽取子序列, 还可以假设
$u_n\to u$ 几乎处处成立.我们有
\begin{equation*}
\norm{u_n}_{p^\star}\le C\norm{\nabla u_n}_p.
\end{equation*}
由 Fatou 引理\footnote{也可以注意到 $(u_n)$ 是 $L^{p^\star}$ 中的 Cauchy
序列来得出结论.}推出
\begin{equation*}
u\in L^{p^\star}\quad\text{且}\quad\norm{u}_{p^\star}\le C\norm{\nabla u}_p.\textquoteleft
\end{equation*}
% 原书此处不等式末尾多印了一个孤立的左单引号 (见印刷页 281); 依技能规范
% 保留原样并记录, 不擅自更正.
\end{proof}

\begin{keycorollary}
\label{cor:9.10}
设 $1\le p<N$.则
\begin{equation*}
W^{1,p}(\R^N)\subset L^q(\R^N)\quad\forall q\in[p,p^\star]
\end{equation*}
且嵌入连续.
\end{keycorollary}

\begin{proof}
给定 $q\in[p,p^\star]$, 我们写
\begin{equation*}
\frac1q=\frac\alpha p+\frac{1-\alpha}{p^\star}
\quad\text{其中}\ \alpha\in[0,1].
\end{equation*}
我们知道 (见第~
\ref{chap:4}~章注记 2)
\begin{equation*}
\norm{u}_q\le\norm{u}_p^\alpha\norm{u}_{p^\star}^{1-\alpha}
\le\norm{u}_p+\norm{u}_{p^\star}
\end{equation*}
(由 Young 不等式).利用定理  \ref{thm:9.9}, 我们推出
\begin{equation*}
\norm{u}_q\le C\norm{u}_{W^{1,p}}\quad\forall u\in W^{1,p}(\R^N).
\end{equation*}
\end{proof}

\begin{keycorollary}[极限情形 $p=N$]
\label{cor:9.11}
我们有
\begin{equation*}
W^{1,p}(\R^N)\subset L^q(\R^N)\quad\forall q\in[N,+\infty).
\end{equation*}
\end{keycorollary}

\begin{proof}
设 $u\in C_c^1(\R^N)$; 把 \eqref{eq:9.20} 用于 $p=N$, 我们得到
\begin{equation*}
\norm{u}_{mN/(N-1)}^m\le m\norm{u}_{(m-1)N/(N-1)}^{m-1}
\norm{\nabla u}_N\quad\forall m\ge1,
\end{equation*}
借助 Young 不等式, 我们有
\begin{equation}
\norm{u}_{mN/(N-1)}\le C\bigl(\norm{u}_{(m-1)N/(N-1)}
+\norm{\nabla u}_N\bigr)\quad\forall m\ge1.
\label{eq:9.21}
\end{equation}
在 \eqref{eq:9.21} 中先取 $m=N$; 它成为
\begin{equation*}
\norm{u}_{N^2/(N-1)}\le C\norm{u}_{W^{1,N}},
\end{equation*}
由插值不等式 (见第~
\ref{chap:4}~章注记 2) 我们有
\begin{equation}
\norm{u}_q\le C\norm{u}_{W^{1,N}}
\label{eq:9.22}
\end{equation}
对一切满足 $N\le q\le N^2/(N-1)$ 的 $q$ 成立.用 $m=N+1$, $m=N+2$ 等
反复迭代这个论证, 我们得到
\begin{equation}
\norm{u}_q\le C\norm{u}_{W^{1,N}}\quad\forall u\in C^1(\R^N)
\label{eq:9.23}
\end{equation}
对一切 $q\in[N,+\infty)$ 成立, 其中常数 $C$ 依赖于 $q$ 与 $N$.\footnote{当
$q\to+\infty$ 时这个常数``爆破''.}不等式 \eqref{eq:9.23} 由稠密性延拓到
$W^{1,N}$.
\end{proof}

\begin{keytheorem}[Morrey]
\index{Morreybudengshi@Morrey 不等式 (Morrey inequality)}
\label{thm:9.12}
设 $p>N$.则
\begin{equation}
W^{1,p}(\R^N)\subset L^\infty(\R^N)
\label{eq:9.24}
\end{equation}
且嵌入连续.此外, 对一切 $u\in W^{1,p}(\R^N)$, 我们有
\begin{equation}
\abs{u(x)-u(y)}\le C\abs{x-y}^\alpha\norm{\nabla u}_p
\quad\text{对几乎一切}\ x,y\in\R^N,
\label{eq:9.25}
\end{equation}
其中 $\alpha=1-(N/p)$, $C$ 为常数 (只依赖于 $p$ 与 $N$).
\end{keytheorem}

\begin{remark}
\label{rem:9.11}
不等式 \eqref{eq:9.25} 蕴含存在函数 $\tilde u\in C(\R^N)$ 使得
$u=\tilde u$ 在 $\R^N$ 上几乎处处成立.(事实上, 设 $A\subset\R^N$ 为零测集
使得 \eqref{eq:9.25} 对 $x,y\in\R^N\setminus A$ 成立; 由于
$\R^N\setminus A$ 在 $\R^N$ 中稠密, 函数 $u|_{\R^N\setminus A}$ 有 (唯一的)
连续延拓到 $\R^N$.)换言之, 每个 $u\in W^{1,p}(\R^N)$ ($p>N$) 都有连续
代表元.在有用的时候, 我们用它的连续代表元代替 $u$, 并且也用 $u$ 表示这个
连续代表元.
\end{remark}

\begin{proof}[定理 \ref{thm:9.12}]
我们先对 $u\in C_c^1(\R^N)$ 建立 \eqref{eq:9.25}.设 $Q$ 为包含 $0$ 的开
立方体, 其边——长为 $r$——平行于坐标轴.对 $x\in Q$ 我们有
\begin{equation*}
u(x)-u(0)=\int_0^1\frac{d}{dt}u(tx)\,dt,
\end{equation*}
于是
\begin{equation}
\abs{u(x)-u(0)}\le\int_0^1\sum_{i=1}^N\abs{x_i}
\left|\frac{\partial u}{\partial x_i}(tx)\right|dt
\le r\sum_{i=1}^N\int_0^1\left|\frac{\partial u}{\partial x_i}(tx)\right|dt.
\label{eq:9.26}
\end{equation}
令
\begin{equation*}
\bar u=\frac{1}{\abs{Q}}\int_Qu(x)\,dx=(\text{$u$ 在 $Q$ 上的平均值}).
\end{equation*}
在 $Q$ 上积分 \eqref{eq:9.26}, 我们得到
\begin{align*}
\abs{\bar u-u(0)}
&\le\frac{r}{\abs{Q}}\int_Qdx\sum_{i=1}^N\int_0^1
\left|\frac{\partial u}{\partial x_i}(tx)\right|dt\\
&=\frac{1}{r^{N-1}}\int_0^1dt\int_Q\sum_{i=1}^N
\left|\frac{\partial u}{\partial x_i}(tx)\right|dx\\
&=\frac{1}{r^{N-1}}\int_0^1dt\int_{tQ}\sum_{i=1}^N
\left|\frac{\partial u}{\partial x_i}(y)\right|\frac{dy}{t^N}.
\end{align*}
于是由 Hölder 不等式, 我们有
\begin{equation*}
\int_{tQ}\left|\frac{\partial u}{\partial x_i}(y)\right|dy
\le\left(\int_Q\left|\frac{\partial u}{\partial x_i}\right|^p\right)^{1/p}
\abs{tQ}^{1/p'}
\end{equation*}
(因为当 $t\in(0,1)$ 时 $tQ\subset Q$).由此推出
\begin{equation*}
\abs{\bar u-u(0)}\le\frac{1}{r^{N-1}}\norm{\nabla u}_{L^p(Q)}
r^{N/p'}\int_0^1\frac{t^{N/p'}}{t^N}\,dt
=\frac{r^{1-(N/p)}}{1-(N/p)}\norm{\nabla u}_{L^p(Q)}.
\end{equation*}
由平移, 这个不等式对一切边——长为 $r$——平行于坐标轴的立方体 $Q$ 仍然成立.
于是我们有
\begin{equation}
\abs{\bar u-u(x)}\le\frac{r^{1-(N/p)}}{1-(N/p)}\norm{\nabla u}_{L^p(Q)}
\quad\forall x\in Q.
\label{eq:9.27}
\end{equation}
把这两个不等式相加 (并利用三角不等式), 我们得到
\begin{equation}
\abs{u(x)-u(y)}\le\frac{2r^{1-(N/p)}}{1-(N/p)}
\norm{\nabla u}_{L^p(Q)}\quad\forall x,y\in Q.
\label{eq:9.28}
\end{equation}
给定 $\R^N$ 中任意两点 $x,y$, 存在这样的立方体 $Q$, 其边长
$r=2\abs{x-y}$ 且包含 $x$ 与 $y$.这就对 $u\in C_c^1(\R^N)$ 推出了
\eqref{eq:9.25}.对一般函数 $u\in W^{1,p}(\R^N)$, 我们使用
$C_c^1(\R^N)$ 中的序列 $(u_n)$ 使得 $u_n\to u$ 在 $W^{1,p}(\R^N)$ 中且
$u_n\to u$ 几乎处处成立.

现在我们证明 \eqref{eq:9.24}.设 $u\in C_c^1(\R^N)$, $x\in\R^N$, $Q$ 为
边长为 $r=1$ 且包含 $x$ 的立方体.由 \eqref{eq:9.27} 与 Hölder 不等式,
我们有
\begin{equation*}
\abs{u(x)}\le\abs{\bar u}+C\norm{\nabla u}_{L^p(Q)}
\le C\norm{u}_{W^{1,p}(Q)}\le C\norm{u}_{W^{1,p}(\R^N)},
\end{equation*}
其中 $C$ 只依赖于 $p$ 与 $N$.于是
\begin{equation*}
\norm{u}_{L^\infty(\R^N)}\le C\norm{u}_{W^{1,p}(\R^N)}
\quad\forall u\in C_c^1(\R^N).
\end{equation*}
对一般函数 $u\in W^{1,p}(\R^N)$, 我们使用标准的稠密性论证.
\end{proof}

\begin{remark}
\label{rem:9.12}
由 \eqref{eq:9.24} 推出: 若 $u\in W^{1,p}(\R^N)$, $N<p<\infty$, 则
\begin{equation*}
\lim_{\abs{x}\to\infty}u(x)=0.
\end{equation*}
事实上, 存在 $C_c^1(\R^N)$ 中的序列 $(u_n)$ 使得 $u_n\to u$ 在
$W^{1,p}(\R^N)$ 中.由 \eqref{eq:9.24}, $u$ 也是 $u_n$ 在 $\R^N$ 上的
一致极限.
\end{remark}

\begin{keycorollary}
\label{cor:9.13}
设 $m\ge1$ 为整数, $p\in[1,+\infty)$.我们有
\begin{equation*}
W^{m,p}(\R^N)\subset L^q(\R^N),
\end{equation*}
其中
\begin{equation*}
\frac1q=\frac1p-\frac mN\quad\text{若}\ \frac1p-\frac mN>0,
\end{equation*}
\begin{equation*}
W^{m,p}(\R^N)\subset L^q(\R^N)\quad\forall q\in[p,+\infty)
\quad\text{若}\ \frac1p-\frac mN=0,
\end{equation*}
\begin{equation*}
W^{m,p}(\R^N)\subset L^\infty(\R^N)\quad\text{若}\
\frac1p-\frac mN<0,
\end{equation*}
且所有这些嵌入都连续.此外, 若 $m-(N/p)>0$ 不是整数, 令\footnote{$[\ ]$
表示整数部分.}
\begin{equation*}
k=[m-(N/p)]\quad\text{且}\quad
\theta=m-(N/p)-k\ (0<\theta<1).
\end{equation*}
则对一切 $u\in W^{m,p}(\R^N)$,
\begin{equation*}
\norm{D^\alpha u}_{L^\infty(\R^N)}\le C\norm{u}_{W^{m,p}(\R^N)}
\quad\forall\alpha,\ \abs{\alpha}\le k,
\end{equation*}
并且\footnote{这蕴含对一切 $\abs{\alpha}<k$, $D^\alpha u$ 是 Lipschitz
连续的, 即
\begin{equation*}
\abs{D^\alpha u(x)-D^\alpha u(y)}\le C\norm{u}_{W^{m,p}}\abs{x-y}
\quad\text{对几乎一切}\ x,y\in\R^N.
\end{equation*}}
\begin{equation*}
\abs{D^\alpha u(x)-D^\alpha u(y)}\le C\norm{u}_{W^{m,p}(\R^N)}
\abs{x-y}^\theta\quad\text{对几乎一切}\ x,y\in\R^N,\
\forall\alpha,\ \abs{\alpha}=k.
\end{equation*}
特别地, $W^{m,p}(\R^N)\subset C^k(\R^N)$.\footnote{这应当理解为在连续
代表元的选择意义下成立.}
\end{keycorollary}

\begin{proof}
所有这些结果都由定理  \ref{thm:9.9}, 推论  \ref{cor:9.11} 与定理
 \ref{thm:9.12} 反复应用得到.
\end{proof}

\begin{remark}
\label{rem:9.13}
$p=1$ 且 $m=N$ 的情形是特殊的.我们有 $W^{N,1}\subset L^\infty$.
(但一般说来, 当 $p>1$ 且 $m=N/p$ 时 $W^{m,p}\subset L^\infty$ 并不成立.)
事实上, 对 $u\in C_c^\infty(\R^N)$, 我们有
\begin{equation*}
u(x_1,x_2,\dots,x_N)=\int_{-\infty}^{x_1}\int_{-\infty}^{x_2}
\cdots\int_{-\infty}^{x_N}\frac{\partial^Nu}
{\partial x_1\partial x_2\cdots\partial x_N}(t_1,t_2,\dots,t_N)
\,dt_1dt_2\cdots dt_N,
\end{equation*}
于是
\begin{equation}
\norm{u}_{L^\infty}\le C\norm{u}_{W^{N,1}}
\quad\forall u\in C_c^\infty(\R^N).
\label{eq:9.29}
\end{equation}
一般函数 $u\in W^{N,1}$ 的情形由稠密性推出.
\end{remark}

现在让我们转向下述情形.

\runin{B. $\Omega\subset\R^N$ 的情形.}

这里我们假设 $\Omega$ 是 $C^1$ 类开集且 $\Gamma$ 有界, 或者 $\Omega=\R^N_+$.

\begin{keycorollary}
\label{cor:9.14}
设 $1\le p\le\infty$.我们有
\begin{equation*}
W^{1,p}(\Omega)\subset L^{p^\star}(\Omega),
\end{equation*}
其中
\begin{equation*}
\frac1{p^\star}=\frac1p-\frac1N,\quad\text{若}\ p<N,
\end{equation*}
\begin{equation*}
W^{1,p}(\Omega)\subset L^q(\Omega)\quad\forall q\in[p,+\infty),
\quad\text{若}\ p=N,
\end{equation*}
\begin{equation*}
W^{1,p}(\Omega)\subset L^\infty(\Omega),\quad\text{若}\ p>N,
\end{equation*}
且所有这些嵌入都连续.此外, 若 $p>N$, 则对一切
$u\in W^{1,p}(\Omega)$,
\begin{equation*}
\abs{u(x)-u(y)}\le C\norm{u}_{W^{1,p}}\abs{x-y}^\alpha
\quad\text{对几乎一切}\ x,y\in\Omega,
\end{equation*}
其中 $\alpha=1-(N/p)$, $C$ 只依赖于 $\Omega$, $p$ 与 $N$.特别地,
$W^{1,p}(\Omega)\subset C(\overline\Omega)$.\footnote{再强调一次, 这应当
理解为在连续代表元的选择意义下成立.}
\end{keycorollary}

\begin{proof}
考虑延拓算子 $P:W^{1,p}(\Omega)\to W^{1,p}(\R^N)$ (见定理
 \ref{thm:9.7}), 然后应用定理  \ref{thm:9.9}, 推论  \ref{cor:9.11} 与定理
 \ref{thm:9.12}.
\end{proof}

\begin{keycorollary}
\label{cor:9.15}
把 $\R^N$ 换成 $\Omega$ 后, 推论  \ref{cor:9.13} 的结论仍然成立.\footnote{确切
地说, 若 $m-(N/p)>0$ 不是整数, 则 $W^{m,p}(\Omega)\subset C^k(\overline\Omega)$,
其中 $k=[m-(N/p)]$, 且
\begin{equation*}
C^k(\overline\Omega)=\{u\in C^k(\Omega);\ D^\alpha u\ \text{对一切}
\ \abs{\alpha}\le k\ \text{在}\ \overline\Omega\ \text{上有连续延拓}\}.
\end{equation*}}
\end{keycorollary}

\begin{proof}
反复应用推论  \ref{cor:9.14}.\footnote{另一种做法是把推论  \ref{cor:9.13}
与延拓算子 $P:W^{m,p}(\Omega)\to W^{m,p}(\R^N)$ 结合起来使用, 但这需要
一个额外的假设: 为构造这个延拓算子, $\Omega$ 必须属于 $C^m$ 类.}
\end{proof}

\begin{keytheorem}[Rellich--Kondrachov]
\label{thm:9.16}
设 $\Omega$ 有界且属于 $C^1$ 类.则下述嵌入是紧的:
\begin{equation*}
W^{1,p}(\Omega)\subset L^q(\Omega)\quad\forall q\in[1,p^\star),
\end{equation*}
其中
\begin{equation*}
\frac1{p^\star}=\frac1p-\frac1N,\quad\text{若}\ p<N,
\end{equation*}
\begin{equation*}
W^{1,p}(\Omega)\subset L^q(\Omega)\quad\forall q\in[p,+\infty),
\quad\text{若}\ p=N,
\end{equation*}
\begin{equation*}
W^{1,p}(\Omega)\subset C(\overline\Omega),\quad\text{若}\ p>N.
\end{equation*}
特别地, 对一切 $p$ (以及一切 $N$), $W^{1,p}(\Omega)\subset L^p(\Omega)$
是紧嵌入.
\end{keytheorem}

\begin{proof}
$p>N$ 的情形由推论  \ref{cor:9.14} 与 Ascoli--Arzelà 定理推出.$p=N$ 的
情形归结为 $p<N$ 的情形.因此只剩下 $p<N$ 的情形.

设 $H$ 为 $W^{1,p}(\Omega)$ 中的单位球.设 $P$ 为定理  \ref{thm:9.7} 的延拓
算子.令 $F=P(H)$, 于是 $H=F|_\Omega$.为证明对 $q\in[1,p^\star)$, $H$
在 $L^q(\Omega)$ 中有紧闭包, 我们引用定理  \ref{thm:4.26}.由于 $\Omega$
有界, 我们总可以假设 $q\ge p$.显然 $F$ 在 $W^{1,p}(\R^N)$ 中有界, 因此
由推论  \ref{cor:9.10} 它在 $L^q(\R^N)$ 中也有界.我们要验证
\begin{equation*}
\lim_{\abs{h}\to0}\norm{\tau_hf-f}_{L^q(\R^N)}=0
\quad\text{关于}\ f\in F\ \text{一致}.
\end{equation*}
由命题  \ref{prop:9.3} 我们有
\begin{equation*}
\norm{\tau_hf-f}_{L^p(\R^N)}\le\abs{h}\norm{\nabla f}_{L^p(\R^N)}
\quad\forall f\in F.
\end{equation*}
由于 $p\le q<p^\star$, 我们可以写
\begin{equation*}
\frac1q=\frac\alpha p+\frac{1-\alpha}{p^\star}
\quad\text{其中}\ \alpha\in(0,1].
\end{equation*}
借助插值不等式 (见第~
\ref{chap:4}~章注记 2), 我们有
\begin{align*}
\norm{\tau_hf-f}_{L^q(\R^N)}
&\le\norm{\tau_hf-f}_{L^p(\R^N)}^\alpha
\norm{\tau_hf-f}_{L^{p^\star}(\R^N)}^{1-\alpha}\\
&\le\abs{h}^\alpha\norm{\nabla f}_{L^p(\R^N)}^\alpha
\bigl(2\norm{f}_{L^{p^\star}(\R^N)}\bigr)^{1-\alpha}\le C\abs{h}^\alpha,
\end{align*}
其中 $C$ 与 $F$ 无关 (因为 $F$ 在 $W^{1,p}(\R^N)$ 中有界).所要结论由此
推出.
\end{proof}

\begin{remark}
\label{rem:9.14}
定理  \ref{thm:9.16} 在下述意义下是``几乎最优的'':

\noindent (i) 若 $\Omega$ 无界, 嵌入 $W^{1,p}(\Omega)\subset L^p(\Omega)$
一般不是紧的.\footnote{关于无界区域上这个嵌入的紧性, 参见 R.~Adams
\bref{Adams}{1} 第 167 页的详细讨论.}

\noindent (ii) 嵌入 $W^{1,p}(\Omega)\subset L^{p^\star}(\Omega)$ 从来不是
紧的, 即使 $\Omega$ 有界且光滑.
\end{remark}

\begin{starremark}
\label{rem:9.15}
设 $\Omega$ 为 $C^1$ 类有界开集.则只要
\begin{equation*}
1\le q\le p^\star\ \text{若}\ 1\le p<N,\quad
1\le q<\infty\ \text{若}\ p=N,\quad
1\le q\le\infty\ \text{若}\ p>N,
\end{equation*}
范数
\begin{equation*}
|||u|||=\norm{\nabla u}_p+\norm{u}_q
\end{equation*}
就与 $W^{1,p}$ 范数等价.
\end{starremark}

\begin{starremark}[极限情形 $p=N$]
\label{rem:9.16}
设 $\Omega$ 为 $C^1$ 类有界开集, $u\in W^{1,N}(\Omega)$.则一般说来
$u\notin L^\infty(\Omega)$.例如, 若
\begin{equation*}
\Omega=\{x\in\R^N;\ \abs{x}<1/2\},
\end{equation*}
则函数
\begin{equation*}
u(x)=(\log 1/\abs{x})^\alpha
\quad\text{其中}\ 0<\alpha<1-(1/N)
\end{equation*}
属于 $W^{1,N}(\Omega)$, 但由于 $x=0$ 处的奇异性它不是有界的.尽管如此, 我们
有 Trudinger 不等式
\begin{equation*}
\int_\Omega e^{\abs{u}^{N/(N-1)}}<\infty\quad\forall u\in W^{1,N}(\Omega)
\end{equation*}
(参见例如 R.~Adams \bref{Adams}{1} 或 D.~Gilbarg--N.~Trudinger
\bref{GilbargTrudinger}{1}).
\end{starremark}

\section{空间 $W_0^{1,p}(\Omega)$}
\label{sec:9.4}

\begin{definition}
设 $1\le p<\infty$; $W_0^{1,p}(\Omega)$ 表示 $C_c^1(\Omega)$ 在
$W^{1,p}(\Omega)$ 中的闭包.令\footnote{有歧义时, 我们将写 $W_0^{1,p}$, $H_0^1$ 而不写 $W_0^{1,p}(\Omega)$, $H_0^1(\Omega)$.}
\begin{equation*}
H_0^1(\Omega)=W_0^{1,2}(\Omega).
\end{equation*}
空间 $W_0^{1,p}$ 配以 $W^{1,p}$ 范数是可分 Banach 空间; 当 $1<p<\infty$ 时
它自反.$H_0^1$ 配以 $H^1$ 标量积是 Hilbert 空间.
\end{definition}

\begin{starremark}
\label{rem:9.17}
由于 $C_c^1(\R^N)$ 在 $W^{1,p}(\R^N)$ 中稠密, 我们有
\begin{equation*}
W_0^{1,p}(\R^N)=W^{1,p}(\R^N).
\end{equation*}
相反, 若 $\Omega\subset\R^N$ 且 $\Omega\ne\R^N$, 则一般说来
$W_0^{1,p}(\Omega)\ne W^{1,p}(\Omega)$.然而, 若 $\R^N\setminus\Omega$
``足够薄''且 $p<N$, 则 $W_0^{1,p}(\Omega)=W^{1,p}(\Omega)$.例如, 若
$\Omega=\R^N\setminus\{0\}$ 且 $N\ge2$, 则可以证明
$H_0^1(\Omega)=H^1(\Omega)$.
\end{starremark}

\begin{remark}
\label{rem:9.18}
利用一列磨光核容易验证 $C_c^\infty(\Omega)$ 在 $W_0^{1,p}(\Omega)$ 中
稠密.换言之, 在 $W_0^{1,p}(\Omega)$ 的定义中本来同样可以用
$C_c^\infty(\Omega)$ 代替 $C_c^1(\Omega)$.

$W_0^{1,p}(\Omega)$ 中的函数``大致''就是那些``在 $\Gamma=\partial\Omega$
上为零''的 $W^{1,p}(\Omega)$ 函数.要确切说明这一点是微妙的, 因为函数
$u\in W^{1,p}(\Omega)$ 只在几乎处处意义下定义 (而 $\Gamma$ 的测度为零!),
并且 $u$ 未必有连续代表元.\footnote{尽管如此, 若
$u\in W^{1,p}(\Omega)$ (当 $\Omega$ 正则时), 可以赋予 $u|_\Gamma$ 一个
意义, 并且可以证明例如 $u|_\Gamma\in L^p(\Gamma)$.这依赖于迹理论 (见本章
末的评注).}下述刻画表明我们``确实''得到了``在 $\Gamma$ 上为零''的函数.
我们先从一个简单事实开始:
\end{remark}

\begin{lemma}
\label{lem:9.5}
设 $u\in W^{1,p}(\Omega)$, $1\le p<\infty$, 并假设 $\supp u$ 是 $\Omega$
的紧子集.则 $u\in W_0^{1,p}(\Omega)$.
\end{lemma}

\begin{proof}
固定开集 $\omega$ 使得 $\supp u\subset\omega\subset\subset\Omega$, 并选取
$\alpha\in C_c^1(\omega)$ 使得 $\alpha=1$ 在 $\supp u$ 上; 于是
$\alpha u=u$.另一方面 (定理  \ref{thm:9.2}), 存在 $C_c^\infty(\R^N)$ 中的
序列 $(u_n)$ 使得 $u_n\to u$ 在 $L^p(\Omega)$ 中且
$\nabla u_n\to\nabla u$ 在 $L^p(\omega)^N$ 中.由此推出
$\alpha u_n\to\alpha u$ 在 $W^{1,p}(\Omega)$ 中.于是
$\alpha u$ 属于 $W_0^{1,p}(\Omega)$, $u$ 也如此.
\end{proof}

\begin{theorem}
\label{thm:9.17}
设 $\Omega$ 属于 $C^1$ 类.设\footnote{回忆若 $p>N$, 则
$u\in W^{1,p}(\Omega)\Rightarrow u\in C(\overline\Omega)$ (见推论
 \ref{cor:9.14}).}
\begin{equation*}
u\in W^{1,p}(\Omega)\cap C(\overline\Omega),
\end{equation*}
其中 $1\le p<\infty$.则下述性质等价:

\noindent (i) $u=0$ 在 $\Gamma$ 上;

\noindent (ii) $u\in W_0^{1,p}(\Omega)$.
\end{theorem}

\begin{proof}
\noindent (i)$\Rightarrow$(ii). 先设 $\supp u$ 有界.固定函数
$G\in C^1(\R)$ 使得
\begin{equation*}
\abs{G(t)}\le\abs{t}\quad\forall t\in\R
\quad\text{且}\quad
G(t)=
\begin{cases}
0 & \text{若}\ \abs{t}\le1,\\
t & \text{若}\ \abs{t}\ge2.
\end{cases}
\end{equation*}
则 $u_n=(1/n)G(nu)$ 属于 $W^{1,p}$ (由命题  \ref{prop:9.5}).容易验证
(利用控制收敛) $u_n\to u$ 在 $W^{1,p}$ 中.另一方面,
\begin{equation*}
\supp u_n\subset\{x\in\Omega;\ \abs{u(x)}\ge1/n\},
\end{equation*}
于是 $\supp u_n$ 是含于 $\Omega$ 的紧集.由引理  \ref{lem:9.5},
$u_n\in W_0^{1,p}$, 由此推出 $u\in W_0^{1,p}$.在 $\supp u$ 无界的一般
情形, 考虑序列 $(\zeta_nu)$ (其中 $(\zeta_n)$ 是定理  \ref{thm:9.2} 的证明
中的一列截断函数).由上述结果 $\zeta_nu\in W_0^{1,p}$, 且由于
$\zeta_nu\to u$ 在 $W^{1,p}$ 中, 我们推出 $u\in W_0^{1,p}$.

\noindent (ii)$\Rightarrow$(i). 利用局部坐标卡, 这归结为下述问题.设
$u\in W_0^{1,p}(Q_+)\cap C(\overline{Q_+})$; 证明 $u=0$ 在 $Q_0$ 上.

设 $(u_n)$ 为 $C_c^1(Q_+)$ 中的序列使得 $u_n\to u$ 在 $W^{1,p}(Q_+)$ 中.
对 $(x',x_N)\in Q_+$, 我们有
\begin{equation*}
\abs{u_n(x',x_N)}\le\int_0^{x_N}\left|\frac{\partial u_n}{\partial x_N}
(x',t)\right|dt,
\end{equation*}
于是对 $0<\varepsilon<1$,
\begin{equation*}
\frac1\varepsilon\int_{\abs{x'}<1}\int_0^\varepsilon
\abs{u_n(x',x_N)}\,dx'\,dx_N
\le\int_{\abs{x'}<1}\int_0^\varepsilon
\left|\frac{\partial u_n}{\partial x_N}(x',t)\right|dx'\,dt.
\end{equation*}
取极限, 当 $n\to\infty$ 时 ($\varepsilon>0$ 固定) 我们得到
\begin{equation*}
\frac1\varepsilon\int_{\abs{x'}<1}\int_0^\varepsilon
\abs{u(x',x_N)}\,dx'\,dx_N
\le\int_{\abs{x'}<1}\int_0^\varepsilon
\left|\frac{\partial u}{\partial x_N}(x',t)\right|dx'\,dt.
\end{equation*}
最后, 令 $\varepsilon\to0$, 我们得到
\begin{equation*}
\int_{\abs{x'}<1}\abs{u(x',0)}\,dx'=0
\end{equation*}
(因为 $u\in C(\overline{Q_+})$ 且 $\frac{\partial u}{\partial x_N}
\in L^1(Q_+)$).于是 $u=0$ 在 $Q_0$ 上.
\end{proof}

\begin{remark}
\label{rem:9.19}
在 (i)$\Rightarrow$(ii) 的证明中我们没有用到 $\Omega$ 的光滑性.然而,
逆命题 (ii)$\Rightarrow$(i) 需要对 $\Omega$ 加光滑性假设 (例如考虑
$\Omega=\R^N\setminus\{0\}$, $N\ge2$, $p\le N$).
\end{remark}

下面是 $W_0^{1,p}$ 的另一个刻画.

\begin{proposition}
\label{prop:9.18}
设 $\Omega$ 属于 $C^1$ 类.设 $u\in L^p(\Omega)$, $1<p<\infty$.下述性质
等价:

\noindent (i) $u\in W_0^{1,p}(\Omega)$;

\noindent (ii) 存在常数 $C$ 使得
\begin{equation*}
\left|\int_\Omega u\frac{\partial\varphi}{\partial x_i}\right|
\le C\norm{\varphi}_{L^{p'}(\Omega)}
\quad\forall\varphi\in C_c^1(\R^N),\ \forall i=1,2,\dots,N;
\end{equation*}

\noindent (iii) 函数
\begin{equation*}
\bar u(x)=
\begin{cases}
u(x) & \text{若}\ x\in\Omega,\\
0 & \text{若}\ x\in\R^N\setminus\Omega
\end{cases}
\end{equation*}
属于 $W^{1,p}(\R^N)$, 并且此时
$\frac{\partial\bar u}{\partial x_i}=\frac{\partial u}{\partial x_i}$.
\end{proposition}

\begin{proof}
\noindent (i)$\Rightarrow$(ii). 设 $(u_n)$ 为 $C_c^1(\Omega)$ 中的序列
使得 $u_n\to u$ 在 $W^{1,p}$ 中.对 $\varphi\in C_c^1(\R^N)$ 我们有
\begin{equation*}
\left|\int_\Omega u_n\frac{\partial\varphi}{\partial x_i}\right|
=\left|\int_\Omega\frac{\partial u_n}{\partial x_i}\varphi\right|
\le\left\|\frac{\partial u_n}{\partial x_i}\right\|_p
\norm{\varphi}_{p'}.
\end{equation*}
取极限即得 (ii).

\noindent (ii)$\Rightarrow$(iii). 设 $\varphi\in C_c^1(\R^N)$; 我们有
\begin{equation*}
\left|\int_{\R^N}\bar u\frac{\partial\varphi}{\partial x_i}\right|
=\left|\int_\Omega u\frac{\partial\varphi}{\partial x_i}\right|
\le C\norm{\varphi}_{L^{p'}(\Omega)}
\le C\norm{\varphi}_{L^{p'}(\R^N)}.
\end{equation*}
于是 $\bar u\in W^{1,p}(\R^N)$ (由命题  \ref{prop:9.3}).

\noindent (iii)$\Rightarrow$(i). 总可以假设 $\Omega$ 有界 (否则考虑序列
$(\zeta_nu)$).利用局部坐标卡与单位分解, 这归结为下述问题.设
$u\in L^p(Q_+)$ 使得函数
\begin{equation*}
\bar u(x)=
\begin{cases}
u(x) & \text{若}\ x\in Q,\ x_N>0,\\
0 & \text{若}\ x\in Q,\ x_N<0
\end{cases}
\end{equation*}
属于 $W^{1,p}(Q)$; 证明
\begin{equation*}
\alpha u\in W_0^{1,p}(Q_+)\quad\forall\alpha\in C_c^1(Q).
\end{equation*}
设 $(\rho_n)$ 为一列磨光核使得
\begin{equation*}
\supp\rho_n\subset\left\{x\in\R^N;\ \frac1{2n}<x_N<\frac1n\right\};
\end{equation*}
例如可以取 $\rho_n(x)=n^N\rho(nx)$, 其中
$\supp\rho\subset\{x\in\R^N;\ (1/2)<x_N<1\}$.于是
$\rho_n\star(\alpha\bar u)\to\alpha\bar u$ 在 $W^{1,p}(\R^N)$ 中 (注意
$\alpha\bar u$ 在 $Q$ 外延拓为 $0$ 后属于 $W^{1,p}(\R^N)$).另一方面,
\begin{equation*}
\supp(\rho_n\star\alpha\bar u)\subset\supp\rho_n+
\supp(\alpha\bar u)\subset Q_+
\end{equation*}
对足够大的 $n$ 成立.由此推出
\begin{equation*}
\rho_n\star(\alpha\bar u)\in C_c^1(Q_+),
\end{equation*}
于是 $\alpha u\in W_0^{1,p}(Q_+)$.
\end{proof}

\begin{remark}
\label{rem:9.20}
推论  \ref{cor:9.14} 的证明使用了延拓算子, 因此必须假设 $\Omega$ 光滑.
若把 $W^{1,p}(\Omega)$ 换成 $W_0^{1,p}(\Omega)$, 则可以使用在 $\Omega$ 外
零延拓的典范延拓, 它对任意区域 $\Omega$ 都有效 (在命题  \ref{prop:9.18}
的证明中, 蕴涵 (i)$\Rightarrow$(iii) 没有用到 $\Omega$ 的任何光滑性假设).
特别地, 推论  \ref{cor:9.14} 的结论对任意开集 $\Omega$ 的
$W_0^{1,p}(\Omega)$ 成立.类似地, 定理  \ref{thm:9.16} 的结论对任意有界
开集 $\Omega$ 的 $W_0^{1,p}(\Omega)$ 成立.由定理  \ref{thm:9.9} 还可以推出:
若 $\Omega$ 是任意开集且 $1\le p<N$, 则
\begin{equation}
\norm{u}_{L^{p^\star}(\Omega)}\le C(p,N)\norm{\nabla u}_{L^p(\Omega)}
\quad\forall u\in W_0^{1,p}(\Omega).
\label{eq:9.30}
\end{equation}
\end{remark}

\begin{keycorollary}[Poincaré 不等式]
\label{cor:9.19}
设 $1\le p<\infty$, $\Omega$ 为有界开集.则存在常数 $C$ (依赖于 $\Omega$
与 $p$) 使得
\begin{equation*}
\norm{u}_{L^p(\Omega)}\le C\norm{\nabla u}_{L^p(\Omega)}
\quad\forall u\in W_0^{1,p}(\Omega).
\end{equation*}
特别地, 表达式 $\norm{\nabla u}_{L^p(\Omega)}$ 是 $W_0^{1,p}(\Omega)$ 上的
范数, 且与范数 $\norm{u}_{W^{1,p}}$ 等价; 在 $H_0^1(\Omega)$ 上, 表达式
$\sum_{i=1}^N\int_\Omega\frac{\partial u}{\partial x_i}
\frac{\partial v}{\partial x_i}$ 是标量积, 它诱导的范数为
$\norm{\nabla u}_{L^2}$, 且与范数 $\norm{u}_{H^1}$ 等价.
\end{keycorollary}

\begin{remark}
\label{rem:9.21}
若 $\Omega$ 有有限测度, 或者 $\Omega$ 在某个轴上的投影有界, Poincaré
不等式仍然成立.
\end{remark}

\begin{remark}
\label{rem:9.22}
对每个整数 $m\ge1$ 与 $1\le p<\infty$, 定义 $W_0^{m,p}(\Omega)$ 为
$C_c^m(\Omega)$ 在 $W^{m,p}(\Omega)$ 中的闭包.``大致''说来, 函数
$u$ 属于 $W_0^{m,p}(\Omega)$, 如果 $u\in W^{m,p}(\Omega)$ 且对一切满足
$\abs{\alpha}\le m-1$ 的多重指标 $\alpha$, $D^\alpha u=0$ 在 $\Gamma$ 上.
重要的是要注意到, 当 $m\ge2$ 时 $W_0^{m,p}(\Omega)$ 与
$W^{m,p}(\Omega)\cap W_0^{1,p}(\Omega)$ 是有区别的.
\end{remark}

\subsection*{$W_0^{1,p}(\Omega)$ 的对偶空间}

\begin{notation}
我们用 $W^{-1,p'}(\Omega)$ 表示 $W_0^{1,p}(\Omega)$ 的对偶空间
($1\le p<\infty$), 用 $H^{-1}(\Omega)$ 表示 $H_0^1(\Omega)$ 的对偶.
$L^2(\Omega)$ 的对偶与 $L^2(\Omega)$ 等同, 但我们不把 $H_0^1(\Omega)$ 与
它的对偶等同 (见第~
\ref{chap:5}~章注记 3).我们有包含关系
\begin{equation*}
H_0^1(\Omega)\subset L^2(\Omega)\subset H^{-1}(\Omega),
\end{equation*}
其中这些嵌入连续且稠密.

若 $\Omega$ 有界, 则
\begin{equation*}
W_0^{1,p}(\Omega)\subset L^2(\Omega)\subset W^{-1,p'}(\Omega)
\quad\text{若}\ 2N/(N+2)\le p<\infty,
\end{equation*}
且嵌入连续且稠密.若 $\Omega$ 无界, 同样的结论成立, 但只对范围
$2N/(N+2)\le p\le2$.
\end{notation}

$W^{-1,p'}$ 中的元素由下述结果完全刻画:

\begin{proposition}
\label{prop:9.20}
设 $F\in W^{-1,p'}(\Omega)$.则存在函数
$f_0,f_1,f_2,\dots,f_N\in L^{p'}(\Omega)$ 使得
\begin{equation*}
\dual{F}{v}=\int_\Omega f_0v+\sum_{i=1}^N\int_\Omega
f_i\frac{\partial v}{\partial x_i}\quad\forall v\in W_0^{1,p}(\Omega),
\end{equation*}
且
\begin{equation*}
\norm{F}=\max_{0\le i\le N}\norm{f_i}_{p'}.
\end{equation*}
若 $\Omega$ 有界, 我们可以取 $f_0=0$.
\end{proposition}

\begin{proof}
仿照命题  \ref{prop:8.14} 的证明.
\end{proof}

\section{\index{tuoyuanbianzhiwenti@椭圆边值问题 (elliptic boundary value problem)}
一些边值问题的变分表述}
\label{sec:9.5}

现在我们将利用前面的框架来研究一些二阶椭圆型偏微分方程 (= PDE).

\begin{example}[Laplace 算子的齐次 Dirichlet 问题]
\label{ex:9.1}
设 $\Omega\subset\R^N$ 为有界开集.我们寻找函数 $u:\Omega\to\R$ 满足
\begin{equation}
\begin{cases}
-\Delta u+u=f & \text{在}\ \Omega\ \text{中},\\
u=0 & \text{在}\ \Gamma=\partial\Omega\ \text{上},
\end{cases}
\label{eq:9.31}
\end{equation}
其中
\begin{equation*}
\Delta u=\sum_{i=1}^N\frac{\partial^2u}{\partial x_i^2}
=\text{$u$ 的 Laplace 算子},
\end{equation*}
而 $f$ 是 $\Omega$ 上给定的函数.边界条件 $u=0$ 在 $\Gamma$ 上称为
(\emph{齐次}) \emph{Dirichlet 条件}.
\end{example}

\begin{definition}
\eqref{eq:9.31} 的\emph{古典解}是满足 \eqref{eq:9.31} (在通常意义下) 的
函数 $u\in C^2(\Omega)$.\eqref{eq:9.31} 的\emph{弱解}是满足
\begin{equation}
\int_\Omega\nabla u\cdot\nabla v+\int_\Omega uv=\int_\Omega fv
\quad\forall v\in H_0^1(\Omega)
\label{eq:9.32}
\end{equation}
的函数 $u\in H_0^1(\Omega)$, 其中
$\nabla u\cdot\nabla v=\sum_{i=1}^N\frac{\partial u}{\partial x_i}
\frac{\partial v}{\partial x_i}$.
\end{definition}

我们实施第~
\ref{chap:8}~章中描述的程序.

\runin{步骤 A: 每个古典解都是弱解.}

事实上, $u\in H^1(\Omega)\cap C(\overline\Omega)$ 且 $u=0$ 在 $\Gamma$ 上,
于是由定理  \ref{thm:9.17} (另见注记 \ref{rem:9.19})
$u\in H_0^1(\Omega)$.另一方面, 若 $v\in C_c^1(\Omega)$ 我们有
\begin{equation*}
\int_\Omega\nabla u\cdot\nabla v+\int_\Omega uv=\int_\Omega fv,
\end{equation*}
由稠密性, 它对一切 $v\in H_0^1(\Omega)$ 仍然成立.

\runin{步骤 B: 弱解的存在性与唯一性.}

这就是下述基本结果的内容.

\begin{keytheorem}[Dirichlet, Riemann, Poincaré, Hilbert]
\label{thm:9.21}
给定任意 $f\in L^2(\Omega)$, \eqref{eq:9.31} 存在唯一的弱解
$u\in H_0^1(\Omega)$.此外, $u$ 由
\begin{equation*}
\min_{v\in H_0^1(\Omega)}\left\{\frac12\int_\Omega
(\abs{\nabla v}^2+\abs{v}^2)-\int_\Omega fv\right\}
\end{equation*}
得到.这就是\emph{Dirichlet 原理}.
\end{keytheorem}

\begin{proof}
在 Hilbert 空间 $H=H_0^1(\Omega)$ 中对双线性形式
\begin{equation*}
a(u,v)=\int_\Omega(\nabla u\cdot\nabla v+uv)
\end{equation*}
与线性泛函 $\varphi:v\mapsto\int_\Omega fv$ 应用 Lax--Milgram 定理.
\end{proof}

\runin{步骤 C: \index{ruojiedezhengzexing@弱解的正则性 (regularity of weak solutions)}
弱解的正则性.}

这个问题是微妙的.我们将在第~\ref{sec:9.6}~节中讨论它.

\runin{步骤 D: 复原古典解.}

假设 \eqref{eq:9.31} 的弱解 $u\in H_0^1(\Omega)$ 属于 $C^2(\Omega)$,
并假设 $\Omega$ 属于 $C^1$ 类.则 $u=0$ 在 $\Gamma$ 上 (由定理
 \ref{thm:9.17}).另一方面, 我们有
\begin{equation*}
\int_\Omega(-\Delta u+u)v=\int_\Omega fv\quad\forall v\in C_c^1(\Omega),
\end{equation*}
于是 $-\Delta u+u=f$ 在 $\Omega$ 上几乎处处成立 (由推论  \ref{cor:4.24}).
事实上, 由于 $u\in C^2(\Omega)$, $-\Delta u+u=f$ 在 $\Omega$ 上处处成立;
于是 $u$ 是古典解.

我们现在描述另外一些例子.在每个情形中, 精确地指明函数空间与适当的弱
表述都是至关重要的.

\begin{example}[非齐次 Dirichlet 条件]
\label{ex:9.2}
设 $\Omega\subset\R^N$ 为有界开集.我们寻找函数 $u:\Omega\to\R$ 满足
\begin{equation}
\begin{cases}
-\Delta u+u=f & \text{在}\ \Omega\ \text{中},\\
u=g & \text{在}\ \Gamma\ \text{上},
\end{cases}
\label{eq:9.33}
\end{equation}
其中 $f$ 在 $\Omega$ 上给定, $g$ 在 $\Gamma$ 上给定.假设存在函数
$\tilde g\in H^1(\Omega)\cap C(\overline\Omega)$ 使得\footnote{例如, 若
$\Omega$ 属于 $C^1$ 类且 $g\in C^1(\Gamma)$, 这个假设就满足.若 $\Omega$
足够正则, 则不必假设 $\tilde g\in C(\overline\Omega)$.应用迹理论 (见本章
末的评注), 只须知道 $\tilde g\in H^1(\Omega)$, 即
$g\in H^{1/2}(\Gamma)$.} $\tilde g=g$ 在 $\Gamma$ 上, 并考虑集合
\begin{equation*}
K=\{v\in H^1(\Omega);\ v-\tilde g\in H_0^1(\Omega)\}.
\end{equation*}
由定理  \ref{thm:9.17} 推出 $K$ 与 $\tilde g$ 的选取无关, 只依赖于 $g$.
$K$ 是 $H^1(\Omega)$ 中的非空闭凸集.
\end{example}

\begin{definition}
\eqref{eq:9.33} 的\emph{古典解}是满足 \eqref{eq:9.33} 的函数
$u\in C^2(\Omega)$.\eqref{eq:9.33} 的\emph{弱解}是满足
\begin{equation}
\int_\Omega(\nabla u\cdot\nabla v+uv)=\int_\Omega fv
\quad\forall v\in H_0^1(\Omega)
\label{eq:9.34}
\end{equation}
的函数 $u\in K$.
\end{definition}

如上所述, 任何古典解都是弱解.

\begin{keyproposition}
\label{prop:9.22}
给定任意 $f\in L^2(\Omega)$, \eqref{eq:9.33} 存在唯一的弱解
$u\in K$.此外, $u$ 由
\begin{equation*}
\min_{v\in K}\left\{\frac12\int_\Omega
(\abs{\nabla v}^2+v^2)-\int_\Omega fv\right\}
\end{equation*}
得到.
\end{keyproposition}

\begin{proof}
我们断言: $u\in K$ 是 \eqref{eq:9.33} 的弱解当且仅当
\begin{equation}
\int_\Omega\nabla u\cdot(\nabla v-\nabla u)+\int_\Omega u(v-u)
\ge\int_\Omega f(v-u)\quad\forall v\in K.
\label{eq:9.35}
\end{equation}
事实上, 若 $u$ 是 \eqref{eq:9.33} 的弱解, 显然 \eqref{eq:9.35} 成立, 甚至
取等号.反之, 若 $u\in K$ 满足 \eqref{eq:9.35}, 在 \eqref{eq:9.35} 中取
$v=u\pm w$, 其中 $w\in H_0^1(\Omega)$, 就得到 \eqref{eq:9.34}.然后我们可以
应用 Stampacchia 定理 (定理  \ref{thm:5.6}) 来完成证明.
\end{proof}

正则性与古典解的复原的研究遵循与例 \ref{ex:9.1} 相同的模式.

\begin{example}[一般二阶椭圆方程]
\label{ex:9.3}
设 $\Omega\subset\R^N$ 为有界开集.给定函数
$a_{ij}(x)\in C^1(\Omega)$, $1\le i,j\le N$, 满足椭圆性条件
\begin{equation}
\sum_{i,j=1}^Na_{ij}(x)\xi_i\xi_j\ge\alpha\abs{\xi}^2,
\quad\forall x\in\Omega,\ \forall\xi\in\R^N,
\label{eq:9.36}
\end{equation}
其中 $\alpha>0$.还给定函数 $a_0\in C(\Omega)$.我们寻找函数
$u:\Omega\to\R$ 满足
\begin{equation}
\begin{cases}
\displaystyle-\sum_{i,j=1}^N\frac{\partial}{\partial x_j}
\left(a_{ij}\frac{\partial u}{\partial x_i}\right)+a_0u=f
& \text{在}\ \Omega\ \text{中},\\[2mm]
u=0 & \text{在}\ \Gamma\ \text{上}.
\end{cases}
\label{eq:9.37}
\end{equation}
\eqref{eq:9.37} 的\emph{古典解}是满足 \eqref{eq:9.37} (在通常意义下) 的
函数 $u\in C^2(\Omega)$.\eqref{eq:9.37} 的\emph{弱解}是满足
\begin{equation}
\int_\Omega\sum_{i,j=1}^Na_{ij}\frac{\partial u}{\partial x_i}
\frac{\partial v}{\partial x_j}+\int_\Omega a_0uv=\int_\Omega fv
\quad\forall v\in H_0^1(\Omega)
\label{eq:9.38}
\end{equation}
的函数 $u\in H_0^1(\Omega)$.

如上所述, 任何古典解都是弱解.若 $a_0(x)\ge0$ 在 $\Omega$ 上, 则对一切
$f\in L^2(\Omega)$ 存在唯一的弱解 $u\in H_0^1$: 只须在空间
$H=H_0^1$ 中对连续双线性形式
\begin{equation*}
a(u,v)=\int_\Omega\sum_{i,j=1}^Na_{ij}\frac{\partial u}{\partial x_i}
\frac{\partial v}{\partial x_j}+\int_\Omega a_0uv
\end{equation*}
应用 Lax--Milgram 定理.$a(\ ,\ )$ 的强制性来自椭圆性假设, 假设
$a_0\ge0$ 以及 Poincaré 不等式.若矩阵 $(a_{ij})$ 还是对称的, 则形式
$a(\ ,\ )$ 对称, 且 $u$ 由
\begin{equation*}
\min_{v\in H_0^1}\left\{\frac12\int_\Omega\left(
\sum_{i,j=1}^Na_{ij}\frac{\partial v}{\partial x_i}
\frac{\partial v}{\partial x_j}+a_0v^2\right)-\int_\Omega fv\right\}
\end{equation*}
得到.

现在我们考虑一个更一般的问题: 寻找函数 $u:\Omega\to\R$ 满足
\begin{equation}
\begin{cases}
\displaystyle-\sum_{i,j}\frac{\partial}{\partial x_j}
\left(a_{ij}\frac{\partial u}{\partial x_i}\right)
+\sum_ia_i\frac{\partial u}{\partial x_i}+a_0u=f
& \text{在}\ \Omega\ \text{中},\\[2mm]
u=0 & \text{在}\ \Gamma\ \text{上},
\end{cases}
\label{eq:9.39}
\end{equation}
其中函数 $(a_{ij})\in L^\infty(\Omega)$ 满足椭圆性条件, 而函数
$(a_i)_{0\le i\le N}$ 给定在 $L^\infty(\Omega)$ 中.\eqref{eq:9.39} 的
\emph{弱解}是使得
\begin{equation}
\int_\Omega\sum_{i,j}a_{ij}\frac{\partial u}{\partial x_i}
\frac{\partial v}{\partial x_j}+\int_\Omega\sum_ia_i
\frac{\partial u}{\partial x_i}v+\int_\Omega a_0uv=\int_\Omega fv
\quad\forall v\in H_0^1
\label{eq:9.40}
\end{equation}
的函数 $u\in H_0^1$.相应的连续双线性形式为
\begin{equation}
a(u,v)=\int_\Omega\sum_{i,j}a_{ij}\frac{\partial u}{\partial x_i}
\frac{\partial v}{\partial x_j}+\int_\Omega\sum_ia_i
\frac{\partial u}{\partial x_i}v+\int_\Omega a_0uv.
\label{eq:9.41}
\end{equation}
一般说来这个形式不是对称的;\footnote{在维数 $N$ 中没有像一维那样已知的
办法把它化为对称情形.}在某些情形下它是强制的: 这时可以用 Lax--Milgram
定理得到弱解的存在性与唯一性.在一般情形——即使没有强制性——我们仍有下述
结果.
\end{example}

\begin{theorem}
\label{thm:9.23}
若 $f=0$, 则 \eqref{eq:9.40} 的解 $u\in H_0^1$ 的集合是有限维向量空间,
设其维数为 $d$.此外, 存在 $d$ 维子空间 $F\subset L^2(\Omega)$ 使得\footnote{
换言之, \eqref{eq:9.40} 有解当且仅当 $f$ 满足 $d$ 个正交性条件.}
\begin{equation*}
[\eqref{eq:9.40}\ \text{有解}]
\iff\left[\int_\Omega fv=0\quad\forall v\in F\right].
\end{equation*}
\end{theorem}

\begin{remark}
\label{rem:9.23}
假设与 \eqref{eq:9.40} 相联系的齐次方程 (即 $f=0$ 的情形) 以 $u=0$ 为
唯一解.则对每个 $f\in L^2$, \eqref{eq:9.40} 存在唯一的解
$u\in H_0^1$.\footnote{注意椭圆问题中解的存在性与唯一性之间的密切关系.
这个值得注意的关系是 Fredholm 二择一定理 (定理  \ref{thm:6.6}) 的推论.}
特别地, 若 $a_0\ge0$ 在 $\Omega$ 上, 则可以用极大值原理型的方法证明
$f=0\Rightarrow u=0$.于是我们推出, 只在假设 $a_0\ge0$ 在 $\Omega$ 上
(而对 $a_i$, $1\le i\le N$ 不作任何假设) 的条件下, 对每个 $f\in L^2$,
\eqref{eq:9.40} 都存在唯一的解 $u\in H_0^1$; 参见例如
D.~Gilbarg--N.~Trudinger \bref{GilbargTrudinger}{1}.
\end{remark}

\begin{proof}[定理 \ref{thm:9.23}]
固定 $\lambda>0$, 使得双线性形式
\begin{equation*}
a(u,v)+\lambda\int_\Omega uv
\end{equation*}
在 $H_0^1$ 上强制.对每个 $f\in L^2$ 存在唯一的 $u\in H_0^1$ 满足
\begin{equation*}
a(u,\varphi)+\lambda\int_\Omega u\varphi=\int_\Omega f\varphi
\quad\forall\varphi\in H_0^1.
\end{equation*}
令 $u=Tf$, 于是 $T:L^2\to L^2$ 是紧线性算子 (由于 $\Omega$ 有界, 嵌入
$H_0^1\subset L^2$ 是紧的; 见定理  \ref{thm:9.16} 与注记 \ref{rem:9.20}).
方程 \eqref{eq:9.40} 等价于
\begin{equation}
u=T(f+\lambda u).
\label{eq:9.42}
\end{equation}
令 $v=f+\lambda u$ 为新未知量, \eqref{eq:9.42} 成为
\begin{equation}
v-\lambda Tv=f.
\label{eq:9.43}
\end{equation}
结论由 Fredholm 二择一定理推出.
\end{proof}

\begin{example}[齐次 Neumann 问题]
\label{ex:9.4}
设 $\Omega\subset\R^N$ 为 $C^1$ 类有界区域.我们寻找函数
$u:\Omega\to\R$ 满足
\begin{equation}
\begin{cases}
-\Delta u+u=f & \text{在}\ \Omega\ \text{中},\\[1mm]
\dfrac{\partial u}{\partial n}=0 & \text{在}\ \Gamma\ \text{上},
\end{cases}
\label{eq:9.44}
\end{equation}
其中 $f$ 在 $\Omega$ 上给定; $\frac{\partial u}{\partial n}$ 表示 $u$ 的
外法向导数, 即 $\frac{\partial u}{\partial n}=\nabla u\cdot n$, 其中
$n$ 是 $\Gamma$ 的单位外法向量.边界条件 $\frac{\partial u}{\partial n}=0$
在 $\Gamma$ 上称为 (\emph{齐次}) \emph{Neumann 条件}.
\end{example}

\begin{definition}
\eqref{eq:9.44} 的\emph{古典解}是满足 \eqref{eq:9.44} 的函数
$u\in C^2(\overline\Omega)$.\eqref{eq:9.44} 的\emph{弱解}是满足
\begin{equation}
\int_\Omega\nabla u\cdot\nabla v+\int_\Omega uv=\int_\Omega fv
\quad\forall v\in H^1(\Omega)
\label{eq:9.45}
\end{equation}
的函数 $u\in H^1(\Omega)$.
\end{definition}

\runin{步骤 A: 每个古典解都是弱解.}

回忆由 Green 公式我们有
\begin{equation}
\int_\Omega(\Delta u)v=\int_\Gamma\frac{\partial u}{\partial n}v\,d\sigma
-\int_\Omega\nabla u\cdot\nabla v
\quad\forall u\in C^2(\overline\Omega),\ \forall v\in C^1(\overline\Omega),
\label{eq:9.46}
\end{equation}
其中 $d\sigma$ 是 $\Gamma$ 上的面测度.若 $u$ 是 \eqref{eq:9.44} 的古典
解, 则 $u\in H^1(\Omega)$, 且我们有
\begin{equation*}
\int_\Omega\nabla u\cdot\nabla v+\int_\Omega uv=\int_\Omega fv
\quad\forall v\in C^1(\overline\Omega).
\end{equation*}
由稠密性 (推论  \ref{cor:9.8}) 我们推出
\begin{equation*}
\int_\Omega\nabla u\cdot\nabla v+\int_\Omega uv=\int_\Omega fv
\quad\forall v\in H^1(\Omega).
\end{equation*}

\runin{步骤 B: 弱解的存在性与唯一性.}

\begin{proposition}
\label{prop:9.24}
对每个 $f\in L^2(\Omega)$, \eqref{eq:9.44} 存在唯一的弱解
$u\in H^1(\Omega)$.此外, $u$ 由
\begin{equation*}
\min_{v\in H^1(\Omega)}\left\{\frac12\int_\Omega
(\abs{\nabla v}^2+v^2)-\int_\Omega fv\right\}
\end{equation*}
得到.
\end{proposition}

\begin{proof}
在 $H=H^1(\Omega)$ 中应用 Lax--Milgram 定理.
\end{proof}

\runin{步骤 C: 弱解的正则性.}

这将在第~\ref{sec:9.6}~节中讨论.

\runin{步骤 D: 复原古典解.}

若 $u\in C^2(\overline\Omega)$ 是 \eqref{eq:9.44} 的弱解, 则由
\eqref{eq:9.46} 我们有
\begin{equation}
\int_\Omega(-\Delta u+u)v+\int_\Gamma\frac{\partial u}{\partial n}
v\,d\sigma=\int_\Omega fv\quad\forall v\in C^1(\overline\Omega).
\label{eq:9.47}
\end{equation}
在 \eqref{eq:9.47} 中先取 $v\in C_c^1(\Omega)$, 推出
\begin{equation*}
-\Delta u+u=f\quad\text{在}\ \Omega\ \text{中}.
\end{equation*}
然后回到 \eqref{eq:9.47}, 取 $v\in C^1(\overline\Omega)$; 得到
\begin{equation*}
\int_\Gamma\frac{\partial u}{\partial n}v\,d\sigma=0
\quad\forall v\in C^1(\overline\Omega),
\end{equation*}
因此 $\frac{\partial u}{\partial n}=0$ 在 $\Gamma$ 上.

\begin{example}[无界区域]
\label{ex:9.5}
当 $\Omega$ 是 $\R^N$ 中的无界开集时, 除了 $\Gamma=\partial\Omega$ 上的
通常边界条件之外, 还要在无穷远处加一个边界条件, 例如
$u(x)\to0$ 当 $\abs{x}\to\infty$.在弱解的水平上, 这``翻译''为条件
$u\in H^1$.当然, 必须首先证明: 若 $u$ 是古典解且
$u(x)\to0$ 当 $\abs{x}\to\infty$, 则 $u$ 必属于 $H^1$ (见第~
\ref{chap:8}~章例
\ref{ex:8.8} 中的讨论).下面是几个典型的例子:

\noindent (a) $\Omega=\R^N$; 给定 $f\in L^2(\R^N)$, 方程
\begin{equation*}
-\Delta u+u=f\quad\text{在}\ \R^N\ \text{中}
\end{equation*}
在下述意义下有唯一的弱解:
\begin{equation*}
u\in H^1(\R^N)\quad\text{且}\quad
\int_{\R^N}\nabla u\cdot\nabla v+\int_{\R^N}uv=\int_{\R^N}fv
\quad\forall v\in H^1(\R^N).
\end{equation*}

\noindent (b) $\Omega=\R^N_+$; 给定 $f\in L^2(\R^N_+)$, 问题
\begin{equation*}
\begin{cases}
-\Delta u+u=f & \text{在}\ \R^N_+\ \text{中},\\
u(x',0)=0 & \text{对}\ x'\in\R^{N-1}
\end{cases}
\end{equation*}
在下述意义下有唯一的弱解:
\begin{equation*}
u\in H_0^1(\Omega)\quad\text{且}\quad
\int_\Omega\nabla u\cdot\nabla v+\int_\Omega uv=\int_\Omega fv
\quad\forall v\in H_0^1(\Omega).
\end{equation*}

\noindent (c) $\Omega=\R^N_+$; 给定 $f\in L^2(\R^N_+)$, 问题
\begin{equation*}
\begin{cases}
-\Delta u+u=f & \text{在}\ \R^N_+\ \text{中},\\
\dfrac{\partial u}{\partial x_N}(x',0)=0 & \text{对}\ x'\in\R^{N-1}
\end{cases}
\end{equation*}
在下述意义下有唯一的弱解:
\begin{equation*}
u\in H^1(\Omega)\quad\text{且}\quad
\int_\Omega\nabla u\cdot\nabla v+\int_\Omega uv=\int_\Omega fv
\quad\forall v\in H^1(\Omega).
\end{equation*}
\end{example}


\section{弱解的正则性}
\label{sec:9.6}

\begin{definition}
设 $m\ge1$ 为整数.我们说开集 $\Omega$ 属于 $C^m$ 类, 如果对每个
$x\in\Gamma$ 存在 $x$ 在 $\R^N$ 中的邻域 $U$ 与双射 $H:Q\to U$ 使得
\begin{equation*}
H\in C^m(Q),\quad H^{-1}\in C^m(U),\quad
H(Q_+)=U\cap\Omega,\quad H(Q_0)=U\cap\Gamma.
\end{equation*}
若 $\Omega$ 对一切 $m$ 都属于 $C^m$ 类, 则称 $\Omega$ 属于 $C^\infty$ 类.
\end{definition}

主要的正则性结果如下.

\begin{keytheorem}[Dirichlet 问题的正则性]
\label{thm:9.25}
设 $\Omega$ 是 $C^2$ 类开集, $\Gamma$ 有界 (或者 $\Omega=\R^N_+$).设
$f\in L^2(\Omega)$, $u\in H_0^1(\Omega)$ 满足
\begin{equation}
\int_\Omega\nabla u\nabla\varphi+\int_\Omega u\varphi
=\int_\Omega f\varphi\quad\forall\varphi\in H_0^1(\Omega).
\label{eq:9.48}
\end{equation}
则 $u\in H^2(\Omega)$ 且 $\norm{u}_{H^2}\le C\norm{f}_{L^2}$, 其中 $C$ 是
只依赖于 $\Omega$ 的常数.此外, 若 $\Omega$ 属于 $C^{m+2}$ 类且
$f\in H^m(\Omega)$, 则
\begin{equation*}
u\in H^{m+2}(\Omega)\quad\text{且}\quad
\norm{u}_{H^{m+2}}\le C\norm{f}_{H^m}.
\end{equation*}
特别地, 若 $f\in H^m(\Omega)$ 且 $m>N/2$, 则 $u\in C^2(\Omega)$.最后, 若
$\Omega$ 属于 $C^\infty$ 类且 $f\in C^\infty(\Omega)$, 则
$u\in C^\infty(\Omega)$.
\end{keytheorem}

\begin{theorem}[Neumann 问题的正则性]
\label{thm:9.26}
在与定理  \ref{thm:9.25} 相同的假设下, 对 Neumann 问题的解, 即对满足
\begin{equation}
\int_\Omega\nabla u\cdot\nabla\varphi+\int_\Omega u\varphi
=\int_\Omega f\varphi\quad\forall\varphi\in H^1(\Omega)
\label{eq:9.49}
\end{equation}
的 $u\in H^1(\Omega)$, 可以得到相同的结论.
\end{theorem}

\begin{remark}
\label{rem:9.24}
对与一般二阶椭圆算子相联系的 Dirichlet (或 Neumann) 问题的解, 也可以
得到相同的结论, 即若 $u\in H_0^1(\Omega)$ 满足
\begin{equation*}
\int_\Omega\sum_{i,j}a_{ij}\frac{\partial u}{\partial x_i}
\frac{\partial\varphi}{\partial x_j}+\int_\Omega\sum_ia_i
\frac{\partial u}{\partial x_i}\varphi+\int_\Omega a_0u\varphi
=\int_\Omega f\varphi\quad\forall\varphi\in H_0^1(\Omega),
\end{equation*}
则\footnote{若 $\Omega$ 无界, 我们还必须假设
$D^\alpha a_{ij}\in L^\infty(\Omega)$ $\forall\alpha$, $\abs{\alpha}\le m+1$
以及 $D^\alpha a_i\in L^\infty(\Omega)$ $\forall\alpha$, $\abs{\alpha}\le m$.}
\begin{equation*}
[f\in L^2(\Omega),\ a_{ij}\in C^1(\Omega)\ \text{且}\ a_i\in C(\Omega)]
\Rightarrow u\in H^2(\Omega),
\end{equation*}
并且对 $m\ge1$,
\begin{equation*}
[f\in H^m(\Omega),\ a_{ij}\in C^{m+1}(\Omega)\ \text{且}\ a_i\in C^m(\Omega)]
\Rightarrow u\in H^{m+2}(\Omega).
\end{equation*}
\end{remark}

我们只证明定理  \ref{thm:9.25}; 定理  \ref{thm:9.26} 的证明完全类似.
证明的主要思想如下.我们先考虑 $\Omega=\R^N$ 的情形, 然后考虑
$\Omega=\R^N_+$ 的情形.对一般的区域 $\Omega$ 我们分两步进行:

\noindent 1. \emph{内部正则性}, 即 $u$ 在每个区域 $\omega\subset\subset\Omega$ 上
正则.这里, 证明遵循与 $\Omega=\R^N$ 相同的模式.

\noindent 2. \emph{边界正则性}, 即 $u$ 在边界的某个邻域上正则.这里, 在局部
坐标卡中, 证明类似于 $\Omega=\R^N_+$ 的情形.

我们建议读者在着手处理一般情形之前先好好研究 $\Omega=\R^N$ 与
$\Omega=\R^N_+$ 这两种情形.本节的安排如下:

\noindent A. $\Omega=\R^N$ 的情形.

\noindent B. $\Omega=\R^N_+$ 的情形.

\noindent C. 一般情形:

\noindent C1. 内部估计.

\noindent C2. 边界附近的估计.

证明的本质要素是 L.~Nirenberg 的\emph{平移方法}\footnote{也称为\emph{差商}
技巧.}.

\subsection*{A. $\Omega=\R^N$ 的情形}

\noindent\emph{记号.} 给定 $h\in\R^N$, $h\neq0$, 令
\begin{equation*}
D_hu=\frac1{\abs{h}}(\tau_hu-u),
\quad\text{即}\quad
D_hu(x)=\frac{u(x+h)-u(x)}{\abs{h}}.
\end{equation*}
在 \eqref{eq:9.48} 中取 $\varphi=D_{-h}(D_hu)$.这是允许的, 因为
$\varphi\in H^1(\R^N)$ (由于 $u\in H^1(\R^N)$); 我们得到
\begin{equation*}
\int\abs{\nabla D_hu}^2+\int\abs{D_hu}^2=\int fD_{-h}(D_hu),
\end{equation*}
从而
\begin{equation}
\norm{D_hu}_{H^1}^2\le\norm{f}_2\norm{D_{-h}(D_hu)}_2.
\label{eq:9.50}
\end{equation}
另一方面, 回忆 (命题  \ref{prop:9.3})
\begin{equation}
\norm{D_{-h}v}_2\le\norm{\nabla v}_2\quad\forall v\in H^1.
\label{eq:9.51}
\end{equation}
把它用于 $v=D_hu$, 得到
\begin{equation*}
\norm{D_hu}_{H^1}^2\le\norm{f}_2\norm{\nabla(D_hu)}_2,
\end{equation*}
从而
\begin{equation*}
\norm{D_hu}_{H^1}\le\norm{f}_2.
\end{equation*}
特别地,
\begin{equation}
\left\|D_h\frac{\partial u}{\partial x_i}\right\|_2\le\norm{f}_2
\quad\forall i=1,2,\dots,N.
\label{eq:9.52}
\end{equation}
再次应用命题  \ref{prop:9.3}, 我们看到
$\frac{\partial u}{\partial x_i}\in H^1$, 从而 $u\in H^2$.

我们现在证明 $f\in H^1\Rightarrow u\in H^3$.记 $Du$ 为导数
$\frac{\partial u}{\partial x_i}$ ($1\le i\le N$) 中的任何一个.我们已经知道
$Du\in H^1$.还须证明 $Du\in H^2$.为此只须验证
\begin{equation}
\int\nabla(Du)\cdot\nabla\varphi+\int(Du)\varphi=\int(Df)\varphi
\quad\forall\varphi\in H^1
\label{eq:9.53}
\end{equation}
(然后我们就可以对 $Du$ 应用前面的分析, 得到 $Du\in H^2$).若
$\varphi\in C_c^\infty(\R^N)$, 我们可以在 \eqref{eq:9.48} 中把 $\varphi$ 换成
$D\varphi$; 它成为
\begin{equation*}
\int\nabla u\cdot\nabla(D\varphi)+\int uD\varphi=\int fD\varphi,
\end{equation*}
从而
\begin{equation*}
\int\nabla(Du)\cdot\nabla\varphi+\int(Du)\varphi=\int(Df)\varphi
\quad\forall\varphi\in C_c^\infty(\R^N).
\end{equation*}
这蕴含 \eqref{eq:9.53}, 因为 $C_c^\infty(\R^N)$ 在 $H^1(\R^N)$ 中稠密
(命题 9.2).
% 原书此处印作 "(Proposition 9.2)" (印刷页 301); 本章的命题编号为 9.1, 9.3,
% 9.4, ... 并无编号 9.2 的命题 (9.2 是定理 9.2, Friedrichs 的逼近定理, 其
% Ω=R^N 的情形恰好给出这一稠密性). 依技能规范保留原样并记录, 不擅自更正,
% 故此处不用 \ref 以免产生未定义引用.

为证明 $f\in H^m\Rightarrow u\in H^{m+2}$, 只须对 $m$ 作归纳并应用
\eqref{eq:9.53}.

\subsection*{B. $\Omega=\R^N_+$ 的情形}

我们再次使用平移, 但只在切向进行, 即沿方向
$h\in\R^{N-1}\times\{0\}$: 我们说 $h$ \emph{平行于边界}, 并记作
$h\parallel\Gamma$.关键的一点是
\begin{equation*}
u\in H_0^1(\Omega)\ \Rightarrow\ \tau_hu\in H_0^1(\Omega)
\quad\text{若}\ h\parallel\Gamma.
\end{equation*}
换言之, $H_0^1(\Omega)$ 在切向平移下不变.

我们取 $h\parallel\Gamma$, 并在 \eqref{eq:9.48} 中代入
$\varphi=D_{-h}(D_hu)$; 得到
\begin{equation*}
\int\abs{\nabla(D_hu)}^2+\int\abs{D_hu}^2=\int fD_{-h}(D_hu),
\end{equation*}
即
\begin{equation}
\norm{D_hu}_{H^1}^2\le\norm{f}_2\norm{D_{-h}(D_hu)}_2.
\label{eq:9.54}
\end{equation}
现在我们使用下述引理.

\begin{lemma}
\label{lem:9.6}
我们有
\begin{equation*}
\norm{D_hv}_{L^2(\Omega)}\le\norm{\nabla v}_{L^2(\Omega)}
\quad\forall v\in H^1(\Omega),\ \forall h\parallel\Gamma.
\end{equation*}
\end{lemma}

\begin{proof}
先取 $v\in C_c^1(\R^N)$, 仿照命题  \ref{prop:9.3} 的证明 (注意
$\Omega+th=\Omega$ 对一切 $t$ 与一切 $h\parallel\Gamma$ 成立).对一般的
$v\in H^1(\Omega)$ 用稠密性论证.
\end{proof}

结合 \eqref{eq:9.54} 与引理  \ref{lem:9.6}, 我们得到
\begin{equation}
\norm{D_hu}_{H^1}\le\norm{f}_2\quad\forall h\parallel\Gamma.
\label{eq:9.55}
\end{equation}
设 $1\le j\le N$, $1\le k\le N-1$, $h=\abs{h}e_k$, 并设
$\varphi\in C_c^\infty(\Omega)$.我们有
\begin{equation*}
\int D_h\left(\frac{\partial u}{\partial x_j}\right)\varphi
=-\int uD_{-h}\left(\frac{\partial\varphi}{\partial x_j}\right),
\end{equation*}
而由 \eqref{eq:9.55},
\begin{equation*}
\left|\int uD_{-h}\left(\frac{\partial\varphi}{\partial x_j}\right)\right|
\le\norm{f}_2\norm{\varphi}_2.
\end{equation*}
令 $h\to0$ 取极限, 它成为
\begin{equation}
\left|\int u\frac{\partial^2\varphi}{\partial x_j\partial x_k}\right|
\le\norm{f}_2\norm{\varphi}_2
\quad\forall1\le j\le N,\ \forall1\le k\le N-1.
\label{eq:9.56}
\end{equation}
最后, 我们断言
\begin{equation}
\left|\int u\frac{\partial^2\varphi}{\partial x_N^2}\right|
\le\norm{f}_2\norm{\varphi}_2\quad\forall\varphi\in C_c^\infty(\Omega).
\label{eq:9.57}
\end{equation}
为证明 \eqref{eq:9.57}, 我们回到方程 \eqref{eq:9.48}, 推出
\begin{equation*}
\left|\int u\frac{\partial^2\varphi}{\partial x_N^2}\right|
\le\sum_{i=1}^{N-1}\left|\int u\frac{\partial^2\varphi}{\partial x_i^2}\right|
+\left|\int(f-u)\varphi\right|\le C\norm{f}_2\norm{\varphi}_2
\end{equation*}
(这里用了 \eqref{eq:9.56}).结合 \eqref{eq:9.56} 与 \eqref{eq:9.57}, 我们得到
\begin{equation*}
\left|\int u\frac{\partial^2\varphi}{\partial x_j\partial x_k}\right|
\le C\norm{f}_2\norm{\varphi}_2
\quad\forall\varphi\in C_c^\infty(\Omega),\ \forall1\le j,k\le N.
\end{equation*}
于是 $u\in H^2(\Omega)$, 因为存在函数 $f_{jk}\in L^2(\Omega)$ 使得
\begin{equation*}
\int u\frac{\partial^2\varphi}{\partial x_j\partial x_k}
=\int f_{jk}\varphi\quad\forall\varphi\in C_c^\infty(\Omega)
\end{equation*}
(如同命题  \ref{prop:8.3} 的证明).

最后我们证明 $f\in H^m(\Omega)\Rightarrow u\in H^{m+2}(\Omega)$.用 $Du$
表示任一切向导数 $Du=\frac{\partial u}{\partial x_j}$, $1\le j\le N-1$.
我们先建立下述结果.

\begin{lemma}
\label{lem:9.7}
设 $u\in H^2(\Omega)\cap H_0^1(\Omega)$ 满足 \eqref{eq:9.48}.则
$Du\in H_0^1(\Omega)$, 并且
\begin{equation}
\int\nabla(Du)\cdot\nabla\varphi+\int(Du)\varphi=\int(Df)\varphi
\quad\forall\varphi\in H_0^1(\Omega).
\label{eq:9.58}
\end{equation}
\end{lemma}

\begin{proof}
唯一微妙之处在于证明 $Du\in H_0^1(\Omega)$; 至于 \eqref{eq:9.58}, 它可由
\eqref{eq:9.48} 取 $D\varphi$ 代替 $\varphi$ (其中
$\varphi\in C_c^\infty(\Omega)$) 然后利用稠密性得到.设 $h=\abs{h}e_j$,
$1\le j\le N-1$, 于是 $D_hu\in H_0^1$ (因为 $H_0^1$ 在切向平移下不变).
由引理  \ref{lem:9.6} 我们有
\begin{equation*}
\norm{D_hu}_{H^1}\le\norm{u}_{H^2}.
\end{equation*}
于是存在序列 $h_n\to0$ 使得 $D_{h_n}u$ 弱收敛于 $H_0^1$ 中某个 $g$
(因为 $H_0^1$ 是 Hilbert 空间).特别地, $D_{h_n}u\rightharpoonup g$ 在
$L^2$ 中弱收敛.对 $\varphi\in C_c^\infty(\Omega)$ 我们有
\begin{equation*}
\int(D_hu)\varphi=\int u(D_{-h}\varphi),
\end{equation*}
令 $h_n\to0$ 取极限, 得到
\begin{equation*}
\int g\varphi=-\int u\frac{\partial\varphi}{\partial x_j}
\quad\forall\varphi\in C_c^\infty(\Omega).
\end{equation*}
因此 $\frac{\partial u}{\partial x_j}=g\in H_0^1(\Omega)$.
\end{proof}

\runin{证明 $f\in H^m\Rightarrow u\in H^{m+2}$.} 对 $m$ 作归纳.假设结论
到 $m$ 阶成立, 并设 $f\in H^{m+1}$.我们已经知道 $u\in H^{m+2}$; 又 $Du$
(任一切向导数) 属于 $H_0^1(\Omega)$ 且满足 \eqref{eq:9.58}.把归纳假设用于
$Du$ 与 $Df$, 我们看到 $Du\in H^{m+2}$.为完成证明, 例如只须验证
$\frac{\partial^2u}{\partial x_N^2}\in H^{m+1}$.为此我们再一次回到方程
\eqref{eq:9.48}, 把它写成
\begin{equation*}
\frac{\partial^2u}{\partial x_N^2}
=-\sum_{i=1}^{N-1}\frac{\partial^2u}{\partial x_i^2}+u-f\in H^{m+1}.
\end{equation*}

\subsection*{C. 一般情形}

我们只证明 $f\in L^2(\Omega)\Rightarrow u\in H^2(\Omega)$; 蕴含关系
$f\in H^m\Rightarrow u\in H^{m+2}$ 如同情形 A 与 B 那样对 $m$ 作归纳.
为简化叙述, 我们假设 $\Omega$ 有界.我们使用单位分解, 并像定理  \ref{thm:9.7}
的证明那样写 $u=\sum_{i=0}^k\theta_iu$.

\runin{C1. 内部估计.} 我们断言 $\theta_0u\in H^2(\Omega)$.由于
$\theta_0|_\Omega\in C_c^\infty(\Omega)$, 把 $\theta_0u$ 在 $\Omega$ 外延拓为
$0$ 所得的函数属于 $H^1(\R^N)$ (见注记 \ref{rem:9.4}(ii)).容易验证
$\theta_0u$ 是 $\R^N$ 中下述方程的弱解:
\begin{equation*}
-\Delta(\theta_0u)+\theta_0u=\theta_0f-2\nabla\theta_0\cdot\nabla u
-(\Delta\theta_0)u\defeq g,
\end{equation*}
其中 $g\in L^2(\R^N)$.由情形 A 推出 $\theta_0u\in H^2(\R^N)$ 且
\begin{equation*}
\norm{\theta_0u}_{H^2}\le C(\norm{f}_2+\norm{u}_{H^1})
\le C'\norm{f}_2
\end{equation*}
(因为由 \eqref{eq:9.48} 有 $\norm{u}_{H^1}\le\norm{f}_2$).

\runin{C2. 边界附近的估计.} 我们断言 $\theta_iu\in H^2(\Omega)$,
$1\le i\le k$.回忆我们有双射 $H:Q\to U_i$ 使得
\begin{equation*}
H\in C^2(Q),\quad J=H^{-1}\in C^2(U_i),\quad
H(Q_+)=\Omega\cap U_i,\quad H(Q_0)=\Gamma\cap U_i.
\end{equation*}
我们写 $x=H(y)$, $y=H^{-1}(x)=J(x)$.容易验证
$v=\theta_iu\in H_0^1(\Omega\cap U_i)$, 且 $v$ 是 $\Omega\cap U_i$ 中下述
方程的弱解:
\begin{equation*}
-\Delta v=\theta_if-\theta_iu-2\nabla\theta_i\cdot\nabla u
-(\Delta\theta_i)u\defeq g,
\end{equation*}
其中 $g\in L^2(\Omega\cap U_i)$ 且 $\norm{g}_2\le C\norm{f}_2$.更精确地说,
\begin{equation}
\int_{\Omega\cap U_i}\nabla v\cdot\nabla\varphi\,dx
=\int_{\Omega\cap U_i}g\varphi\,dx
\quad\forall\varphi\in H_0^1(\Omega\cap U_i).
\label{eq:9.59}
\end{equation}
现在我们把 $v|_{\Omega\cap U_i}$ 转移到 $Q_+$ 上.令
\begin{equation*}
w(y)=v(H(y))\quad\text{对}\ y\in Q_+,
\quad\text{即}\quad
w(Jx)=v(x)\quad\text{对}\ x\in\Omega\cap U_i.
\end{equation*}
下述基本引理表明方程 \eqref{eq:9.59} 变成 $w$ 在 $Q_+$ 上的一个二阶椭圆
方程.\footnote{更一般地, 如果我们从 $v$ 的椭圆方程出发, 就会得到 $w$ 的椭圆
方程: 椭圆性条件在变量替换下保持不变.}

\begin{lemma}
\label{lem:9.8}
在上述记号下, $w$ 属于 $H_0^1(Q_+)$ 且满足
\begin{equation}
\sum_{k,\ell=1}^N\int_{Q_+}a_{k\ell}\frac{\partial w}{\partial y_k}
\frac{\partial\psi}{\partial y_\ell}\,dy
=\int_{Q_+}\tilde g\psi\,dy\quad\forall\psi\in H_0^1(Q_+),
\label{eq:9.60}
\end{equation}
其中\footnote{$\operatorname{detJac}H$ 表示 Jacobi 行列式, 即 Jacobi 矩阵
$\operatorname{Jac}H=\left(\frac{\partial H_i}{\partial y_j}\right)$ 的行列式.}
$\tilde g=(g\circ H)\abs{\operatorname{detJac}H}\in L^2(Q_+)$, 且函数
$a_{k\ell}\in C^1(Q_+)$ 满足椭圆性条件 \eqref{eq:9.36}.
\end{lemma}

\begin{proof}
设 $\psi\in H_0^1(Q_+)$, 并对 $x\in\Omega\cap U_i$ 令 $\varphi(x)=\psi(Jx)$.
则 $\varphi\in H_0^1(\Omega\cap U_i)$ 且
\begin{equation*}
\frac{\partial v}{\partial x_j}=\sum_k\frac{\partial w}{\partial y_k}
\frac{\partial J_k}{\partial x_j},\qquad
\frac{\partial\varphi}{\partial x_j}=\sum_\ell\frac{\partial\psi}{\partial y_\ell}
\frac{\partial J_\ell}{\partial x_j}.
\end{equation*}
于是由积分中通常的变量替换公式,
\begin{equation*}
\begin{aligned}
\int_{\Omega\cap U_i}\nabla v\cdot\nabla\varphi\,dx
&=\int_{\Omega\cap U_i}\sum_{j,k,\ell}\frac{\partial J_k}{\partial x_j}
\frac{\partial J_\ell}{\partial x_j}\frac{\partial w}{\partial y_k}
\frac{\partial\psi}{\partial y_\ell}\,dx\\
&=\int_{Q_+}\sum_{j,k,\ell}\frac{\partial J_k}{\partial x_j}
\frac{\partial J_\ell}{\partial x_j}\frac{\partial w}{\partial y_k}
\frac{\partial\psi}{\partial y_\ell}\abs{\operatorname{detJac}H}\,dy.
\end{aligned}
\end{equation*}
于是
\begin{equation}
\int_{\Omega\cap U_i}\nabla v\cdot\nabla\varphi\,dx
=\int_{Q_+}\sum_{k,\ell}a_{k\ell}\frac{\partial w}{\partial y_k}
\frac{\partial\psi}{\partial y_\ell}\,dy,
\label{eq:9.61}
\end{equation}
其中
\begin{equation*}
a_{k\ell}=\sum_j\frac{\partial J_k}{\partial x_j}
\frac{\partial J_\ell}{\partial x_j}\abs{\operatorname{detJac}H}.
\end{equation*}
我们注意到 $a_{k\ell}\in C^1(Q_+)$, 且椭圆性条件成立, 因为对一切
$\xi\in\R^N$ 有
\begin{equation*}
\sum_{k,\ell}a_{k\ell}\xi_k\xi_\ell
=\abs{\operatorname{detJac}H}\sum_j
\left|\sum_k\frac{\partial J_k}{\partial x_j}\xi_k\right|^2
\ge\alpha\abs{\xi}^2,
\end{equation*}
其中 $\alpha>0$, 因为 Jacobi 矩阵 $\operatorname{Jac}H$ 与
$\operatorname{Jac}J$ 都非奇异.

另一方面, 我们有
\begin{equation}
\int_{\Omega\cap U_i}g\varphi\,dx
=\int_{Q_+}(g\circ H)\psi\abs{\operatorname{detJac}H}\,dy.
\label{eq:9.62}
\end{equation}
结合 \eqref{eq:9.59}, \eqref{eq:9.61} 与 \eqref{eq:9.62} 即得
\eqref{eq:9.60}.引理  \ref{lem:9.8} 证毕.
\end{proof}

我们现在回到边界估计的证明, 并证明
$w\in H^2(Q_+)$\footnote{下文中我们用 $C$ 表示只依赖于 $a_{k\ell}$ 的各种
常数.} 且 $\norm{w}_{H^2}\le C\norm{\tilde g}_2$.回到
$\Omega\cap U_i$, 这将推出 $\theta_iu$ 属于 $H^2(\Omega\cap U_i)$, 从而
实际上属于 $H^2(\Omega)$ 且 $\norm{\theta_iu}_{H^2}\le C\norm{f}_2$.

如同情形 B ($\Omega=\R^N_+$), 我们使用切向平移.在 \eqref{eq:9.60} 中取
$\psi=D_{-h}(D_hw)$, 其中 $h\parallel Q_0$, 且 $\abs{h}$ 足够小使得
$\psi\in H_0^1(Q_+)$.\footnote{回忆
$\operatorname{supp}w\subset\{(x',x_N);\ \abs{x'}<1-\delta\ \text{且}\
0\le x_N<1-\delta\}$ 对某个 $\delta>0$ 成立.} 我们得到
\begin{equation}
\sum_{k,\ell}\int_{Q_+}D_h\left(a_{k\ell}\frac{\partial w}{\partial y_k}\right)
\frac{\partial}{\partial y_\ell}(D_hw)
=\int_{Q_+}\tilde gD_{-h}(D_hw).
\label{eq:9.63}
\end{equation}
但
\begin{equation}
\int_{Q_+}\tilde gD_{-h}(D_hw)\le\norm{\tilde g}_2
\norm{D_{-h}(D_hw)}_2\le\norm{\tilde g}_2\norm{\nabla D_hw}_2
\label{eq:9.64}
\end{equation}
(由引理  \ref{lem:9.6}).另一方面, 写
\begin{equation*}
D_h\left(a_{k\ell}\frac{\partial w}{\partial y_k}\right)(y)
=a_{k\ell}(y+h)\frac{\partial}{\partial y_k}D_hw(y)
+(D_ha_{k\ell}(y))\frac{\partial w}{\partial y_k}(y),
\end{equation*}
于是
\begin{equation}
\sum_{k,\ell}\int_{Q_+}D_h\left(a_{k\ell}\frac{\partial w}{\partial y_k}\right)
\frac{\partial}{\partial y_\ell}(D_hw)
\ge\alpha\norm{\nabla(D_hw)}_2^2-C\norm{w}_{H^1}\norm{\nabla D_hw}_2.
\label{eq:9.65}
\end{equation}
结合 \eqref{eq:9.64} 与 \eqref{eq:9.65}, 我们得到
\begin{equation}
\norm{\nabla D_hw}_2\le C(\norm{w}_{H^1}+\norm{\tilde g}_2)
\le C\norm{\tilde g}_2
\label{eq:9.66}
\end{equation}
(注意由 \eqref{eq:9.60} 与 Poincaré 不等式有
$\norm{w}_{H^1}\le C\norm{\tilde g}_2$).由 \eqref{eq:9.66} 推出——如同
情形 B 那样——
\begin{equation}
\left|\int_{Q_+}\frac{\partial w}{\partial y_k}
\frac{\partial\psi}{\partial y_\ell}\right|
\le C\norm{\tilde g}_2\norm{\psi}_2
\quad\forall\psi\in C_c^1(Q_+),\ \forall(k,\ell)\neq(N,N).
\label{eq:9.67}
\end{equation}
为推出 $w\in H^2(Q_+)$ (且 $\norm{w}_{H^2}\le C\norm{\tilde g}_2$),
还须证明
\begin{equation}
\left|\int_{Q_+}\frac{\partial w}{\partial y_N}
\frac{\partial\psi}{\partial y_N}\right|
\le C\norm{\tilde g}_2\norm{\psi}_2\quad\forall\psi\in C_c^1(Q_+).
\label{eq:9.68}
\end{equation}
为此我们回到方程, 在其中把 $\psi$ 换成 $(1/a_{NN})\psi$
($\psi\in C_c^1(Q_+)$); 这是允许的, 因为 $a_{NN}\in C^1(Q_+)$ 且
$a_{NN}\ge\alpha>0$.方程成为
\begin{equation*}
\int a_{NN}\frac{\partial w}{\partial y_N}
\frac{\partial}{\partial y_N}\left(\frac{\psi}{a_{NN}}\right)
=\int\frac{\tilde g}{a_{NN}}\psi
-\sum_{(k,\ell)\neq(N,N)}\int a_{k\ell}\frac{\partial w}{\partial y_k}
\frac{\partial}{\partial y_\ell}\left(\frac{\psi}{a_{NN}}\right),
\end{equation*}
即
\begin{equation}
\begin{aligned}
\int\frac{\partial w}{\partial y_N}\frac{\partial\psi}{\partial y_N}
={}&\int\frac1{a_{NN}}\frac{\partial a_{NN}}{\partial y_N}
\frac{\partial w}{\partial y_N}\psi+\int\frac{\tilde g}{a_{NN}}\psi\\
&+\sum_{(k,\ell)\neq(N,N)}\int\frac{\partial w}{\partial y_k}
\frac{\partial a_{k\ell}}{\partial y_\ell}\frac{\psi}{a_{NN}}
-\sum_{(k,\ell)\neq(N,N)}\int\frac{\partial w}{\partial y_k}
\frac{\partial}{\partial y_\ell}\left(\frac{a_{k\ell}\psi}{a_{NN}}\right).
\end{aligned}
\label{eq:9.69}
\end{equation}
结合 \eqref{eq:9.67}\footnote{我们对 $(a_{k\ell}a_{NN})\psi$ 代替 $\psi$
使用 \eqref{eq:9.67}.} 与 \eqref{eq:9.69}, 我们得到
\begin{equation*}
\left|\int_{Q_+}\frac{\partial w}{\partial y_N}
\frac{\partial\psi}{\partial y_N}\right|
\le C(\norm{w}_{H^1}+\norm{\tilde g}_2)\norm{\psi}_2
\quad\forall\psi\in C_c^1(Q_+).
\end{equation*}
这就建立了 \eqref{eq:9.68}, 并完成了边界附近的估计.

\begin{remark}
\label{rem:9.25}
设 $\Omega$ 为任意开集, $u\in H^1(\Omega)$ 使得
\begin{equation*}
\int_\Omega\nabla u\cdot\nabla\varphi=\int_\Omega f\varphi
\quad\forall\varphi\in C_c^\infty(\Omega).
\end{equation*}
我们假设 $f\in H^m(\Omega)$.则对每个 $\theta\in C_c^\infty(\Omega)$ 有
$\theta u\in H^{m+2}(\Omega)$: 我们说 $u\in H^{m+2}_{\mathrm{loc}}(\Omega)$.
为证明这一点, 只须像情形 C1 那样进行并对 $m$ 作归纳.特别地,
$f\in C^\infty(\Omega)\Rightarrow u\in C^\infty(\Omega)$.\footnote{但一般
说来, 我们不能断言例如 $u\in C(\Omega)$ (即使 $\Omega$ 与 $f$ 都非常光滑),
因为没有给定任何边界条件.}

对\emph{很弱}意义下的解也有同样的结论, 即 $u\in L^2(\Omega)$ 使得
\begin{equation*}
-\int_\Omega u\Delta\varphi=\int_\Omega f\varphi
\quad\forall\varphi\in C_c^\infty(\Omega).
\end{equation*}
(证明要稍微微妙一些; 参见例如 S.~Agmon \bref{Agmon}{1}.) 我们要强调椭圆
问题中正则性结果的\emph{局部}性质.更精确地说, 设 $f\in L^2(\Omega)$,
$u\in H_0^1(\Omega)$ 是
\begin{equation*}
\int_\Omega\nabla u\cdot\nabla\varphi+\int_\Omega u\varphi
=\int_\Omega f\varphi\quad\forall\varphi\in H_0^1(\Omega)
\end{equation*}
的唯一弱解.固定 $\omega\subset\subset\Omega$; 则 $u|_\omega$ 依赖于 $f$
在\emph{整个} $\Omega$ 上的值——而不仅是 $f$ 在 $\omega$ 上的值.\footnote{例如,
若 $f\ge0$ 在 $\Omega$ 上, $f=0$ 在 $\omega$ 上, 而 $f>0$ 在 $\Omega$ 的某个
开子集上, 则 $u>0$ 在 $\Omega$ 上 (从而在 $\omega$ 上); 见本章末评注中的强
极大值原理.} 相反, $u|_\omega$ 的\emph{正则性}只依赖于 $f|_\omega$ 的正则性.
例如, 即使 $f$ 在 $\omega$ 外非常不规则, 也有
$f\in C^\infty(\omega)\Rightarrow u\in C^\infty(\omega)$.这个性质称为
\emph{亚椭圆性}.
\end{remark}

\begin{remark}
\label{rem:9.26}
从某种观点看, 正则性结果有点令人惊讶.事实上, 对 $\Delta u$ 即对导数之和
$\sum_k\frac{\partial^2u}{\partial x_k^2}$ 所作的假设, 迫使我们对\emph{每个}
导数 $\frac{\partial^2u}{\partial x_i\partial x_j}$ 分别得到同一性质的结论.
\end{remark}

\section{极大值原理}
\label{sec:9.7}

极大值原理是一个非常有用的工具, 它有许多种表述形式.这里我们介绍几种
简单的形式.

设 $\Omega$ 为 $\R^N$ 的一般开子集.

\begin{keytheorem}[Dirichlet 问题的极大值原理]
\index{Nweijidazhiyuanli@N 维极大值原理 (maximum principle in dimension N)}
\label{thm:9.27}
假设\footnote{若 $\Omega$ 属于 $C^1$ 类, 则可以借助迹理论 (它赋予
$u|_\Gamma$ 以意义; 见本章末的评注) 去掉假设 $u\in C(\overline\Omega)$;
此外若 $u\in H_0^1(\Omega)$, 也可以去掉假设 $u\in C(\overline\Omega)$.}
$f\in L^2(\Omega)$ 且 $u\in H^1(\Omega)\cap C(\overline\Omega)$ 满足
\begin{equation}
\int_\Omega\nabla u\cdot\nabla\varphi+\int_\Omega u\varphi
=\int_\Omega f\varphi\quad\forall\varphi\in H_0^1(\Omega).
\label{eq:9.70}
\end{equation}
则对一切 $x\in\Omega$,
\begin{equation*}
\min\left\{\inf_\Gamma u,\ \inf_\Omega f\right\}\le u(x)
\le\max\left\{\sup_\Gamma u,\ \sup_\Omega f\right\}.
\end{equation*}
(这里以及下文中, $\sup$ 与 $\inf$ 分别表示本性上确界与本性下确界.)
\end{keytheorem}

\begin{proof}
我们使用 Stampacchia 的截断方法.固定函数 $G\in C^1(\R)$ 使得
\begin{equation*}
\begin{array}{ll}
\text{(i)}\ \abs{G'(s)}\le M & \forall s\in\R,\\[1mm]
\text{(ii)}\ G\ \text{在}\ (0,+\infty)\ \text{上严格递增}, &\\[1mm]
\text{(iii)}\ G(s)=0 & \forall s\le0.
\end{array}
\end{equation*}
令
\begin{equation*}
K=\max\left\{\sup_\Gamma u,\ \sup_\Omega f\right\},
\end{equation*}
并假设 $K<\infty$ (否则无须证明).令 $v=G(u-K)$.我们区分两种情形:

\noindent (a) $\abs{\Omega}<\infty$ 的情形.

则 $v\in H^1(\Omega)$ (由命题  \ref{prop:9.5} 用于函数
$t\mapsto G(t-K)-G(-K)$).另一方面, $v\in H_0^1(\Omega)$, 因为
$v\in C(\overline\Omega)$ 且 $v=0$ 在 $\Gamma$ 上 (见定理  \ref{thm:9.17}).
把这个 $v$ 代入 \eqref{eq:9.70}, 并像定理  \ref{thm:8.19} 的证明那样进行.
% 原书此处印作 "as in the proof of Theorem 8.18" (印刷页 308); 第~
% \ref{chap:8}~章并无编号
% 8.18 的定理 (8.18 是命题, 讲二阶边值问题的存在唯一性), 用 Stampacchia 截断
% 方法证明极大值原理的是定理~
% \ref{thm:8.19}. 依技能规范保留原样并记录, 不擅自更正.

\noindent (b) $\abs{\Omega}=\infty$ 的情形.

此时 $K\ge0$ (因为 $f(x)\le K$ 在 $\Omega$ 上几乎处处成立且 $f\in L^2$
蕴含 $K\ge0$).固定 $K'>K$.由命题  \ref{prop:9.5} 用于函数
$t\mapsto G(t-K')$, 我们看到 $v=G(u-K')\in H^1(\Omega)$.此外
$v\in C(\overline\Omega)$ 且 $v=0$ 在 $\Gamma$ 上; 因此 $v\in H_0^1(\Omega)$.
把这个 $v$ 代入 \eqref{eq:9.70}, 我们有
\begin{equation}
\int_\Omega\abs{\nabla u}^2G'(u-K')+\int_\Omega uG(u-K')
=\int_\Omega fG(u-K').
\label{eq:9.71}
\end{equation}
另一方面, $G(u-K')\in L^1(\Omega)$, 因为\footnote{因为
$G(u-K')-G(-K')\le M\abs{u}$ 且 $G(-K')=0$ (由于 $-K'<0$).}
\begin{equation*}
0\le G(u-K')\le M\abs{u},
\end{equation*}
而在集合 $[u\ge K']=\{x\in\Omega;\ u(x)\ge K'\}$ 上有
\begin{equation*}
K'\int_{[u\ge K']}\abs{u}\le\int_\Omega u^2<\infty.
\end{equation*}
由 \eqref{eq:9.71} 我们推出
\begin{equation*}
\int_\Omega(u-K')G(u-K')\le\int_\Omega(f-K')G(u-K')\le0.
\end{equation*}
于是 $u\le K'$ 在 $\Omega$ 上几乎处处成立, 从而 $u\le K$ 在 $\Omega$ 上几乎
处处成立 (因为 $K'>K$ 是任意的).
\end{proof}

\begin{keycorollary}
\label{cor:9.28}
设 $f\in L^2(\Omega)$, $u\in H^1(\Omega)\cap C(\overline\Omega)$\footnote{同
上, 在某些情形下可以去掉假设 $u\in C(\overline\Omega)$.} 满足
\eqref{eq:9.70}.则
\begin{equation}
[u\ge0\ \text{在}\ \Gamma\ \text{上且}\ f\ge0\ \text{在}\ \Omega\ \text{中}]
\Rightarrow[u\ge0\ \text{在}\ \Omega\ \text{中}],
\label{eq:9.72}
\end{equation}
\begin{equation}
\norm{u}_{L^\infty(\Omega)}\le\max\left\{
\norm{u}_{L^\infty(\Gamma)},\ \norm{f}_{L^\infty(\Omega)}\right\}.
\label{eq:9.73}
\end{equation}
特别地,
\begin{equation}
\text{若}\ f=0\ \text{在}\ \Omega\ \text{中, 则}\
\norm{u}_{L^\infty(\Omega)}\le\norm{u}_{L^\infty(\Gamma)},
\label{eq:9.74}
\end{equation}
\begin{equation}
\text{若}\ u=0\ \text{在}\ \Gamma\ \text{上, 则}\
\norm{u}_{L^\infty(\Omega)}\le\norm{f}_{L^\infty(\Omega)}.
\label{eq:9.75}
\end{equation}
\end{keycorollary}

\begin{remark}
\label{rem:9.27}
若 $\Omega$ 有界且 $u$ 是方程
\begin{equation}
-\Delta u+u=f\quad\text{在}\ \Omega\ \text{中}
\label{eq:9.76}
\end{equation}
的古典解, 则可以对定理  \ref{thm:9.27} 给出另一种证明.事实上, 设
$x_0\in\overline\Omega$ 使得 $u(x_0)=\max_{\overline\Omega}u$.

\noindent (i) 若 $x_0\in\Gamma$, 则 $u(x_0)\le\sup_\Gamma u\le K$.

\noindent (ii) 若 $x_0\in\Omega$, 则 $\nabla u(x_0)=0$ 且
$\frac{\partial^2u}{\partial x_i^2}(x_0)\le0$ 对一切 $1\le i\le N$ 成立,
从而 $\Delta u(x_0)\le0$.由此, 利用方程 \eqref{eq:9.76} 我们得到
\begin{equation*}
u(x_0)=f(x_0)+\Delta u(x_0)\le f(x_0)\le K.
\end{equation*}

这个方法的好处是它适用于一般的二阶椭圆方程.例如, 定理  \ref{thm:9.27} 的
结论对
\begin{equation}
-\sum_{i,j=1}^N\frac{\partial}{\partial x_j}
\left(a_{ij}\frac{\partial u}{\partial x_i}\right)
+\sum_{i=1}^Na_i\frac{\partial u}{\partial x_i}+u=f
\quad\text{在}\ \Omega\ \text{中}
\label{eq:9.77}
\end{equation}
也成立.注意若 $x_0\in\Omega$, 则
\begin{equation*}
\sum_{i,j=1}^Na_{ij}(x_0)\frac{\partial^2u}{\partial x_i\partial x_j}(x_0)\le0;
\end{equation*}
事实上, 通过 (依赖于 $x_0$ 的) 坐标变换可以把问题化为矩阵 $a_{ij}(x_0)$
为对角形的情形.定理  \ref{thm:9.27} 的结论对 \eqref{eq:9.77} 的弱解仍然
成立, 但证明要微妙一些; 参见 D.~Gilbarg--N.~Trudinger
\bref{GilbargTrudinger}{1}.
\end{remark}

\begin{proposition}
\label{prop:9.29}
设函数 $a_{ij}\in L^\infty(\Omega)$ 满足椭圆性条件 \eqref{eq:9.36}, 且
$a_i,a_0\in L^\infty(\Omega)$, $a_0\ge0$ 在 $\Omega$ 上.设
$f\in L^2(\Omega)$, $u\in H^1(\Omega)\cap C(\overline\Omega)$\footnote{同
上, 在某些情形下可以去掉假设 $u\in C(\overline\Omega)$.} 使得
\begin{equation}
\int_\Omega\sum_{i,j}a_{ij}\frac{\partial u}{\partial x_i}
\frac{\partial\varphi}{\partial x_j}
+\int_\Omega\sum_ia_i\frac{\partial u}{\partial x_i}\varphi
+\int_\Omega a_0u\varphi=\int_\Omega f\varphi
\quad\forall\varphi\in H_0^1(\Omega).
\label{eq:9.78}
\end{equation}
则
\begin{equation}
[u\ge0\ \text{在}\ \Gamma\ \text{上且}\ f\ge0\ \text{在}\ \Omega\ \text{中}]
\Rightarrow[u\ge0\ \text{在}\ \Omega\ \text{中}].
\label{eq:9.79}
\end{equation}
设另外 $a_0\equiv0$ 且 $\Omega$ 有界.则
\begin{equation}
[f\ge0\ \text{在}\ \Omega\ \text{中}]
\Rightarrow\left[u\ge\inf_\Gamma u\ \text{在}\ \Omega\ \text{中}\right],
\label{eq:9.80}
\end{equation}
且
\begin{equation}
[f=0\ \text{在}\ \Omega\ \text{中}]
\Rightarrow\left[\inf_\Gamma u\le\sup_\Gamma u\ \text{在}\ \Omega\
\text{中}\right].
\label{eq:9.81}
\end{equation}
\end{proposition}

\begin{proof}
我们在 $a_i\equiv0$ ($1\le i\le N$) 的情形下证明这个结果; 一般情形更为
微妙 (参见 D.~Gilbarg--N.~Trudinger \bref{GilbargTrudinger}{1}, 定理 8.1).
证明 \eqref{eq:9.79} 相当于证明
\begin{equation*}
[u\le0\ \text{在}\ \Gamma\ \text{上且}\ f\le0\ \text{在}\ \Omega\ \text{中}]
\Rightarrow[u\le0\ \text{在}\ \Omega\ \text{中}].
\tag*{(79$'$)}
\end{equation*}
在 \eqref{eq:9.78} 中取 $\varphi=G(u)$, 其中 $G$ 如定理  \ref{thm:9.27} 的
证明中那样; 我们于是得到
\begin{equation*}
\int_\Omega\sum_{i,j}a_{ij}\frac{\partial u}{\partial x_i}
\frac{\partial u}{\partial x_j}G'(u)\le0,
\end{equation*}
从而
\begin{equation*}
\int_\Omega\abs{\nabla u}^2G'(u)\le0.
\end{equation*}
令 $H(t)=\int_0^t[G'(s)]^{1/2}ds$, 于是 $H(u)\in H_0^1(\Omega)$ 且
\begin{equation*}
\abs{\nabla H(u)}^2=\abs{\nabla u}^2G'(u)=0.
\end{equation*}
由此推出\footnote{注意若 $f\in W_0^{1,p}(\Omega)$, $1\le p<\infty$ 且
$\nabla f=0$ 在 $\Omega$ 中, 则 $f=0$ 在 $\Omega$ 中.事实上, 设 $\bar f$
为 $f$ 在 $\Omega$ 外按 $0$ 延拓所得; 则 $\bar f\in W^{1,p}(\R^N)$ 且
$\nabla\bar f=\nabla f=0$ (见命题  \ref{prop:9.18}).于是 $\bar f$ 为常数
(见注记 \ref{rem:9.7}), 而由 $\bar f\in L^p(\R^N)$ 得
$\bar f\equiv0$.} $H(u)=0$ 在 $\Omega$ 中, 从而 $u\le0$ 在 $\Omega$ 中.

现在我们以下述形式证明 \eqref{eq:9.80}:
\begin{equation*}
[f\le0\ \text{在}\ \Omega\ \text{中}]
\Rightarrow\left[u\le\sup_\Gamma u\ \text{在}\ \Omega\ \text{中}\right].
\tag*{(80$'$)}
\end{equation*}
令 $K=\sup_\Gamma u$; 则 $(u-K)$ 满足 \eqref{eq:9.78}, 因为 $a_0\equiv0$
且 $(u-K)\in H^1(\Omega)$ (由于 $\Omega$ 有界).应用
(79$'$), 我们得到 $u-K\le0$ 在 $\Omega$ 中, 即
(80$'$).最后, \eqref{eq:9.81} 由 \eqref{eq:9.80} 与
(80$'$) 推出.
\end{proof}

\begin{proposition}[Neumann 问题的极大值原理]
\label{prop:9.30}
设 $f\in L^2(\Omega)$, $u\in H^1(\Omega)$ 使得
\begin{equation*}
\int_\Omega\nabla u\cdot\nabla\varphi+\int_\Omega u\varphi
=\int_\Omega f\varphi\quad\forall\varphi\in H^1(\Omega).
\end{equation*}
则对几乎一切 $x\in\Omega$,
\begin{equation*}
\inf_\Omega f\le u(x)\le\sup_\Omega f.
\end{equation*}
\end{proposition}

\begin{proof}
仿照定理  \ref{thm:9.27} 的证明.
\end{proof}

\section{特征函数与谱分解}
\label{sec:9.8}

本节中我们假设 $\Omega$ 是有界开集.

\begin{keytheorem}
\label{thm:9.31}
存在 $L^2(\Omega)$ 的 Hilbert 基 $(e_n)_{n\ge1}$ 以及实数列
$(\lambda_n)_{n\ge1}$, 使得 $\lambda_n>0$ $\forall n$ 且
$\lambda_n\to+\infty$, 并且
\begin{equation}
e_n\in H_0^1(\Omega)\cap C^\infty(\Omega),
\label{eq:9.82}
\end{equation}
\begin{equation}
-\Delta e_n=\lambda_ne_n\quad\text{在}\ \Omega\ \text{中}.
\label{eq:9.83}
\end{equation}
我们说 $\lambda_n$ 是 $-\Delta$ (带 Dirichlet 边界条件) 的\emph{特征值},
$e_n$ 是相应的\emph{特征函数}.
\end{keytheorem}

\begin{proof}
给定 $f\in L^2(\Omega)$, 令 $u=Tf$ 为问题
\begin{equation}
\int_\Omega\nabla u\cdot\nabla\varphi=\int_\Omega f\varphi
\quad\forall\varphi\in H_0^1(\Omega)
\label{eq:9.84}
\end{equation}
的唯一解 $u\in H_0^1(\Omega)$.我们把 $T$ 看作从 $L^2(\Omega)$ 到
$L^2(\Omega)$ 的算子.则 $T$ 是自伴紧算子 (重复定理  \ref{thm:8.22} 的证明,
并利用 $H_0^1(\Omega)\subset L^2(\Omega)$ 为紧嵌入这一事实).另一方面,
% 原书此处印作 "repeat the proof of Theorem 8.21" (印刷页 311); 第~
% \ref{chap:8}~章中
% 8.21 是命题 (Neumann 情形的极大值原理), 讲特征函数与 Hilbert 基的是定理
% 8.22. 依技能规范保留原样并记录, 不擅自更正.
$N(T)=\{0\}$ 且 $(Tf,f)_{L^2}\ge0$ $\forall f\in L^2$.我们推出 (应用定理
 \ref{thm:6.11}) $L^2$ 存在由 $T$ 的特征函数组成的 Hilbert 基 $(e_n)$,
相应的特征值 $(\mu_n)$ 满足 $\mu_n>0$ $\forall n$ 且 $\mu_n\to0$.于是
$e_n\in H_0^1(\Omega)$ 且
\begin{equation*}
\int_\Omega\nabla e_n\cdot\nabla\varphi
=\frac1{\mu_n}\int_\Omega e_n\varphi\quad\forall\varphi\in H_0^1(\Omega).
\end{equation*}
换言之, $e_n$ 是 \eqref{eq:9.83} 的弱解, 其中 $\lambda_n=1/\mu_n$.由第~\ref{sec:9.6}~节的正则性结果 (见注记 \ref{rem:9.25}) 我们知道
$e_n\in H^2(\omega)$ 对每个 $\omega\subset\subset\Omega$ 成立.由此推出
$e_n\in H^4(\omega)$ 对每个 $\omega\subset\subset\Omega$ 成立, 然后
$e_n\in H^6(\omega)$ 对每个 $\omega\subset\subset\Omega$ 成立, 等等.于是
$e_n\in\cap_{m\ge1}H^m(\omega)$ 对一切 $\omega\subset\subset\Omega$ 成立.
因此 $e_n\in C^\infty(\omega)$ 对一切 $\omega\subset\subset\Omega$ 成立, 即
$e_n\in C^\infty(\Omega)$.
\end{proof}

\begin{remark}
\label{rem:9.28}
在定理  \ref{thm:9.31} 的假设下, 序列 $(e_n/\sqrt{\lambda_n})$ 是赋以标量积
$\int_\Omega\nabla u\cdot\nabla v$ 的 $H_0^1(\Omega)$ 的 Hilbert 基, 而
$(e_n/\sqrt{\lambda_n+1})$ 是赋以标量积
$\int_\Omega(\nabla u\cdot\nabla v+uv)$ 的 $H_0^1(\Omega)$ 的 Hilbert 基.
事实上, 显然 $(e_n/\sqrt{\lambda_n})$ 在 $H_0^1(\Omega)$ 中正交规范 (利用
\eqref{eq:9.83}).还须验证 $e_n$ 张成的向量空间在 $H_0^1(\Omega)$ 中稠密.
为此, 设 $f\in H_0^1(\Omega)$ 使得 $(f,e_n)_{H_0^1}=0$ $\forall n$.我们要
证明 $f=0$.由 \eqref{eq:9.83} 有 $\lambda_n\int_\Omega e_nf=0$ $\forall n$,
从而 $f=0$ (因为 $(e_n)$ 是 $L^2(\Omega)$ 的 Hilbert 基).
\end{remark}

\begin{remark}
\label{rem:9.29}
在定理  \ref{thm:9.31} 的假设下 (对一般的有限区域 $\Omega$), 可以证明
$e_n\in L^\infty(\Omega)$.另一方面, 若 $\Omega$ 属于 $C^\infty$ 类, 则
$e_n\in C^\infty(\overline\Omega)$; 这容易由定理  \ref{thm:9.25} 推出.
\end{remark}

\begin{remark}
\label{rem:9.30}
设 $a_{ij}\in L^\infty(\Omega)$ 满足椭圆性条件 \eqref{eq:9.36}, 且
$a_0\in L^\infty(\Omega)$.则存在 $L^2(\Omega)$ 的 Hilbert 基 $(e_n)$ 以及
实数列 $(\lambda_n)$, $\lambda_n\to+\infty$, 使得 $e_n\in H_0^1(\Omega)$ 且
\begin{equation*}
\int_\Omega\sum_{i,j}a_{ij}\frac{\partial e_n}{\partial x_i}
\frac{\partial\varphi}{\partial x_j}+\int_\Omega a_0e_n\varphi
=\lambda_n\int_\Omega e_n\varphi\quad\forall\varphi\in H_0^1(\Omega).
\end{equation*}
\end{remark}

\section*{第 9 章评注}
\addcontentsline{toc}{section}{第 9 章评注}

本章是 Sobolev 空间理论与椭圆方程理论的引论.希望深入钻研这个广阔课题的
读者可以查阅一份内容丰富的文献目录; 我们谨援引其中一部分:
S.~Agmon \bref{Agmon}{1}, L.~Bers--F.~John--M.~Schechter
\bref{BersJohnSchechter}{1}, J.-L.~Lions \bref{Lions}{1},
J.-L.~Lions--E.~Magenes \bref{LionsMagenes}{1}, A.~Friedman
\bref{Friedman}{2}, M.~Miranda \bref{Miranda}{1}, G.~Folland
\bref{Folland}{1}, F.~Treves \bref{Treves}{4}, R.~Adams \bref{Adams}{1},
D.~Gilbarg--N.~Trudinger \bref{GilbargTrudinger}{1},
G.~Stampacchia \bref{Stampacchia}{1}, R.~Courant--D.~Hilbert
\bref{CourantHilbert}{1} 第 2 卷, H.~Weinberger \bref{Weinberger}{1},
L.~Nirenberg \bref{Nirenberg}{1}, E.~Giusti \bref{Giusti}{2},
L.~C.~Evans \bref{Evans}{1}, M.~Giaquinta \bref{Giaquinta}{1},
E.~Lieb--M.~Loss \bref{Lieb}{1}, M.~Taylor \bref{Taylor}{1},
W.~Ziemer \bref{Ziemer}{1}, O.~Ladyzhenskaya--N.~Uraltseva
\bref{Ladyzhenskaya}{1}, N.~Krylov \bref{Krylov}{1}, \bref{Krylov}{2},
V.~Maz'ja \bref{Mazja}{1}, C.~Morrey \bref{Morrey}{1},
Y.~Z.~Chen--L.~C.~Wu \bref{ChenWu}{1}, E.~DiBenedetto
\bref{DiBenedetto}{1}, Q.~Han--F.~H.~Lin \bref{HanLin}{1}, J.~Jost
\bref{Jost}{1}, W.~Strauss \bref{Strauss}{1}, 以及这些著作中援引的文献.

\cmtnum{1.} 在第~
\ref{chap:9}~章中我们常常假设 $\Omega$ 属于 $C^1$ 类; 这排除了例如带
``角''的区域.在各种情形下都可以减弱这个假设, 用某种``异乎寻常''的条件来
代替它: $\Omega$ 分片属于 $C^1$ 类, $\Omega$ 是 Lipschitz 区域, $\Omega$
具有锥性质, $\Omega$ 具有线段性质, 等等; 参见例如 R.~Adams \bref{Adams}{1}
与 S.~Agmon \bref{Agmon}{1}.

\cmtnum{2.} 定理  \ref{thm:9.7} (延拓算子的存在性) 可以借助反射延拓技巧的
适当推广, 移植到空间 $W^{m,p}(\Omega)$ ($\Omega$ 属于 $C^m$ 类); 参见例如
R.~Adams \bref{Adams}{1} 与 S.~Agmon \bref{Agmon}{1}.

\cmtnum{3.} 一些涉及 Sobolev 范数的非常有用的不等式.

\cmtnum{$\bullet$\ A. Poincaré--Wirtinger 不等式.}

设 $\Omega$ 为 $C^1$ 类连通开集, $1\le p\le\infty$.则存在常数 $C$ 使得
\begin{equation*}
\norm{u-\bar u}_p\le C\norm{\nabla u}_p\quad\forall u\in W^{1,p}(\Omega),
\end{equation*}
其中 $\bar u=\frac1{\abs{\Omega}}\int_\Omega u$.由此, 利用 Sobolev 不等式
推出: 若 $p<N$, 则
\begin{equation*}
\norm{u-\bar u}_{p^\star}\le C\norm{\nabla u}_p
\quad\forall u\in W^{1,p}(\Omega).
\end{equation*}

\cmtnum{$\bullet$\ B. Hardy 不等式.}

设 $\Omega$ 为 $C^1$ 类有界开集, $1<p<\infty$.令
$d(x)=\dist(x,\Gamma)$.则存在常数 $C$ 使得
\begin{equation*}
\left\|\frac ud\right\|_p\le C\norm{\nabla u}_p
\quad\forall u\in W_0^{1,p}(\Omega).
\end{equation*}
反之,
\begin{equation*}
[u\in W^{1,p}(\Omega)\ \text{且}\ (u/d)\in L^p(\Omega)]
\Rightarrow[u\in W_0^{1,p}(\Omega)];
\end{equation*}
见 J.-L.~Lions--E.~Magenes \bref{LionsMagenes}{1}.

\cmtnum{$\bullet$\ C. Gagliardo--Nirenberg 插值不等式.}

这里我们只举出几个在应用中经常遇到的例子.一般情形见
L.~Nirenberg \bref{Nirenberg}{1} 或 A.~Friedman \bref{Friedman}{2}.

为确定起见, 设 $\Omega\subset\R^N$ 为正则有界开集.

\noindent\emph{例 1.} 设 $u\in L^p(\Omega)\cap W^{2,r}(\Omega)$, 其中
$1\le p\le\infty$, $1\le r\le\infty$.则 $u\in W^{1,q}(\Omega)$, 其中 $q$ 是
$p$ 与 $r$ 的调和平均, 即 $\frac1q=\frac12(\frac1p+\frac1r)$, 并且
\begin{equation*}
\norm{Du}_{L^q}\le C\norm{u}_{W^{2,r}}^{1/2}\norm{u}_{L^p}^{1/2}.
\end{equation*}
特例:

\noindent (a) $p=\infty$, 从而 $q=2r$.我们有
\begin{equation*}
\norm{Du}_{L^q}\le C\norm{u}_{W^{2,r}}^{1/2}\norm{u}_{L^\infty}^{1/2}.
\end{equation*}
除别的用处之外, 这个不等式可以用来证明 $W^{2,r}\cap L^\infty$ 是一个代数,
也就是说,
\begin{equation*}
u,v\in W^{2,r}\cap L^\infty\Rightarrow uv\in W^{2,r}\cap L^\infty
\end{equation*}
(这个性质对 $W^{m,r}\cap L^\infty$ ($m$ 为整数, $m\ge2$) 仍然成立).

\noindent (b) $p=q=r$.我们有
\begin{equation*}
\norm{Du}_{L^p}\le C\norm{u}_{W^{2,p}}^{1/2}\norm{u}_{L^p}^{1/2},
\end{equation*}
由此特别推出
\begin{equation*}
\norm{Du}_{L^p}\le\varepsilon\norm{D^2u}_{L^p}
+C_\varepsilon\norm{u}_{L^p}\quad\forall\varepsilon>0.
\end{equation*}

\noindent\emph{例 2.} 设 $1\le q\le p<\infty$.则
\begin{equation}
\norm{u}_{L^p}\le C\norm{u}_{L^q}^{1-a}\norm{u}_{W^{1,N}}^a
\quad\forall u\in W^{1,N}(\Omega),\ \text{其中}\ a=1-(q/p).
\label{eq:9.85}
\end{equation}
我们指出经常用到的一个特例
\begin{equation*}
N=2,\quad p=4,\quad q=2,\quad\text{且}\ a=1/2,
\end{equation*}
也就是说,
\begin{equation*}
\norm{u}_{L^4}\le C\norm{u}_{L^2}^{1/2}\norm{u}_{H^1}^{1/2}
\quad\forall u\in H^1(\Omega).
\end{equation*}
在这方面我们注意到, 还有通常的插值不等式 (第~
\ref{chap:4}~章注记 2)
\begin{equation*}
\norm{u}_{L^p}\le\norm{u}_{L^q}^{1-a}\norm{u}_{L^\infty}^a
\quad\text{其中}\ a=1-(q/p),
\end{equation*}
但它不能推出 \eqref{eq:9.85}, 因为 $W^{1,N}$ 不包含在 $L^\infty$ 中.

\noindent\emph{例 3.} 设 $1\le q\le p\le\infty$ 且 $r>N$.则
\begin{equation}
\norm{u}_{L^p}\le C\norm{u}_{L^q}^{1-a}\norm{u}_{W^{1,r}}^a
\quad\forall u\in W^{1,r}(\Omega),
\label{eq:9.86}
\end{equation}
其中 $a=\left(\frac1q-\frac1p\right)\Big/\left(\frac1q+\frac1N-\frac1r\right)$.

\cmtnum{$\bullet$\ 4.} 下述性质有时是有用的.设 $u\in W^{1,p}(\Omega)$,
$1\le p\le\infty$, $\Omega$ 为任意开集.则 $\nabla u=0$ 在集合
$\{x\in\Omega;\ u(x)=k\}$ 上几乎处处成立, 其中 $k$ 为任意常数.

\cmtnum{$\star$5.} 当 $p>N$ 时, $W^{1,p}(\Omega)$ 中的函数在 $\Omega$ 中
几乎处处按通常意义可微.更精确地说, 设 $u\in W^{1,p}(\Omega)$, $p>N$.则
存在零测集 $A\subset\Omega$ 使得
\begin{equation*}
\lim_{h\to0}\frac{u(x+h)-u(x)-h\cdot\nabla u(x)}{\abs{h}}=0
\quad\forall x\in\Omega\setminus A.
\end{equation*}
当 $u\in W^{1,p}(\Omega)$ 且 $p\le N$ ($N>1$) 时这个性质不成立.关于这个
问题可查阅 E.~Stein \bref{Stein}{1} (第~
\ref{chap:8}~章).

\cmtnum{6.} \emph{分数阶 Sobolev 空间.}

可以定义一族介于 $L^p(\Omega)$ 与 $W^{1,p}(\Omega)$ 之间的空间.更精确地说,
若 $0<s<1$ ($s\in\R$) 且 $1\le p<\infty$, 令
\begin{equation*}
W^{s,p}(\Omega)=\left\{u\in L^p(\Omega);\
\frac{\abs{u(x)-u(y)}}{\abs{x-y}^{s+(N/p)}}\in L^p(\Omega\times\Omega)
\right\},
\end{equation*}
赋以自然的范数.令 $H^s(\Omega)=W^{s,2}(\Omega)$.关于这些空间的研究, 参见
例如 R.~Adams \bref{Adams}{1}, J.-L.~Lions--E.~Magenes
\bref{LionsMagenes}{1}, P.~Malliavin \bref{Malliavin}{1},
H.~Triebel \bref{Triebel}{1} 以及 L.~Grafakos \bref{Grafakos}{1}.
空间 $W^{s,p}(\Omega)$ 也可以定义为 $W^{1,p}$ 与 $L^p$ 之间的插值空间;
当 $p=2$ 且 $\Omega=\R^N$ 时还可以用 Fourier 变换来定义.最后, 对实数
$s$ (非整数) 且 $s>1$, 我们如下定义 $W^{s,p}(\Omega)$: 写
$s=m+\sigma$, 其中 $m$ 是 $s$ 的整数部分, 并令
\begin{equation*}
W^{s,p}(\Omega)=\{u\in W^{m,p}(\Omega);\
D^\alpha u\in W^{\sigma,p}(\Omega)\ \forall\alpha,\ \abs{\alpha}=m\}.
\end{equation*}
借助局部坐标卡还可以定义 $W^{s,p}(\Gamma)$, 其中 $\Gamma$ 是光滑流形
(例如正则开集的边界).这些空间在迹理论中起重要作用 (见评注 7).

\cmtnum{$\bullet$\ 7.} \emph{迹理论.}

设 $1\le p<\infty$.我们先从一个基本引理开始.

\begin{lemma}
\label{lem:9.9}
设 $\Omega=\R^N_+$.则存在常数 $C$ 使得
\begin{equation*}
\left(\int_{\R^{N-1}}\abs{u(x',0)}^p\,dx'\right)^{1/p}
\le C\norm{u}_{W^{1,p}(\Omega)}\quad\forall u\in C_c^1(\R^N).
\end{equation*}
\end{lemma}

\begin{proof}
令 $G(t)=\abs{t}^{p-1}t$, 并设 $u\in C_c^1(\R^N)$.我们有
\begin{equation*}
\begin{aligned}
G(u(x',0))&=-\int_0^{+\infty}\frac{\partial}{\partial x_N}
G(u(x',x_N))\,dx_N\\
&=-\int_0^{+\infty}G'(u(x',x_N))
\frac{\partial u}{\partial x_N}(x',x_N)\,dx_N.
\end{aligned}
\end{equation*}
于是
\begin{equation*}
\begin{aligned}
\abs{u(x',0)}^p
&\le p\int_0^\infty\abs{u(x',x_N)}^{p-1}
\left|\frac{\partial u}{\partial x_N}(x',x_N)\right|dx_N\\
&\le C\left(\int_0^\infty\abs{u(x',x_N)}^p\,dx_N
+\int_0^\infty\left|\frac{\partial u}{\partial x_N}(x',x_N)\right|^p
dx_N\right),
\end{aligned}
\end{equation*}
对 $x'\in\R^{N-1}$ 积分即得结论.
\end{proof}

由引理  \ref{lem:9.9} 可以推出: 由
$u\mapsto u|_\Gamma$ ($\Gamma=\partial\Omega=\R^{N-1}\times\{0\}$) 定义在
$C_c^1(\R^N)$ 上, 取值于 $L^p(\Gamma)$ 的映射, 按稠密性可以延拓为从
$W^{1,p}(\Omega)$ 到 $L^p(\Gamma)$ 的有界线性算子.这个算子按定义就是 $u$ 在
$\Gamma$ 上的\emph{迹}; 它也记作 $u|_\Gamma$.

我们指出, $L^p(\R^N_+)$ 与 $W^{1,p}(\R^N_+)$ 之间有一个本质区别:
$L^p(\R^N_+)$ 中的函数在 $\Gamma$ 上没有迹.利用局部坐标卡, 容易想象当
$\Omega$ 是 $\R^N$ 中的正则开集 (例如 $\Omega$ 属于 $C^1$ 类且 $\Gamma$ 有界)
时, 如何对函数 $u\in W^{1,p}(\Omega)$ 定义它在 $\Gamma=\partial\Omega$ 上的
迹.在这种情形下 $u|_\Gamma\in L^p(\Gamma)$ (关于面测度 $d\sigma$).迹的最
重要的性质如下:

\noindent (i) 若 $u\in W^{1,p}(\Omega)$, 则实际上
$u|_\Gamma\in W^{1-(1/p),p}(\Gamma)$ 且
\begin{equation*}
\norm{u|_\Gamma}_{W^{1-(1/p),p}(\Gamma)}\le C\norm{u}_{W^{1,p}(\Omega)}
\quad\forall u\in W^{1,p}(\Omega).
\end{equation*}
此外, \index{jisuanzi@迹算子 (trace operator)}
迹算子 $u\mapsto u|_\Gamma$ 是从 $W^{1,p}(\Omega)$ 到
$W^{1-(1/p),p}(\Gamma)$ 上的满射.

\noindent (ii) 迹算子的核是 $W_0^{1,p}(\Omega)$, 即
\begin{equation*}
W_0^{1,p}(\Omega)=\{u\in W^{1,p}(\Omega);\ u|_\Gamma=0\}.
\end{equation*}

\noindent (iii) 我们有 Green 公式
\begin{equation*}
\int_\Omega\frac{\partial u}{\partial x_i}v
=-\int_\Omega u\frac{\partial v}{\partial x_i}
+\int_\Gamma uv(\vec n\cdot\vec e_i)\,d\sigma\quad\forall u,v\in H^1(\Omega),
\end{equation*}
其中 $\vec n$ 是 $\Gamma$ 的单位外法向量.注意面积分是有意义的, 因为
$u,v\in L^2(\Gamma)$.

同样地, 对函数 $u\in W^{2,p}(\Omega)$ 我们可以谈论 $\frac{\partial u}{\partial n}$:
令 $\frac{\partial u}{\partial n}=(\nabla u)|_\Gamma\cdot\vec n$, 这是有意义的,
因为 $(\nabla u)|_\Gamma\in L^p(\Gamma)^N$, 且
$\frac{\partial u}{\partial n}\in L^p(\Gamma)$ (实际上
$\frac{\partial u}{\partial n}\in W^{1-(1/p),p}(\Gamma)$).Green 公式也成立:
\begin{equation*}
-\int_\Omega(\Delta u)v=\int_\Omega\nabla u\cdot\nabla v
-\int_\Gamma\frac{\partial u}{\partial n}v\,d\sigma\quad\forall u,v\in H^2(\Omega).
\end{equation*}

\noindent (iv) 算子
$u\mapsto\{u|_\Gamma,\frac{\partial u}{\partial n}\}$ 是从 $W^{2,p}(\Omega)$
到 $W^{2-(1/p),p}(\Gamma)\times W^{1-(1/p),p}(\Gamma)$ 上的有界线性满射.
关于这些问题, $p=2$ 的情形见 J.-L.~Lions--E.~Magenes
\bref{LionsMagenes}{1} (关于 $p\neq2$ 的情形见其中援引的文献).

\cmtnum{8.} \emph{2$m$ 阶算子与椭圆组.}

第~
\ref{chap:9}~章中证明的存在性与正则性结果可以推广到 $2m$ 阶椭圆算子与椭圆组.
\footnote{但极大值原理不能推广, 除非在非常特殊的情形.} 本质要素之一是
G\aa rding 不等式.关于这些问题, 参见 S.~Agmon \bref{Agmon}{1},
J.-L.~Lions--E.~Magenes \bref{LionsMagenes}{1},
S.~Agmon--A.~Douglis--L.~Nirenberg \bref{AgmonDouglisNirenberg}{1}.
$2m$ 阶算子与某些组在力学与物理学中起重要作用.我们特别指出双调和算子
$\Delta^2$ (板的弯曲理论), 弹性力学组与 Stokes 组 (流体力学); 参见例如
Ph.~Ciarlet \bref{Ciarlet}{1}, G.~Duvaut--J.-L.~Lions
\bref{DuvautLions}{1}, R.~Temam \bref{Temam}{1},
J.~Ne\v cas--L.~Hlava\v cek \bref{NecasHlavacek}{1}, M.~Gurtin
\bref{Gurtin}{1}.

\cmtnum{9.} \emph{$L^p$ 与 $C^{0,\alpha}$ 空间中的正则性.}

第~
\ref{chap:9}~章中关于 $p=2$ 证明的正则性定理可以推广到 $p\neq2$ 的情形.

\begin{keytheorem}[Agmon--Douglis--Nirenberg]
\label{thm:9.32}
设 $\Omega$ 属于 $C^2$ 类且 $\Gamma$ 有界, $1<p<\infty$.则对一切
$f\in L^p(\Omega)$, 方程
\begin{equation}
-\Delta u+u=f\quad\text{在}\ \Omega\ \text{中}
\label{eq:9.87}
\end{equation}
存在唯一的解 $u\in W^{2,p}(\Omega)\cap W_0^{1,p}(\Omega)$.此外, 若
$\Omega$ 属于 $C^{m+2}$ 类且 $f\in W^{m,p}(\Omega)$ ($m\ge1$ 为整数), 则
\begin{equation*}
u\in W^{m+2,p}(\Omega)\quad\text{且}\quad
\norm{u}_{W^{m+2,p}}\le C\norm{f}_{W^{m,p}}.
\end{equation*}
把 \eqref{eq:9.87} 换成带光滑系数的二阶椭圆方程, 有类似的结果.定理
 \ref{thm:9.32} 的证明比 $p=2$ 的情形 (定理  \ref{thm:9.25}) 复杂得多.
``古典''的方法本质上依赖于两个要素:
\end{keytheorem}

\noindent (a) 利用基本解给出 $u$ 的显式表示公式.例如, 若 $\Omega=\R^3$,
则 \eqref{eq:9.87} 的解由 $u=G\star f$ 给出, 其中
$G(x)=\frac{c}{\abs{x}}e^{-\abs{x}}$.于是形式上有
$\frac{\partial^2u}{\partial x_i\partial x_j}
=\frac{\partial^2G}{\partial x_i\partial x_j}\star f$;``不幸的是''
$\frac{\partial^2G}{\partial x_i\partial x_j}$ 不属于 $L^1(\R^3)$,\footnote{
但几乎属于!} 因为在 $x=0$ 处有奇性, 因此不能对卷积积应用初等估计 (例如定理
4.15).

\noindent (b) 为克服这个困难, 人们使用 Calderón--Zygmund 的 $L^p$ 奇异积分
理论 (参见例如 E.~Stein \bref{Stein}{1} 与 L.~Bers--F.~John--M.~Schechter
\bref{BersJohnSchechter}{1}).

\noindent\emph{注意:} 定理  \ref{thm:9.32} 的结论对 $p=1$ 与 $p=\infty$
不成立.

在 H\"older 空间框架中的另一个基本正则性结果如下.\footnote{回忆当
$0<\alpha<1$ 且 $m$ 为整数时,
\begin{equation*}
C^{0,\alpha}(\Omega)=\left\{u\in C(\Omega);\
\sup_{\substack{x,y\in\Omega\\x\neq y}}
\frac{\abs{u(x)-u(y)}}{\abs{x-y}^\alpha}<\infty\right\},
\end{equation*}
且 $C^{m,\alpha}(\Omega)=\{u\in C(\Omega);\ D^\beta u\in C^{0,\alpha}(\Omega)\
\forall\beta,\ \abs{\beta}=m\}$.}

\begin{keytheorem}[Schauder]
\label{thm:9.33}
设 $\Omega$ 有界且属于 $C^{2,\alpha}$ 类, $0<\alpha<1$.则对每个
$f\in C^{0,\alpha}(\Omega)$, 问题
\begin{equation}
\begin{cases}
-\Delta u+u=f & \text{在}\ \Omega\ \text{中},\\
u=0 & \text{在}\ \Gamma\ \text{上}
\end{cases}
\label{eq:9.88}
\end{equation}
存在唯一的解 $u\in C^{2,\alpha}(\Omega)$.此外, 若 $\Omega$ 属于
$C^{m+2,\alpha}$ 类 ($m\ge1$ 为整数) 且 $f\in C^{m,\alpha}(\Omega)$, 则
$u\in C^{m+2,\alpha}(\Omega)$ 且
$\norm{u}_{C^{m+2,\alpha}}\le C\norm{f}_{C^{m,\alpha}}$.
\end{keytheorem}

把 \eqref{eq:9.88} 换成带光滑系数的二阶椭圆算子, 有类似的结果.定理
 \ref{thm:9.33} 的证明——如同定理  \ref{thm:9.32} 的证明一样——依赖于 $u$ 的
显式表示以及 H\"older, Korn, Lichtenstein, Giraud 的 $C^{0,\alpha}$ 空间
奇异积分理论.关于这个课题, 参见
S.~Agmon--A.~Douglis--L.~Nirenberg \bref{AgmonDouglisNirenberg}{1},
L.~Bers--F.~John--M.~Schechter \bref{BersJohnSchechter}{1},
C.~Morrey \bref{Morrey}{1}, D.~Gilbarg--N.~Trudinger
\bref{GilbargTrudinger}{1}.另一种避免奇异积分理论的方法由 Campanato 与
Stampacchia 提出 (参见例如 Y.~Z.~Chen--L.~C.~Wu \bref{ChenWu}{1} 与
E.~Giusti \bref{Giusti}{2}).其它初等方法由 A.~Brandt \bref{Brandt}{1}
(仅基于极大值原理) 与 L.~Simon \bref{Simon}{2} 发展起来.

设 $\Omega$ 为有界正则开集, $f\in C(\Omega)$.由定理  \ref{thm:9.32},
存在 $u\in W^{2,p}(\Omega)\cap W_0^{1,p}(\Omega)$ (对一切 $1<p<\infty$),
它是 \eqref{eq:9.87} 的唯一解.特别地, $u\in C^{1,\alpha}(\Omega)$ 对一切
$0<\alpha<1$ 成立 (由 Morrey 定理 (定理  \ref{thm:9.12})).一般说来, $u$
不属于 $C^2$, 甚至不属于 $W^{2,\infty}$.这就解释了为什么人们常常避免在
空间 $L^1(\Omega)$, $L^\infty(\Omega)$ 与 $C(\Omega)$ 中工作——对于这些
空间我们没有最优的正则性结果.

定理  \ref{thm:9.32} 与  \ref{thm:9.33} 可以推广到 $2m$ 阶椭圆算子与椭圆组;
见 S.~Agmon--A.~Douglis--L.~Nirenberg \bref{AgmonDouglisNirenberg}{1}.
最后我们指出, 从另一个方向看, 带不连续系数的二阶椭圆方程是大量工作的
主题.我们举出下述著名结果为例.

\begin{keytheorem}[De Giorgi, Nash, Stampacchia]
\label{thm:9.34}
设 $\Omega\subset\R^N$, $N\ge2$, 为有界正则开集.假设函数
$a_{ij}\in L^\infty(\Omega)$ 满足椭圆性条件 \eqref{eq:9.36}.设
$f\in L^p(\Omega)$, $p>N/2$, 且 $u\in H_0^1(\Omega)$ 使得
\begin{equation*}
\int_\Omega\sum_{i,j}a_{ij}\frac{\partial u}{\partial x_i}
\frac{\partial\varphi}{\partial x_j}=\int_\Omega f\varphi
\quad\forall\varphi\in H_0^1(\Omega).
\end{equation*}
则 $u\in C^{0,\alpha}(\Omega)$, 其中 $0<\alpha<1$ 依赖于 $\Omega$, $a_{ij}$
与 $p$.
\end{keytheorem}

关于这些问题, 参见 G.~Stampacchia \bref{Stampacchia}{1},
D.~Gilbarg--N.~Trudinger \bref{GilbargTrudinger}{1},
O.~Ladyzhenskaya--N.~Uraltseva \bref{Ladyzhenskaya}{1} 以及
E.~Giusti \bref{Giusti}{2}.

\cmtnum{10.} \emph{变分方法的一些缺点以及如何绕开它们!}

变分方法非常容易地给出弱解的存在性.它并不总是适用, 但可以补充完善.
我们指出两个例子.设 $\Omega\subset\R^N$ 为有界正则开集.

\noindent (a) \emph{对偶方法.} 设 $f\in L^1(\Omega)$——甚至设 $f$ 是
$\Omega$ 上的一个 (Radon) 测度——并寻找问题
\begin{equation}
\begin{cases}
-\Delta u+u=f & \text{在}\ \Omega\ \text{中},\\
u=0 & \text{在}\ \Gamma\ \text{上}
\end{cases}
\label{eq:9.89}
\end{equation}
的解.只要 $N>1$, 线性泛函 $\varphi\mapsto\int_\Omega f\varphi$ 对每个
$\varphi\in H_0^1(\Omega)$ 就没有定义, 因此变分方法无效.另一方面, 可以使用
下述技巧.用 $T:L^2(\Omega)\to L^2(\Omega)$ 记算子 $f\mapsto u$ (其中 $u$ 是
\eqref{eq:9.89} 的解, 它对 $f\in L^2(\Omega)$ 存在).我们知道 $T$ 是自伴的.
另一方面 (定理  \ref{thm:9.32}), 对 $2\le p<\infty$ 有
$T:L^p(\Omega)\to W^{2,p}(\Omega)$, 而由 Sobolev 与 Morrey 定理, 当
$p>N/2$ 时 $T:L^p(\Omega)\to C^0(\Omega)$.由对偶性推出, 当 $p>N/2$ 时
\begin{equation*}
T^\star:\mathcal M(\Omega)=C^0(\Omega)^\star\to L^{p'}(\Omega).
\end{equation*}
由于 $T$ 在 $L^2$ 中自伴, $T^\star$ 是 $T$ 的延拓; 于是可以把
$u=T^\star f$ 看作 \eqref{eq:9.89} 的广义解.事实上, 若
$f\in L^1(\Omega)$, 则 $u=T^\star f\in L^q(\Omega)$ 对一切
$q<N/(N-2)$ 成立; $u$ 是 \eqref{eq:9.89} 在
\begin{equation*}
-\int_\Omega u\Delta\varphi+\int_\Omega u\varphi=\int_\Omega f\varphi
\quad\forall\varphi\in C^2(\Omega),\ \varphi=0\ \text{在}\ \Gamma\ \text{上}
\end{equation*}
意义下的唯一 (\emph{很}) 弱解.按照同样的想法, 可以研究 $f$ 给定于
$H^{-m}(\Omega)$ 时的 \eqref{eq:9.89}; 见 J.-L.~Lions--E.~Magenes
\bref{LionsMagenes}{1}.

\noindent (b) \emph{稠密性方法.} 设 $g\in C(\Gamma)$ 并寻找问题
\begin{equation}
\begin{cases}
-\Delta u+u=0 & \text{在}\ \Omega\ \text{中},\\
u=g & \text{在}\ \Gamma\ \text{上}
\end{cases}
\label{eq:9.90}
\end{equation}
的解.一般说来, 若 $g\in C(\Gamma)$, 并不存在函数 $\tilde g\in H^1(\Omega)$
使得 $\tilde g|_\Gamma=g$ (见评注 7, 并注意 $C(\Gamma)$ 不包含在
$H^{1/2}(\Gamma)$ 中).因此不可能在 $H^1(\Omega)$ 中寻找 \eqref{eq:9.90} 的
解: 变分方法无效.然而, 我们有下述结果.

\begin{keytheorem}
\label{thm:9.35}
\eqref{eq:9.90} 存在唯一的解 $u\in C(\overline\Omega)\cap C^\infty(\Omega)$.
\end{keytheorem}

\begin{proof}
固定 $\tilde g\in C_c(\R^N)$ 使得 $\tilde g|_\Gamma=g$; 由
Tietze--Urysohn 定理 (参见例如 J.~Dieudonn\'e \bref{Dieudonne}{1},
J.~Dugundji \bref{Dugundji}{1}, J.~Munkres \bref{Munkres}{1})
这样的 $\tilde g$ 存在.设 $(\tilde g_n)$ 是 $C_c^\infty(\R^N)$ 中的序列,
使得 $\tilde g_n\to g$ 在 $\R^N$ 上一致收敛.令
$g_n=\tilde g_n|_\Gamma$.应用变分方法与正则性结果, 我们看到问题
\begin{equation*}
\begin{cases}
-\Delta u_n+u_n=0 & \text{在}\ \Omega\ \text{中},\\
u_n=g_n & \text{在}\ \Gamma\ \text{上}
\end{cases}
\end{equation*}
存在古典解 $u_n\in C^2(\Omega)$.由极大值原理 (推论  \ref{cor:9.28}) 我们有
\begin{equation*}
\norm{u_m-u_n}_{L^\infty(\Omega)}\le\norm{g_m-g_n}_{L^\infty(\Gamma)}.
\end{equation*}
于是 $(u_n)$ 是 $C(\overline\Omega)$ 中的 Cauchy 序列, 且
$u_n\to u$ 在 $C(\overline\Omega)$ 中.显然有
\begin{equation*}
\int_\Omega u(-\Delta\varphi+\varphi)=0\quad\forall\varphi\in C_c^\infty(\Omega),
\end{equation*}
从而 $u\in C^\infty(\Omega)$ (见注记 \ref{rem:9.25}).于是
$u\in C(\overline\Omega)\cap C^\infty(\Omega)$ 满足 \eqref{eq:9.90}.
\eqref{eq:9.90} 的解的唯一性由极大值原理推出 (见注记 \ref{rem:9.27}).
\end{proof}

\begin{starremark}
\label{rem:9.31}
在定理  \ref{thm:9.35} 中, 假设 $\Omega$ 足够光滑是本质的.当 $\Omega$ 有
``病态''边界时, 我们会遇到位势论中的问题 (正则点, Wiener 判别法等).

求解 \eqref{eq:9.90} 的另一种方法是在位势论中很经典的 Perron 方法.定义
\begin{equation*}
u(x)=\sup\{v(x);\ v\in C(\overline\Omega)\cap C^2(\Omega),\
-\Delta v+v\le0\ \text{在}\ \Omega\ \text{中},\ v\le g\ \text{在}\
\Gamma\ \text{上}\},
\end{equation*}
并 (直接) 证明 $u$ 满足 \eqref{eq:9.90}.满足
$-\Delta v+v\le0$ 在 $\Omega$ 中且 $v\le g$ 在 $\Gamma$ 上的函数 $v$
称为 \eqref{eq:9.90} 的\emph{下解}.
\end{starremark}

\cmtnum{11.} \emph{强极大值原理.}

当 $u$ 是古典解时, 我们可以加强命题  \ref{prop:9.29} 的结论.更精确地说,
设 $\Omega$ 为连通的有界正则开集, $a_{ij}\in C^1(\Omega)$ 满足椭圆性条件
\eqref{eq:9.36}, $a_i,a_0\in C(\Omega)$, $a_0\ge0$ 在 $\Omega$ 上.

\begin{keytheorem}[Hopf]
\index{qiangjidazhiyuanli@强极大值原理 (strong maximum principle)}
\label{thm:9.36}
设 $u\in C(\overline\Omega)\cap C^2(\Omega)$ 满足
\begin{equation}
-\sum_{i,j}\frac{\partial}{\partial x_j}
\left(a_{ij}\frac{\partial u}{\partial x_i}\right)
+\sum_ia_i\frac{\partial u}{\partial x_i}+a_0u=f
\quad\text{在}\ \Omega\ \text{中}.
\label{eq:9.91}
\end{equation}
假设 $f\ge0$ 在 $\Omega$ 上.若存在 $x_0\in\Omega$ 使得
$u(x_0)=\min_{\overline\Omega}u$, 且 $u(x_0)\le0$,\footnote{若 $a_0=0$,
则假设 $u(x_0)\le0$ 是多余的.} 则 $u$ 在 $\Omega$ 中为常数 (并且进一步
$f=0$ 在 $\Omega$ 中).
\end{keytheorem}

关于证明, 参见例如 L.~Bers--F.~John--M.~Schechter
\bref{BersJohnSchechter}{1}, D.~Gilbarg--N.~Trudinger
\bref{GilbargTrudinger}{1}, M.~Protter--H.~Weinberger
\bref{ProtterWeinberger}{1} 以及 P.~Pucci--J.~Serrin
\bref{PucciSerrin}{1}.

\begin{keycorollary}
\label{cor:9.37}
设 $u\in C(\overline\Omega)\cap C^2(\Omega)$ 满足 \eqref{eq:9.91},
其中 $f\ge0$ 在 $\Omega$ 上.假设 $u\ge0$ 在 $\Gamma$ 上.则
\begin{itemize}
	\item 或者 $u>0$ 在 $\Omega$ 中,
	\item 或者 $u\equiv0$ 在 $\Omega$ 中.
\end{itemize}
\end{keycorollary}

关于与极大值原理相关的其它结果 (Harnack 不等式等), 参见例如
G.~Stampacchia \bref{Stampacchia}{1},
D.~Gilbarg--N.~Trudinger \bref{GilbargTrudinger}{1},
M.~Protter--H.~Weinberger \bref{ProtterWeinberger}{1},
R.~Sperb \bref{Sperb}{1} 以及 P.~Pucci--J.~Serrin
\bref{PucciSerrin}{1}.

\cmtnum{12.} \emph{Laplace--Beltrami 算子.}

定义在 Riemann 流形 (带边界或不带边界) 上的椭圆算子, 特别是
Laplace--Beltrami 算子, 在微分几何与物理学中起重要作用; 参见例如
Y.~Choquet--C.~Dewitt--M.~Dillard \bref{ChoquetDewittDillard}{1}.

\cmtnum{13.} \emph{谱性质. 反问题.}

二阶椭圆算子的特征值与特征函数有许多值得注意的性质.这里我们举出其中
一些.设 $\Omega\subset\R^N$ 为连通的有界正则开集.设
$a_{ij}\in C^1(\Omega)$ 满足椭圆性条件 \eqref{eq:9.36}, $a_0\in C(\Omega)$.
设 $A$ 为算子
\begin{equation*}
Au=-\sum_{i,j}\frac{\partial}{\partial x_j}
\left(a_{ij}\frac{\partial u}{\partial x_i}\right)+a_0u
\end{equation*}
并带齐次 Dirichlet 条件 ($u=0$ 在 $\Gamma$ 上).用 $(\lambda_n)$ 记 $A$
的特征值按递增顺序排列的序列, 当 $n\to\infty$ 时
$\lambda_n\to+\infty$.则第一个特征值 $\lambda_1$ 的重数为 1 (这时说
$\lambda_1$ 是\emph{单}特征值),\footnote{当维数 $N\ge2$ 时, 其余特征值
的重数可以大于 1.} 且我们可以选取相应的特征函数 $e_1$ 使 $e_1>0$ 在
$\Omega$ 上; 这由 Krein--Rutman 定理推出 (见第~
\ref{chap:6}~章评注与问题 41).此外,
可以证明当 $n\to\infty$ 时 $\lambda_n\sim cn^{2/N}$, 其中 $c>0$; 见
S.~Agmon \bref{Agmon}{1}.

$\Omega$ 的几何性质\footnote{特别是当 $\Omega$ 是不带边界的 Riemann 流形
且 $A$ 是 Laplace--Beltrami 算子时.}与 $A$ 的谱之间的关系是大量研究的
主题; 参见例如 M.~Kac \bref{Kac}{1}, Marcel Berger \bref{Berger}{1},
R.~Osserman \bref{Osserman}{1}, I.~M.~Singer \bref{Singer}{1},
P.~B\'erard \bref{Berard}{1}, I.~Chavel \bref{Chavel}{1}.谱几何的目标
是仅从谱 $(\lambda_n)$ 的知识出发, ``还原''关于 $\Omega$ 的最大量的信息.

一个极其简单的问题是下述问题.设 $\Omega_1$ 与 $\Omega_2$ 是 $\R^2$ 中
两个有界区域; 假设算子 $-\Delta$ (带 Dirichlet 边界条件) 在 $\Omega_1$ 与
$\Omega_2$ 上的特征值相同.$\Omega_1$ 与 $\Omega_2$ 是否等距? 这个问题被
M.~Kac 戏称为: ``人能听出鼓的形状吗?''\footnote{因为附着在边界 $\Gamma$ 上
的膜振动的谐波是函数 $e_n(x)\sin\sqrt{\lambda_n}t$, 其中 $(\lambda_n,e_n)$
是 $-\Delta$ 带 Dirichlet 边界条件的特征值与特征函数.} 已知若 $\Omega_1$
是圆盘, 答案是肯定的.1991 年, C.~Gordon--L.~Webb--S.~Wolpert
\bref{GordonWebbWolpert}{1} 对带角的区域给出了否定的答案.Kac 问题对光滑
区域仍未解决.

另一类重要的``反问题''是仅从边界上的测量 (例如 Dirichlet-to-Neumann 映射)
或``无穷远处''的测量 (逆散射) 来确定偏微分方程中的系数与参数, 或者确定
内部物体的形状与特征.这类问题出现在许多领域 (医学成像, 地震学等); 参见
例如 G.~Uhlmann \bref{Uhlmann}{1}, C.~B.~Croke 等 \bref{CrokeEtAl}{1}.

\cmtnum{14.} \emph{退化椭圆问题.}

考虑形如
\begin{equation*}
\begin{cases}
\displaystyle-\sum_{i,j}\frac{\partial}{\partial x_j}
\left(a_{ij}\frac{\partial u}{\partial x_i}\right)
+\sum_ia_i\frac{\partial u}{\partial x_i}+a_0u=f
& \text{在}\ \Omega\ \text{中},\\[2mm]
\text{边界条件给定在}\ \Gamma\ \text{上, 或给定在}\ \Gamma\ \text{的一部分上}
\end{cases}
\end{equation*}
的问题, 其中函数 $a_{ij}$ 不满足椭圆性条件 \eqref{eq:9.36}, 而只满足
\begin{equation}
\sum_{i,j}a_{ij}(x)\xi_i\xi_j\ge0\quad\forall x\in\Omega,\ \forall\xi\in\R^N.
\tag*{(36$'$)}
\end{equation}
参见例如 J.~Kohn--L.~Nirenberg \bref{KohnNirenberg}{1},
M.~S.~Baouendi--C.~Goulaouic \bref{BaouendiGoulaouic}{1},
O.~Oleinik--E.~Radkevitch \bref{OleinikRadkevitch}{1} 的著作.

\cmtnum{15.} \emph{非线性椭圆问题.}

这是一个由几何, 力学, 物理学, 最优控制, 概率论等中数不清的问题所激发的
庞大研究领域.自 20 世纪 30 年代初 Leray 与 Schauder 的开创性工作以来,
它有了惊人的发展.我们区分几个类别:

\cmtsub{(a) 半线性问题.} 例如, 这包括形如
\begin{equation}
\begin{cases}
-\Delta u=f(x,u) & \text{在}\ \Omega\ \text{中},\\
u=0 & \text{在}\ \Gamma\ \text{上}
\end{cases}
\label{eq:9.92}
\end{equation}
的问题, 其中 $f(x,u)$ 是给定的函数.

这一类中包括所谓的\emph{分支问题}, 其中人们研究问题
\begin{equation*}
\begin{cases}
-\Delta u=f_\lambda(x,u) & \text{在}\ \Omega\ \text{中},\\
u=0 & \text{在}\ \Gamma\ \text{上}
\end{cases}
\tag*{(91$'$)}
\end{equation*}
的解集 $(\lambda,u)$ 的结构, 其中 $\lambda$ 是变化的参数.

\cmtsub{(b) 拟线性问题.} 考虑形如
\begin{equation}
\begin{cases}
\displaystyle-\sum_{i,j}\frac{\partial}{\partial x_j}
\left(a_{ij}(x,u,\nabla u)\frac{\partial u}{\partial x_i}\right)
=f(x,u,\nabla u) & \text{在}\ \Omega\ \text{中},\\[2mm]
u=0 & \text{在}\ \Gamma\ \text{上}
\end{cases}
\label{eq:9.93}
\end{equation}
的问题, 其中函数 $a_{ij}(x,u,p)$ 是椭圆的, 但可能是退化的; 例如我们有
\begin{equation*}
\sum_{i,j}a_{ij}(x,u,p)\xi_i\xi_j\ge\alpha(u,p)\abs{\xi}^2
\quad\forall x\in\Omega,\ \forall\xi\in\R^N,\ \forall u\in\R,\
\forall p\in\R^N,
\end{equation*}
其中 $\alpha(u,p)>0$ $\forall u\in\R$, $\forall p\in\R^N$, 但
$\alpha(u,p)$ 不能被一个常数 $\alpha>0$ 一致地控制在下.特别地, 著名的
\emph{极小曲面方程}属于这一类, 其中
$a_{ij}(x,u,p)=\delta_{ij}(1+\abs{p}^2)^{-1/2}$.更一般地, 人们考虑形如
\begin{equation}
F(x,u,Du,D^2u)=0
\label{eq:9.94}
\end{equation}
的\emph{完全非线性}椭圆问题, 其中矩阵
$\frac{\partial F}{\partial q_{ij}}(x,u,p,q)$ 是椭圆的 (可能是退化的).
例如, \emph{Monge--Amp\`ere 方程}属于这一类.

\cmtsub{(c) 自由边界问题.} 问题是要在一个\emph{事先未给定}的开集 $\Omega$
中求解一个线性椭圆方程.$\Omega$ 未知这一事实常常通过在 $\Gamma$ 上给定
\emph{两个}边界条件 (例如 Dirichlet 条件与 Neumann 条件) 来``补偿''.
问题在于同时找出开集 $\Omega$ 与函数 $u$, 使得 $\ldots$.

\runin{技巧:}

\noindent (a) 对问题 \eqref{eq:9.92} 或 (92$'$) 有几种技巧可用:

\begin{itemize}
	\item \emph{单调性方法}, 见 F.~Browder \bref{Browder}{1} 与
	      J.-L.~Lions \bref{Lions}{3}.
	\item \emph{拓扑方法} (Schauder 不动点定理, Leray--Schauder 度理论
	      等); 见 J.~T.~Schwartz \bref{Schwartz}{1},
	      M.~Krasnoselskii \bref{Krasnoselskii}{1}, 以及
	      L.~Nirenberg \bref{Nirenberg}{2}, \bref{Nirenberg}{3}.
	\item \emph{变分方法} (临界点理论, 极小极大技巧, Morse 理论等); 见
	      P.~Rabinowitz \bref{Rabinowitz}{1}, \bref{Rabinowitz}{2},
	      Melvyn Berger \bref{Berger}{1},
	      M.~Krasnoselskii \bref{Krasnoselskii}{1},
	      L.~Nirenberg \bref{Nirenberg}{3},
	      J.~Mawhin--M.~Willem \bref{MawhinWillem}{1},
	      M.~Willem \bref{Willem}{1}, M.~Struwe \bref{Struwe}{1}.
\end{itemize}

一般性的综述参见例如 A.~Ambrosetti--G.~Prodi \bref{AmbrosettiProdi}{1}
与 E.~Zeidler \bref{Zeidler}{1} 的书.

\noindent (b) 求解 \eqref{eq:9.93} 型问题可能要涉及精细的估计技巧;\footnote{
例如极小曲面方程就是这种情形.} 见 E.~De Giorgi,
O.~Ladyzhenskaya--N.~Uraltseva \bref{Ladyzhenskaya}{1},
J.~Serrin \bref{Serrin}{1}, E.~Bombieri \bref{Bombieri}{1} 以及
D.~Gilbarg--N.~Trudinger \bref{GilbargTrudinger}{1} 的著作.关于完全非线性
方程特别是 Monge--Amp\`ere 方程, 近来也取得了重要进展; 参见例如
S.~T.~Yau \bref{Yau}{1}, L.~Caffarelli--L.~Nirenberg--J.~Spruck
\bref{CaffarelliNirenbergSpruck}{1} 以及 X.~Cabr\'e--L.~Caffarelli
\bref{CabreCaffarelli}{1}.

\noindent (c) 关于自由边界问题, 近年来出现了许多新结果, 它们常常与变分
不等式理论相联系; 参见例如 D.~Kinderlehrer--G.~Stampacchia
\bref{KinderlehrerStampacchia}{1}, C.~Baiocchi--A.~Capelo
\bref{BaiocchiCapelo}{1}, A.~Friedman \bref{Friedman}{4},
J.~Crank \bref{Crank}{1} 以及 L.~Caffarelli--S.~Salsa
\bref{CaffarelliSalsa}{1}.

\cmtnum{16.} \emph{几何测度论.}

这个领域处于几何与偏微分方程的交叉处, 自 20 世纪 60 年代以来有了广泛的
发展, 其起点是 H.~Federer, E.~De Giorgi, A.~I.~Volpert 与 F.~Almgren 的
基本贡献, 它们与变分法, 等周不等式等中产生的问题相联系.它在物理问题中有
大量应用, 例如相变, 力学中的断裂, 图像处理中的边缘检测, 液晶中的线涡, 超导体与超流体.空间 $BV$ (有界变差函数) 在这些问题中起特殊的作用.我们
参见例如 L.~Ambrosio--N.~Fusco--D.~Pallara \bref{Ambrosio}{1},
L.~Simon \bref{Simon}{1}, L.~C.~Evans--R.~Gariepy \bref{EvansGariepy}{1}
以及 F.~H.~Lin--X.~P.~Yang \bref{LinYang}{1}.
